\documentclass[11pt,a4paper]{report}

% Packages
\usepackage[utf8]{inputenc}
\usepackage[T1]{fontenc}
\usepackage{lmodern}     % scalable (Type1) fonts -> crisp, selectable, searchable
\usepackage{microtype}  % better line-breaking & readability, fewer overfull lines
\usepackage{geometry}
\usepackage{graphicx}
\usepackage{longtable}
\usepackage{booktabs}
\usepackage{array}
\usepackage{xcolor}
\usepackage{colortbl}
\usepackage{amssymb}
\usepackage{pifont}
\usepackage{textcomp}
\usepackage{hyperref}
\usepackage{xurl} % allow long URLs to break across lines
\usepackage{fancyhdr}
\usepackage{titlesec}
\usepackage{enumitem}
\usepackage{multirow}
\usepackage{pdflscape}
\usepackage{float}
\usepackage{caption}

% Map Unicode glyphs used in the body so pdfLaTeX compiles cleanly
\DeclareUnicodeCharacter{2713}{\ding{51}}   % checkmark
\DeclareUnicodeCharacter{2714}{\ding{52}}   % heavy checkmark
\DeclareUnicodeCharacter{2717}{\ding{55}}   % cross
\DeclareUnicodeCharacter{26A0}{\textbf{!}}  % warning sign
\DeclareUnicodeCharacter{00B0}{\textdegree} % degree
\DeclareUnicodeCharacter{2265}{$\geq$}
\DeclareUnicodeCharacter{2264}{$\leq$}
\DeclareUnicodeCharacter{2192}{$\rightarrow$}
\DeclareUnicodeCharacter{00D7}{$\times$}

% Page geometry
\geometry{
    a4paper,
    left=25mm,
    right=25mm,
    top=30mm,
    bottom=30mm,
    headheight=15pt
}

% Colors
\definecolor{bvblue}{RGB}{0,51,102}
\definecolor{bvlightblue}{RGB}{51,102,153}
\definecolor{bvgray}{RGB}{128,128,128}
\definecolor{lightgray}{RGB}{240,240,240}
\definecolor{warningred}{RGB}{204,0,0}

% Hyperref setup
\hypersetup{
    colorlinks=true,
    linkcolor=bvblue,
    filecolor=bvblue,
    urlcolor=bvblue,
    citecolor=bvblue,
    bookmarks=true,
    bookmarksopen=true,
    bookmarksopenlevel=1,
    bookmarksnumbered=true,
    breaklinks=true,
    pdfstartview=FitH,
    pdfdisplaydoctitle=true,
    pdftitle={Bureau Veritas Thruster Requirements - Surveyor Catalog},
    pdfauthor={Bureau Veritas},
    pdfsubject={Thruster Plan Approval Requirements},
    pdfkeywords={Bureau Veritas, Thrusters, NR467, Classification, Marine}
}

% Slightly more open line spacing for on-screen readability
\linespread{1.05}

% Header and footer
\pagestyle{fancy}
\fancyhf{}
\fancyhead[L]{\small\textcolor{bvblue}{BV Thruster Requirements Catalog}}
\fancyhead[R]{\small\textcolor{bvblue}{NR467 Sections 7, 14, 15}}
\fancyfoot[C]{\thepage}
\renewcommand{\headrulewidth}{0.5pt}
\renewcommand{\footrulewidth}{0.5pt}

% Title formatting
\titleformat{\chapter}[display]
{\normalfont\huge\bfseries\color{bvblue}}
{\chaptertitlename\ \thechapter}{20pt}{\Huge}

\titleformat{\section}
{\normalfont\Large\bfseries\color{bvblue}}
{\thesection}{1em}{}

\titleformat{\subsection}
{\normalfont\large\bfseries\color{bvlightblue}}
{\thesubsection}{1em}{}

% Custom commands
\newcommand{\fa}{\textbf{\textcolor{warningred}{FA}}}
\newcommand{\fic}{\textbf{\textcolor{bvgray}{FI}}}
\newcommand{\checkitem}{\item[$\square$]}

% --- Submission cover-sheet helpers (Appendix) ---
% Identification block for each ship-type cover sheet
\newcommand{\CoverID}{%
\noindent{\footnotesize
\begin{tabular}{|p{3.1cm}|p{4.1cm}|p{3.1cm}|p{3.1cm}|}
\hline
\rowcolor{lightgray}\textbf{Vessel / Yard No} & & \textbf{Class notation} & \\
\hline
\textbf{IMO number} & & \textbf{Thruster type / qty} & \\
\hline
\textbf{Prime mover} & & \textbf{Reviewer / date} & \\
\hline
\end{tabular}}\par\vspace{2mm}}
% Common core submissions applicable to every thruster package
\newcommand{\CoreCoverItems}{%
\item General Arrangement + Technical Specification sheet (\fa{})
\item \textbf{Type-specific mechanical/structural package} --- attach the matching Chapter \ref{ch:checklists} checklist (\fa{})
\item Materials, welding procedures, NDT plan (\fa{})
\item Thrust / vibration / noise / thermal calculations (\fa{})
\item Control system package + electrical single-line diagram (\fa{})
\item \textbf{Prime-mover package} (Chapter \ref{ch:primemover}): if \emph{electric} --- motor + VFD + harmonic study + cabling/segregation + protection; if \emph{diesel} --- engine approval + \textbf{full-train TVC} + clutch/coupling + FO/LO/cooling/starting/exhaust + controls (\fa{})
\item FAT / dock-trial / sea-trial procedures (\fa{}); test reports (\fic{})
\item O\&M, spare-parts \& installation manuals (\fic{})}
% Section header that always starts on a fresh page with an ID block
\newcommand{\CoverSheet}[1]{\clearpage\section{#1}\CoverID}

% Document begins
\begin{document}

% Title page
\begin{titlepage}
    \centering
    \vspace*{2cm}
    
    {\Huge\bfseries\color{bvblue} Bureau Veritas\par}
    \vspace{0.5cm}
    {\LARGE\bfseries Thruster Requirements\par}
    \vspace{0.3cm}
    {\LARGE\bfseries Comprehensive Surveyor Catalog\par}
    
    \vspace{2cm}
    
    {\Large\bfseries Based on NR467 Rules for Classification of Steel Ships\par}
    \vspace{0.5cm}
    {\large Sections 7, 14, and 15\par}
    
    \vspace{2cm}
    
    {\large\textit{Complete Plan Approval Reference Guide}\par}
    
    \vfill
    
    {\large Document Version 1.6\par}
    {\large 2026\par}
    
    \vspace{1cm}
    
    {\small\textcolor{warningred}{\textbf{FOR BUREAU VERITAS PLAN APPROVAL SURVEYORS}}\par}
    
\end{titlepage}

% Table of contents
\tableofcontents
\newpage

% Document Purpose
\chapter*{Document Purpose and Usage}
\addcontentsline{toc}{chapter}{Document Purpose and Usage}

This catalog provides a comprehensive overview of thruster requirements per Bureau Veritas Rules for Classification (NR467) and related guidance, designed as a practical reference for plan approval surveyors.

\section*{Key Features}
\begin{itemize}[leftmargin=*]
    \item Complete drawing and document lists for all thruster types
    \item Explicit approval categories (\fa{} or \fic{})
    \item Cross-references for overlapping requirements (steering gear, main shafting, hydraulics)
    \item Surveyor checklists by thruster type
    \item Testing and commissioning requirements
    \item Common submission errors and review tips
\end{itemize}

\section*{Document Structure}
The catalog is organized into the following chapters:
\begin{enumerate}[leftmargin=*]
    \item General Requirements (All Thrusters >110kW)
    \item Transverse Thrusters (Tunnel \& Retractable)
    \item Rotating and Azimuth Thrusters (incl. Podded Units)
    \item Water-Jet Propulsion Systems
    \item Voith-Schneider (Cycloidal) Propulsion Systems
    \item Other Thruster Types (Rim-Driven, Pump-Jet, Gill/Jet, Steering Grid)
    \item Prime Mover Considerations --- Diesel-Engine vs. Electric-Motor Driven Thrusters
    \item Thrusters as Part of Steering Gear
    \item Thrusters as Part of Main Shafting
    \item Hydraulic Systems \& Pitch Control
    \item Dynamic Positioning (DP) Thrusters
    \item Ice Class \& Cold-Climate (Winterisation) Requirements
    \item Additional BV Class Notations Affecting Thruster Packages
    \item Ship-Type \& Service-Notation Specific Requirements
    \item Testing \& Commissioning
    \item Surveyor Checklists by Type
    \item Submission Cover Sheets by Ship Type (one-page tick-off)
    \item Common Submission Errors
    \item Quick Reference Tables
    \item Regulatory References
    \item Contact Information
\end{enumerate}

\section*{How to Use This Catalog (Surveyor Workflow)}
\begin{enumerate}[leftmargin=*]
    \item \textbf{Identify the thruster type} (Chapters \ref{ch:transverse}--\ref{ch:othertypes}) and confirm rated power. Thrusters $>$110~kW attract the full Sec.\,15 submission set.
    \item \textbf{Identify the prime mover} --- electric motor or diesel engine (Chapter \ref{ch:primemover}) --- and add the corresponding drive-train drawings.
    \item \textbf{Identify the functional role(s)}: propulsion, steering (Chapter \ref{ch:steering}), main shaft line (Chapter \ref{ch:shafting}), or dynamic positioning (Chapter \ref{ch:dp}). Each role \emph{adds} to the base package.
    \item \textbf{Overlay additional notations \& service constraints}: ice class / winterisation (Chapter \ref{ch:ice}), other BV additional class notations (Chapter \ref{ch:notations}), and the ship-type / service-notation requirements incl.\ SRtP, hazardous-area and offshore rules (Chapter \ref{ch:shiptype}).
    \item \textbf{Verify hydraulics/pitch} (Chapter \ref{ch:hydraulics}) and \textbf{testing \& commissioning} (Chapter \ref{ch:testing}).
    \item \textbf{Run the type-specific checklist} (Chapter \ref{ch:checklists}) and screen against common errors (Chapter \ref{ch:errors}).
\end{enumerate}

\section*{Navigating this PDF}
This document is fully hyperlinked for quick reference:
\begin{itemize}[leftmargin=*]
    \item \textbf{Bookmarks panel} --- open your PDF reader's bookmarks/outline panel to jump to any chapter or section.
    \item \textbf{Table of contents} --- every entry is clickable.
    \item \textbf{Cross-references} --- ``Chapter~X'', ``Section~Y'' and table references in the text are clickable links.
    \item \textbf{Search} --- the text is fully selectable and searchable (e.g.\ search ``MON-SHAFT'', ``escort tug'', ``ice'').
\end{itemize}

\section*{Legend}
\begin{itemize}[leftmargin=*]
    \item \fa{} = \textbf{For Approval} — Must be reviewed and approved before construction/installation
    \item \fic{} = \textbf{For Information} — Submitted for record, not requiring formal approval
    \item \textbf{During Construction} = Certificates/reports generated during manufacturing
\end{itemize}

\vspace{0.3cm}
\noindent\textit{Note on rule references:} Thruster requirements are found in \textbf{NR467, Part C, Chapter 1, Section 15} (``Thrusters''); steering gear in \textbf{Section 14}; shaft lines and propellers in \textbf{Sections 7--8}. Bracketed sub-references such as [2.3] are structural pointers within this catalog. Ice class, comfort, automation and DP requirements sit in \textbf{NR467 Part F} (Additional Class Notations, e.g.\ DYNAPOS for DP); service notations are in \textbf{Part D} (general) and \textbf{Part E} (offshore service vessels \& tugs). Note \textbf{NR216} is the Materials \& Welding rule note (not a DP rule). Always confirm the exact article against the rule edition applicable to the vessel's contract/build date.

\newpage

% Chapter 1: General Requirements
\chapter{General Requirements - All Thrusters >110kW}
\label{ch:general}

\section{Applicability}
All thrusters with power output exceeding 110 kW must comply with NR467 Section 15 requirements. Thrusters $\leq$110 kW may have reduced submission requirements (verify with surveyor).

\section{Two Axes to Classify Before Listing Drawings}
\label{sec:duty-config}

Before applying any type checklist, fix two attributes of the unit --- they change
\emph{which} rules and drawings apply as much as the thruster type does:

\subsection{Duty: manoeuvring vs.\ propulsion}
\begin{itemize}[leftmargin=*]
\item \textbf{Manoeuvring (auxiliary)} --- e.g.\ a bow/stern tunnel thruster used
only for low-speed manoeuvring. Assessed under Sec.\,15 as auxiliary machinery;
it is \emph{not} part of the main shaft line.
\item \textbf{Propulsion (main)} --- e.g.\ an azimuth on a tug, a podded unit, a
water-jet or a Voith-Schneider that propels the ship. The propeller, shaft and
gear are classed as \textbf{main propulsion}: the main-shaft-line requirements of
Chapter \ref{ch:shafting} apply, a \textbf{torsional vibration calculation is
mandatory}, and barred/continuous-speed ranges must be confirmed.
\item \textbf{Take-me-home / emergency propulsion} --- a retractable or auxiliary
unit providing get-home capability; assess independence from the main train
(cross-reference redundant-propulsion notation, Chapter \ref{ch:notations}).
\end{itemize}
The same physical thruster therefore attracts a heavier package when it is the
ship's propulsion than when it is purely for manoeuvring.

\subsection{Propeller configuration: open vs.\ ducted (nozzle)}
\begin{itemize}[leftmargin=*]
\item \textbf{Open (non-ducted) propeller} --- blade and hub strength, cavitation
and the usual shaft/bearing package.
\item \textbf{Ducted (nozzle) propeller} --- \emph{adds} the nozzle/duct structure
and strength drawing, nozzle profile/geometry with propeller-tip clearance, and
the nozzle-to-unit/hull connection. Distinguish:
  \begin{itemize}
  \item \textbf{Fixed nozzle} (Kort-type) --- structural item only; improves
  bollard pull.
  \item \textbf{Steering nozzle} --- the nozzle also steers, so the
  \textbf{steering-gear package} of Chapter \ref{ch:steering} applies (actuator,
  control, redundancy, SOLAS II-1/29 where it is the ship's means of steering).
  \end{itemize}
\end{itemize}

\noindent\textbf{Surveyor note:} record duty and propeller configuration on the
submission cover sheet (Chapter \ref{ch:coversheets}). The interactive
Quick-Reference tool builds the exact drawing list once type, duty, propeller
configuration, prime mover, role, ship type and notations are selected.

\section{Mandatory Submissions for All Thruster Types}

\begin{longtable}{|p{0.5cm}|p{5cm}|p{1.5cm}|p{5.5cm}|p{2cm}|}
\caption{General Requirements - All Thrusters} \label{tab:general-req} \\
\hline
\rowcolor{bvblue}
\textcolor{white}{\textbf{\#}} & \textcolor{white}{\textbf{Drawing/Document}} & \textcolor{white}{\textbf{Category}} & \textcolor{white}{\textbf{Content Requirements}} & \textcolor{white}{\textbf{Reference}} \\
\hline
\endfirsthead

\multicolumn{5}{c}%
{{\tablename\ \thetable{} -- continued from previous page}} \\
\hline
\rowcolor{bvblue}
\textcolor{white}{\textbf{\#}} & \textcolor{white}{\textbf{Drawing/Document}} & \textcolor{white}{\textbf{Category}} & \textcolor{white}{\textbf{Content Requirements}} & \textcolor{white}{\textbf{Reference}} \\
\hline
\endhead

\hline \multicolumn{5}{|r|}{{Continued on next page}} \\ \hline
\endfoot

\hline
\endlastfoot

1 & General Arrangement Drawing & \fa{} & 
\begin{itemize}[leftmargin=*,nosep]
\item Thruster location on vessel
\item Installation orientation
\item Access provisions
\item Clearances to other equipment
\item Coordinate system reference
\end{itemize} & NR467 Sec.15, [1] \\
\hline

2 & Technical Specification Sheet & \fa{} & 
\begin{itemize}[leftmargin=*,nosep]
\item Rated power (kW)
\item Rated thrust (kN)
\item Rated speed (RPM)
\item Operating voltage/pressure
\item Weight and dimensions
\item Operating depth/pressure rating
\end{itemize} & NR467 Sec.15, [1] \\
\hline

3 & Material Specifications & \fa{} & 
\begin{itemize}[leftmargin=*,nosep]
\item Material grades for all major components
\item Compliance with NR467 Part B, Ch.3
\item Corrosion protection details
\item Material certificates (during construction)
\end{itemize} & NR467 Part B, Ch.3 \\
\hline

4 & Welding Procedures & \fa{} & 
\begin{itemize}[leftmargin=*,nosep]
\item Welding Procedure Specifications (WPS)
\item Welder qualification records
\item NDT requirements and extent
\item Weld inspection plan
\end{itemize} & NR467 Part B, Ch.4 \\
\hline

5 & Manufacturer's Quality Plan & \fa{} & 
\begin{itemize}[leftmargin=*,nosep]
\item Manufacturing sequence
\item Inspection and test points
\item Quality control procedures
\item Traceability system
\end{itemize} & NR467 Sec.15 \\
\hline

6 & Operation \& Maintenance Manual & \fic{} & 
\begin{itemize}[leftmargin=*,nosep]
\item Operating instructions
\item Maintenance schedule
\item Troubleshooting guide
\item Safety precautions
\end{itemize} & NR467 Sec.15 \\
\hline

7 & Spare Parts List & \fic{} & 
\begin{itemize}[leftmargin=*,nosep]
\item Recommended spare parts
\item Critical component identification
\item Part numbers and suppliers
\end{itemize} & NR467 Sec.15 \\
\hline

8 & Installation Manual & \fic{} & 
\begin{itemize}[leftmargin=*,nosep]
\item Installation procedure
\item Alignment requirements
\item Foundation specifications
\item Connection details
\end{itemize} & NR467 Sec.15 \\
\hline

\end{longtable}

\newpage

% Chapter 2: Transverse Thrusters
\chapter{Transverse Thrusters (Tunnel \& Retractable)}
\label{ch:transverse}

\section{Tunnel Thrusters (Fixed Installation)}

\subsection{Structural Drawings \& Calculations}

\begin{longtable}{|p{0.5cm}|p{5cm}|p{1.5cm}|p{5.5cm}|p{2cm}|}
\caption{Tunnel Thrusters - Structural Requirements} \label{tab:tunnel-struct} \\
\hline
\rowcolor{bvblue}
\textcolor{white}{\textbf{\#}} & \textcolor{white}{\textbf{Drawing/Document}} & \textcolor{white}{\textbf{Category}} & \textcolor{white}{\textbf{Content Requirements}} & \textcolor{white}{\textbf{Reference}} \\
\hline
\endfirsthead

\multicolumn{5}{c}%
{{\tablename\ \thetable{} -- continued from previous page}} \\
\hline
\rowcolor{bvblue}
\textcolor{white}{\textbf{\#}} & \textcolor{white}{\textbf{Drawing/Document}} & \textcolor{white}{\textbf{Category}} & \textcolor{white}{\textbf{Content Requirements}} & \textcolor{white}{\textbf{Reference}} \\
\hline
\endhead

\hline \multicolumn{5}{|r|}{{Continued on next page}} \\ \hline
\endfoot

\hline
\endlastfoot

1 & Tunnel Structural Arrangement & \fa{} & 
\begin{itemize}[leftmargin=*,nosep]
\item Tunnel diameter and length
\item Tunnel position relative to baseline/centerline
\item Plate thickness and material
\item Reinforcement arrangement
\item Connection to hull structure
\end{itemize} & NR467 Sec.15, [2.1] \\
\hline

2 & Tunnel Structural Calculations & \fa{} & 
\begin{itemize}[leftmargin=*,nosep]
\item Pressure loads (internal/external)
\item Structural stress analysis
\item Buckling analysis
\item Fatigue assessment
\item FEA results (if applicable)
\end{itemize} & NR467 Sec.15, [2.1] \\
\hline

3 & Tunnel End Grating/Grid Design & \fa{} & 
\begin{itemize}[leftmargin=*,nosep]
\item Grating dimensions and material
\item Attachment method
\item Structural strength verification
\item Debris protection
\end{itemize} & NR467 Sec.15, [2.2] \\
\hline

4 & Hull Penetration Details & \fa{} & 
\begin{itemize}[leftmargin=*,nosep]
\item Opening dimensions
\item Reinforcement around opening
\item Welding details
\item Watertight integrity provisions
\end{itemize} & NR467 Part B, Ch.5 \\
\hline

\end{longtable}

\subsection{Thruster Unit - Mechanical Components}

\begin{longtable}{|p{0.5cm}|p{5cm}|p{1.5cm}|p{5.5cm}|p{2cm}|}
\caption{Tunnel Thrusters - Mechanical Components} \label{tab:tunnel-mech} \\
\hline
\rowcolor{bvblue}
\textcolor{white}{\textbf{\#}} & \textcolor{white}{\textbf{Drawing/Document}} & \textcolor{white}{\textbf{Category}} & \textcolor{white}{\textbf{Content Requirements}} & \textcolor{white}{\textbf{Reference}} \\
\hline
\endfirsthead

\multicolumn{5}{c}%
{{\tablename\ \thetable{} -- continued from previous page}} \\
\hline
\rowcolor{bvblue}
\textcolor{white}{\textbf{\#}} & \textcolor{white}{\textbf{Drawing/Document}} & \textcolor{white}{\textbf{Category}} & \textcolor{white}{\textbf{Content Requirements}} & \textcolor{white}{\textbf{Reference}} \\
\hline
\endhead

\hline \multicolumn{5}{|r|}{{Continued on next page}} \\ \hline
\endfoot

\hline
\endlastfoot

5 & Thruster Unit General Assembly & \fa{} & 
\begin{itemize}[leftmargin=*,nosep]
\item Complete unit cross-section
\item Component identification
\item Material specifications
\item Assembly sequence
\end{itemize} & NR467 Sec.15, [2.3] \\
\hline

6 & Propeller Design Drawing & \fa{} & 
\begin{itemize}[leftmargin=*,nosep]
\item Propeller diameter and pitch
\item Number of blades
\item Blade profile and geometry
\item Hub design
\item Material specification
\end{itemize} & NR467 Sec.15, [2.4] \\
\hline

7 & Propeller Strength Calculations & \fa{} & 
\begin{itemize}[leftmargin=*,nosep]
\item Blade root stress analysis
\item Hub stress analysis
\item Centrifugal force calculations
\item Fatigue life assessment
\item Cavitation analysis
\end{itemize} & NR467 Sec.15, [2.4] \\
\hline

8 & Propeller Shaft Drawing & \fa{} & 
\begin{itemize}[leftmargin=*,nosep]
\item Shaft diameter and length
\item Material specification
\item Surface finish requirements
\item Keyway details
\item Seal locations
\end{itemize} & NR467 Sec.15, [2.5] \\
\hline

9 & Shaft Strength Calculations & \fa{} & 
\begin{itemize}[leftmargin=*,nosep]
\item Torsional strength analysis
\item Bending stress calculations
\item Combined stress analysis
\item Deflection calculations
\item Critical speed analysis
\item Fatigue assessment
\end{itemize} & NR467 Sec.15, [2.5]; Sec.7 \\
\hline

10 & Bearing Arrangement Drawing & \fa{} & 
\begin{itemize}[leftmargin=*,nosep]
\item Bearing type and manufacturer
\item Bearing housing design
\item Lubrication arrangement
\item Seal design
\item Cooling provisions (if applicable)
\end{itemize} & NR467 Sec.15, [2.6] \\
\hline

11 & Bearing Load Calculations & \fa{} & 
\begin{itemize}[leftmargin=*,nosep]
\item Radial and axial loads
\item Bearing life calculation (L10)
\item Lubrication requirements
\item Temperature analysis
\end{itemize} & NR467 Sec.15, [2.6] \\
\hline

12 & Gearbox Assembly Drawing & \fa{} & 
\begin{itemize}[leftmargin=*,nosep]
\item Gear arrangement (bevel, helical, etc.)
\item Gear ratios
\item Casing design
\item Lubrication system
\item Mounting arrangement
\end{itemize} & NR467 Sec.15, [2.7] \\
\hline

13 & Gear Tooth Calculations & \fa{} & 
\begin{itemize}[leftmargin=*,nosep]
\item Gear tooth strength (ISO 6336)
\item Contact stress analysis
\item Bending stress analysis
\item Lubrication film thickness
\item Efficiency calculations
\end{itemize} & NR467 Sec.15, [2.7] \\
\hline

14 & Gearbox Lubrication System & \fa{} & 
\begin{itemize}[leftmargin=*,nosep]
\item Oil type and capacity
\item Circulation system (if forced)
\item Cooling arrangement
\item Filtration system
\item Oil level and temperature monitoring
\end{itemize} & NR467 Sec.15, [2.7] \\
\hline

\end{longtable}

\subsection{Drive System (Electric or Hydraulic)}

\subsubsection{For Electric Motor Drive}

\begin{longtable}{|p{0.5cm}|p{5cm}|p{1.5cm}|p{5.5cm}|p{2cm}|}
\caption{Tunnel Thrusters - Electric Drive System} \label{tab:tunnel-electric} \\
\hline
\rowcolor{bvblue}
\textcolor{white}{\textbf{\#}} & \textcolor{white}{\textbf{Drawing/Document}} & \textcolor{white}{\textbf{Category}} & \textcolor{white}{\textbf{Content Requirements}} & \textcolor{white}{\textbf{Reference}} \\
\hline
\endfirsthead

\multicolumn{5}{c}%
{{\tablename\ \thetable{} -- continued from previous page}} \\
\hline
\rowcolor{bvblue}
\textcolor{white}{\textbf{\#}} & \textcolor{white}{\textbf{Drawing/Document}} & \textcolor{white}{\textbf{Category}} & \textcolor{white}{\textbf{Content Requirements}} & \textcolor{white}{\textbf{Reference}} \\
\hline
\endhead

\hline \multicolumn{5}{|r|}{{Continued on next page}} \\ \hline
\endfoot

\hline
\endlastfoot

15 & Electric Motor Specification & \fa{} & 
\begin{itemize}[leftmargin=*,nosep]
\item Motor type (induction, synchronous)
\item Rated power and speed
\item Voltage and frequency
\item Insulation class
\item Protection class (IP rating)
\item Cooling method
\item Mounting arrangement
\end{itemize} & NR467 Sec.15, [2.8]; IEC 60092-301 \\
\hline

16 & Motor Foundation Drawing & \fa{} & 
\begin{itemize}[leftmargin=*,nosep]
\item Foundation structure design
\item Vibration isolation details
\item Bolt arrangement
\item Alignment provisions
\item Structural attachment to hull
\end{itemize} & NR467 Sec.15, [2.8] \\
\hline

17 & Motor-Gearbox Coupling & \fa{} & 
\begin{itemize}[leftmargin=*,nosep]
\item Coupling type and manufacturer
\item Alignment tolerances
\item Torque capacity
\item Installation procedure
\end{itemize} & NR467 Sec.15, [2.8] \\
\hline

\end{longtable}

\subsubsection{For Hydraulic Motor Drive}

\begin{longtable}{|p{0.5cm}|p{5cm}|p{1.5cm}|p{5.5cm}|p{2cm}|}
\caption{Tunnel Thrusters - Hydraulic Drive System} \label{tab:tunnel-hydraulic} \\
\hline
\rowcolor{bvblue}
\textcolor{white}{\textbf{\#}} & \textcolor{white}{\textbf{Drawing/Document}} & \textcolor{white}{\textbf{Category}} & \textcolor{white}{\textbf{Content Requirements}} & \textcolor{white}{\textbf{Reference}} \\
\hline
\endfirsthead

\multicolumn{5}{c}%
{{\tablename\ \thetable{} -- continued from previous page}} \\
\hline
\rowcolor{bvblue}
\textcolor{white}{\textbf{\#}} & \textcolor{white}{\textbf{Drawing/Document}} & \textcolor{white}{\textbf{Category}} & \textcolor{white}{\textbf{Content Requirements}} & \textcolor{white}{\textbf{Reference}} \\
\hline
\endhead

\hline \multicolumn{5}{|r|}{{Continued on next page}} \\ \hline
\endfoot

\hline
\endlastfoot

18 & Hydraulic Motor Specification & \fa{} & 
\begin{itemize}[leftmargin=*,nosep]
\item Motor type (axial piston, radial piston, gear)
\item Displacement and speed range
\item Rated pressure
\item Efficiency curves
\item Mounting details
\end{itemize} & NR467 Sec.15, [2.9] \\
\hline

19 & Hydraulic System Schematic & \fa{} & 
\begin{itemize}[leftmargin=*,nosep]
\item Complete hydraulic circuit
\item Component symbols and identification
\item Pressure and flow ratings
\item See Chapter \ref{ch:hydraulics} for detailed requirements
\end{itemize} & NR467 Sec.15, [2.9] \\
\hline

\end{longtable}

\subsection{Control System}

\begin{longtable}{|p{0.5cm}|p{5cm}|p{1.5cm}|p{5.5cm}|p{2cm}|}
\caption{Tunnel Thrusters - Control System} \label{tab:tunnel-control} \\
\hline
\rowcolor{bvblue}
\textcolor{white}{\textbf{\#}} & \textcolor{white}{\textbf{Drawing/Document}} & \textcolor{white}{\textbf{Category}} & \textcolor{white}{\textbf{Content Requirements}} & \textcolor{white}{\textbf{Reference}} \\
\hline
\endfirsthead

\multicolumn{5}{c}%
{{\tablename\ \thetable{} -- continued from previous page}} \\
\hline
\rowcolor{bvblue}
\textcolor{white}{\textbf{\#}} & \textcolor{white}{\textbf{Drawing/Document}} & \textcolor{white}{\textbf{Category}} & \textcolor{white}{\textbf{Content Requirements}} & \textcolor{white}{\textbf{Reference}} \\
\hline
\endhead

\hline \multicolumn{5}{|r|}{{Continued on next page}} \\ \hline
\endfoot

\hline
\endlastfoot

20 & Control System Block Diagram & \fa{} & 
\begin{itemize}[leftmargin=*,nosep]
\item Control architecture
\item Input/output signals
\item Control logic overview
\item Interface with ship systems
\end{itemize} & NR467 Sec.15, [2.10]; Part C, Ch.4 \\
\hline

21 & Control Panel Layout & \fa{} & 
\begin{itemize}[leftmargin=*,nosep]
\item Panel arrangement
\item Component locations
\item Wiring diagram
\item Protection devices
\end{itemize} & NR467 Sec.15, [2.10] \\
\hline

22 & Electrical Single-Line Diagram & \fa{} & 
\begin{itemize}[leftmargin=*,nosep]
\item Power supply arrangement
\item Circuit breakers and protection
\item Emergency stop circuit
\item Monitoring and alarm systems
\end{itemize} & NR467 Sec.15, [2.10]; IEC 60092-504 \\
\hline

23 & Control Software Description & \fa{} & 
\begin{itemize}[leftmargin=*,nosep]
\item Software functionality
\item Control algorithms
\item Safety functions
\item Software validation documentation
\end{itemize} & NR467 Part C, Ch.4 \\
\hline

\end{longtable}

\subsection{Additional Calculations \& Documentation}

\begin{longtable}{|p{0.5cm}|p{5cm}|p{1.5cm}|p{5.5cm}|p{2cm}|}
\caption{Tunnel Thrusters - Additional Calculations} \label{tab:tunnel-calc} \\
\hline
\rowcolor{bvblue}
\textcolor{white}{\textbf{\#}} & \textcolor{white}{\textbf{Drawing/Document}} & \textcolor{white}{\textbf{Category}} & \textcolor{white}{\textbf{Content Requirements}} & \textcolor{white}{\textbf{Reference}} \\
\hline
\endfirsthead

\multicolumn{5}{c}%
{{\tablename\ \thetable{} -- continued from previous page}} \\
\hline
\rowcolor{bvblue}
\textcolor{white}{\textbf{\#}} & \textcolor{white}{\textbf{Drawing/Document}} & \textcolor{white}{\textbf{Category}} & \textcolor{white}{\textbf{Content Requirements}} & \textcolor{white}{\textbf{Reference}} \\
\hline
\endhead

\hline \multicolumn{5}{|r|}{{Continued on next page}} \\ \hline
\endfoot

\hline
\endlastfoot

24 & Thrust Calculation & \fa{} & 
\begin{itemize}[leftmargin=*,nosep]
\item Theoretical thrust calculation
\item Efficiency factors
\item Operating curves (thrust vs. speed)
\item Bollard pull prediction
\end{itemize} & NR467 Sec.15, [2.11] \\
\hline

25 & Vibration Analysis & \fa{} & 
\begin{itemize}[leftmargin=*,nosep]
\item Natural frequency calculations
\item Forced vibration analysis
\item Resonance avoidance verification
\item Vibration isolation design
\end{itemize} & NR467 Sec.15, [2.12] \\
\hline

26 & Noise Level Prediction & \fa{} & 
\begin{itemize}[leftmargin=*,nosep]
\item Airborne noise prediction
\item Structure-borne noise analysis
\item Compliance with IMO noise code
\item Noise reduction measures
\end{itemize} & NR467 Sec.15, [2.12] \\
\hline

27 & Thermal Analysis & \fa{} & 
\begin{itemize}[leftmargin=*,nosep]
\item Heat generation calculations
\item Cooling requirements
\item Temperature rise predictions
\item Ventilation requirements
\end{itemize} & NR467 Sec.15, [2.13] \\
\hline

\end{longtable}

\newpage

\section{Retractable Thrusters}

\subsection{Additional Requirements Beyond Fixed Tunnel Thrusters}

\textbf{All submissions from Section 2.1 apply, PLUS the following:}

\begin{longtable}{|p{0.5cm}|p{5cm}|p{1.5cm}|p{5.5cm}|p{2cm}|}
\caption{Retractable Thrusters - Additional Requirements} \label{tab:retractable} \\
\hline
\rowcolor{bvblue}
\textcolor{white}{\textbf{\#}} & \textcolor{white}{\textbf{Drawing/Document}} & \textcolor{white}{\textbf{Category}} & \textcolor{white}{\textbf{Content Requirements}} & \textcolor{white}{\textbf{Reference}} \\
\hline
\endfirsthead

\multicolumn{5}{c}%
{{\tablename\ \thetable{} -- continued from previous page}} \\
\hline
\rowcolor{bvblue}
\textcolor{white}{\textbf{\#}} & \textcolor{white}{\textbf{Drawing/Document}} & \textcolor{white}{\textbf{Category}} & \textcolor{white}{\textbf{Content Requirements}} & \textcolor{white}{\textbf{Reference}} \\
\hline
\endhead

\hline \multicolumn{5}{|r|}{{Continued on next page}} \\ \hline
\endfoot

\hline
\endlastfoot

28 & Retraction Mechanism Assembly & \fa{} & 
\begin{itemize}[leftmargin=*,nosep]
\item Complete mechanism design
\item Lifting/lowering system
\item Guide rail arrangement
\item Travel limits and stops
\item Position indicators
\end{itemize} & NR467 Sec.15, [2.14] \\
\hline

29 & Hydraulic Cylinder Specifications & \fa{} & 
\begin{itemize}[leftmargin=*,nosep]
\item Cylinder bore and stroke
\item Operating pressure
\item Rod diameter and material
\item Seal specifications
\item Mounting arrangement
\end{itemize} & NR467 Sec.15, [2.14] \\
\hline

30 & Retraction System Calculations & \fa{} & 
\begin{itemize}[leftmargin=*,nosep]
\item Lifting force requirements
\item Hydraulic pressure calculations
\item Cylinder sizing verification
\item Retraction/extension time calculations
\item Dynamic load analysis
\end{itemize} & NR467 Sec.15, [2.14] \\
\hline

31 & Hull Opening \& Closure Design & \fa{} & 
\begin{itemize}[leftmargin=*,nosep]
\item Opening dimensions
\item Closure plate/door design
\item Sealing arrangement (extended position)
\item Sealing arrangement (retracted position)
\item Locking mechanism
\end{itemize} & NR467 Sec.15, [2.15]; Part B, Ch.5 \\
\hline

32 & Hull Opening Structural Analysis & \fa{} & 
\begin{itemize}[leftmargin=*,nosep]
\item Local hull strength calculations
\item Reinforcement design around opening
\item Stress concentration analysis
\item Fatigue assessment
\item FEA results (typically required)
\end{itemize} & NR467 Sec.15, [2.15]; Part B, Ch.5 \\
\hline

33 & Watertight Integrity Documentation & \fa{} & 
\begin{itemize}[leftmargin=*,nosep]
\item Seal design and material
\item Pressure testing requirements
\item Leakage criteria
\item Compliance with subdivision requirements
\end{itemize} & NR467 Part B, Ch.5 \\
\hline

34 & Locking System Design & \fa{} & 
\begin{itemize}[leftmargin=*,nosep]
\item Mechanical locks (extended position)
\item Retracted position locks
\item Locking force calculations
\item Fail-safe provisions
\item Interlock system
\end{itemize} & NR467 Sec.15, [2.16] \\
\hline

35 & Position Indication System & \fa{} & 
\begin{itemize}[leftmargin=*,nosep]
\item Position sensors (type and location)
\item Indication on bridge/control panel
\item Alarm system (intermediate positions)
\item Redundancy provisions
\end{itemize} & NR467 Sec.15, [2.16] \\
\hline

36 & Emergency Lowering System & \fa{} & 
\begin{itemize}[leftmargin=*,nosep]
\item Emergency power source
\item Manual lowering provisions
\item Emergency procedure
\item Time to deploy
\end{itemize} & NR467 Sec.15, [2.17] \\
\hline

37 & Hydraulic System for Retraction & \fa{} & 
\begin{itemize}[leftmargin=*,nosep]
\item HPU specifications
\item Control valve arrangement
\item Piping layout
\item Accumulator sizing (if applicable)
\item Emergency lowering circuit
\item See Chapter \ref{ch:hydraulics} for detailed requirements
\end{itemize} & NR467 Sec.15, [2.17] \\
\hline

38 & Retraction/Extension Procedures & \fic{} & 
\begin{itemize}[leftmargin=*,nosep]
\item Step-by-step operating procedure
\item Pre-operation checks
\item Safety precautions
\item Troubleshooting
\end{itemize} & NR467 Sec.15, [2.18] \\
\hline

39 & Emergency Recovery Procedures & \fic{} & 
\begin{itemize}[leftmargin=*,nosep]
\item Emergency lowering procedure
\item Manual operation instructions
\item Recovery from failure modes
\item Contact information for support
\end{itemize} & NR467 Sec.15, [2.18] \\
\hline

\end{longtable}

\newpage

% Chapter 3: Azimuth Thrusters
\chapter{Rotating and Azimuth Thrusters}
\label{ch:azimuth}

\section{Azimuth Thrusters (Fixed Pitch \& Controllable Pitch)}

\subsection{General Arrangement \& Structure}

\begin{longtable}{|p{0.5cm}|p{5cm}|p{1.5cm}|p{5.5cm}|p{2cm}|}
\caption{Azimuth Thrusters - General Arrangement} \label{tab:azimuth-general} \\
\hline
\rowcolor{bvblue}
\textcolor{white}{\textbf{\#}} & \textcolor{white}{\textbf{Drawing/Document}} & \textcolor{white}{\textbf{Category}} & \textcolor{white}{\textbf{Content Requirements}} & \textcolor{white}{\textbf{Reference}} \\
\hline
\endfirsthead

\multicolumn{5}{c}%
{{\tablename\ \thetable{} -- continued from previous page}} \\
\hline
\rowcolor{bvblue}
\textcolor{white}{\textbf{\#}} & \textcolor{white}{\textbf{Drawing/Document}} & \textcolor{white}{\textbf{Category}} & \textcolor{white}{\textbf{Content Requirements}} & \textcolor{white}{\textbf{Reference}} \\
\hline
\endhead

\hline \multicolumn{5}{|r|}{{Continued on next page}} \\ \hline
\endfoot

\hline
\endlastfoot

1 & General Arrangement Drawing & \fa{} & 
\begin{itemize}[leftmargin=*,nosep]
\item Thruster location on vessel
\item Installation orientation
\item Azimuth rotation range (typically 360°)
\item Clearance envelope (rotating)
\item Access for maintenance
\end{itemize} & NR467 Sec.15, [3.1] \\
\hline

2 & Azimuth Unit Complete Assembly & \fa{} & 
\begin{itemize}[leftmargin=*,nosep]
\item Longitudinal and transverse sections
\item All major components identified
\item Material specifications
\item Assembly sequence
\item Lifting points
\end{itemize} & NR467 Sec.15, [3.1] \\
\hline

3 & Hull Penetration \& Foundation & \fa{} & 
\begin{itemize}[leftmargin=*,nosep]
\item Hull opening dimensions
\item Structural reinforcement
\item Foundation design
\item Watertight integrity details
\item Load transfer to hull structure
\end{itemize} & NR467 Sec.15, [3.2]; Part B, Ch.5 \\
\hline

4 & Hull Structure Calculations & \fa{} & 
\begin{itemize}[leftmargin=*,nosep]
\item Local hull strength analysis
\item Foundation load distribution
\item Stress analysis
\item Fatigue assessment
\item FEA results (typically required)
\end{itemize} & NR467 Sec.15, [3.2]; Part B, Ch.5 \\
\hline

\end{longtable}

\subsection{Propeller \& Shaft System}

\subsubsection{For Fixed Pitch Propellers}

\begin{longtable}{|p{0.5cm}|p{5cm}|p{1.5cm}|p{5.5cm}|p{2cm}|}
\caption{Azimuth Thrusters - Fixed Pitch Propeller} \label{tab:azimuth-fp} \\
\hline
\rowcolor{bvblue}
\textcolor{white}{\textbf{\#}} & \textcolor{white}{\textbf{Drawing/Document}} & \textcolor{white}{\textbf{Category}} & \textcolor{white}{\textbf{Content Requirements}} & \textcolor{white}{\textbf{Reference}} \\
\hline
\endfirsthead

\multicolumn{5}{c}%
{{\tablename\ \thetable{} -- continued from previous page}} \\
\hline
\rowcolor{bvblue}
\textcolor{white}{\textbf{\#}} & \textcolor{white}{\textbf{Drawing/Document}} & \textcolor{white}{\textbf{Category}} & \textcolor{white}{\textbf{Content Requirements}} & \textcolor{white}{\textbf{Reference}} \\
\hline
\endhead

\hline \multicolumn{5}{|r|}{{Continued on next page}} \\ \hline
\endfoot

\hline
\endlastfoot

5 & Propeller Design Drawing & \fa{} & 
\begin{itemize}[leftmargin=*,nosep]
\item Propeller diameter and pitch
\item Number of blades
\item Blade geometry (sections)
\item Hub design
\item Material specification
\item Blade attachment method
\end{itemize} & NR467 Sec.15, [3.3] \\
\hline

6 & Propeller Strength Calculations & \fa{} & 
\begin{itemize}[leftmargin=*,nosep]
\item Blade stress analysis
\item Hub stress analysis
\item Centrifugal forces
\item Hydrodynamic forces
\item Fatigue life assessment
\item Cavitation analysis
\end{itemize} & NR467 Sec.15, [3.3] \\
\hline

\end{longtable}

\subsubsection{For Controllable Pitch Propellers (ADDITIONAL to Fixed Pitch)}

\begin{longtable}{|p{0.5cm}|p{5cm}|p{1.5cm}|p{5.5cm}|p{2cm}|}
\caption{Azimuth Thrusters - Controllable Pitch Mechanism} \label{tab:azimuth-cp} \\
\hline
\rowcolor{bvblue}
\textcolor{white}{\textbf{\#}} & \textcolor{white}{\textbf{Drawing/Document}} & \textcolor{white}{\textbf{Category}} & \textcolor{white}{\textbf{Content Requirements}} & \textcolor{white}{\textbf{Reference}} \\
\hline
\endfirsthead

\multicolumn{5}{c}%
{{\tablename\ \thetable{} -- continued from previous page}} \\
\hline
\rowcolor{bvblue}
\textcolor{white}{\textbf{\#}} & \textcolor{white}{\textbf{Drawing/Document}} & \textcolor{white}{\textbf{Category}} & \textcolor{white}{\textbf{Content Requirements}} & \textcolor{white}{\textbf{Reference}} \\
\hline
\endhead

\hline \multicolumn{5}{|r|}{{Continued on next page}} \\ \hline
\endfoot

\hline
\endlastfoot

7 & CP Hub Mechanism Assembly & \fa{} & 
\begin{itemize}[leftmargin=*,nosep]
\item Hub internal mechanism
\item Pitch control rod/linkage
\item Blade retention system
\item Bearing arrangement
\item Seal design
\end{itemize} & NR467 Sec.15, [3.4] \\
\hline

8 & Blade Retention Calculations & \fa{} & 
\begin{itemize}[leftmargin=*,nosep]
\item Blade root connection strength
\item Centrifugal force analysis
\item Hydrodynamic moment analysis
\item Bolt/pin strength verification
\item Fatigue assessment
\end{itemize} & NR467 Sec.15, [3.4] \\
\hline

9 & Pitch Control Mechanism & \fa{} & 
\begin{itemize}[leftmargin=*,nosep]
\item Hydraulic cylinder (in hub or external)
\item Control rod design
\item Linkage system
\item Pitch range and indicators
\item Fail-safe position
\end{itemize} & NR467 Sec.15, [3.4]; Ch.\ref{ch:hydraulics} \\
\hline

10 & Pitch Control Hydraulic System & \fa{} & 
\begin{itemize}[leftmargin=*,nosep]
\item Hydraulic supply arrangement
\item Rotary joint/slip ring design
\item Control valve system
\item Pitch feedback system
\item See Chapter \ref{ch:hydraulics} for detailed requirements
\end{itemize} & NR467 Sec.15, [3.4]; Ch.\ref{ch:hydraulics} \\
\hline

\end{longtable}

\subsection{Vertical Shaft \& Column}

\begin{longtable}{|p{0.5cm}|p{5cm}|p{1.5cm}|p{5.5cm}|p{2cm}|}
\caption{Azimuth Thrusters - Vertical Shaft \& Column} \label{tab:azimuth-shaft} \\
\hline
\rowcolor{bvblue}
\textcolor{white}{\textbf{\#}} & \textcolor{white}{\textbf{Drawing/Document}} & \textcolor{white}{\textbf{Category}} & \textcolor{white}{\textbf{Content Requirements}} & \textcolor{white}{\textbf{Reference}} \\
\hline
\endfirsthead

\multicolumn{5}{c}%
{{\tablename\ \thetable{} -- continued from previous page}} \\
\hline
\rowcolor{bvblue}
\textcolor{white}{\textbf{\#}} & \textcolor{white}{\textbf{Drawing/Document}} & \textcolor{white}{\textbf{Category}} & \textcolor{white}{\textbf{Content Requirements}} & \textcolor{white}{\textbf{Reference}} \\
\hline
\endhead

\hline \multicolumn{5}{|r|}{{Continued on next page}} \\ \hline
\endfoot

\hline
\endlastfoot

11 & Vertical Shaft Assembly & \fa{} & 
\begin{itemize}[leftmargin=*,nosep]
\item Shaft dimensions and material
\item Upper and lower bearing locations
\item Thrust bearing arrangement
\item Seal system
\item Coupling details
\end{itemize} & NR467 Sec.15, [3.5]; Sec.7 \\
\hline

12 & Vertical Shaft Calculations & \fa{} & 
\begin{itemize}[leftmargin=*,nosep]
\item Torsional strength
\item Bending stress (from propeller forces)
\item Combined stress analysis
\item Deflection calculations
\item Critical speed analysis
\item Fatigue assessment
\end{itemize} & NR467 Sec.15, [3.5]; Sec.7 \\
\hline

13 & Thrust Bearing Design & \fa{} & 
\begin{itemize}[leftmargin=*,nosep]
\item Bearing type and manufacturer
\item Load capacity
\item Lubrication system
\item Cooling arrangement
\item Monitoring provisions
\end{itemize} & NR467 Sec.15, [3.6] \\
\hline

14 & Thrust Bearing Calculations & \fa{} & 
\begin{itemize}[leftmargin=*,nosep]
\item Axial load analysis
\item Bearing life calculation
\item Temperature rise analysis
\item Lubrication film thickness
\end{itemize} & NR467 Sec.15, [3.6] \\
\hline

15 & Radial Bearing Arrangement & \fa{} & 
\begin{itemize}[leftmargin=*,nosep]
\item Upper and lower bearing design
\item Bearing housing
\item Lubrication system
\item Seal arrangement
\end{itemize} & NR467 Sec.15, [3.7] \\
\hline

16 & Column Structure Design & \fa{} & 
\begin{itemize}[leftmargin=*,nosep]
\item Column/strut dimensions
\item Material specification
\item Hydrodynamic fairing
\item Structural strength
\item Corrosion protection
\end{itemize} & NR467 Sec.15, [3.8] \\
\hline

17 & Seal System Design & \fa{} & 
\begin{itemize}[leftmargin=*,nosep]
\item Seal type and material
\item Seal arrangement (upper/lower)
\item Pressure rating
\item Leakage monitoring
\item Replacement procedure
\end{itemize} & NR467 Sec.15, [3.9] \\
\hline

\end{longtable}

\subsection{Gear System}

\begin{longtable}{|p{0.5cm}|p{5cm}|p{1.5cm}|p{5.5cm}|p{2cm}|}
\caption{Azimuth Thrusters - Gear System} \label{tab:azimuth-gear} \\
\hline
\rowcolor{bvblue}
\textcolor{white}{\textbf{\#}} & \textcolor{white}{\textbf{Drawing/Document}} & \textcolor{white}{\textbf{Category}} & \textcolor{white}{\textbf{Content Requirements}} & \textcolor{white}{\textbf{Reference}} \\
\hline
\endfirsthead

\multicolumn{5}{c}%
{{\tablename\ \thetable{} -- continued from previous page}} \\
\hline
\rowcolor{bvblue}
\textcolor{white}{\textbf{\#}} & \textcolor{white}{\textbf{Drawing/Document}} & \textcolor{white}{\textbf{Category}} & \textcolor{white}{\textbf{Content Requirements}} & \textcolor{white}{\textbf{Reference}} \\
\hline
\endhead

\hline \multicolumn{5}{|r|}{{Continued on next page}} \\ \hline
\endfoot

\hline
\endlastfoot

18 & Gear Train Assembly & \fa{} & 
\begin{itemize}[leftmargin=*,nosep]
\item Gear arrangement (bevel, helical, planetary)
\item Gear ratios
\item Gear casing design
\item Lubrication system
\item Access for inspection
\end{itemize} & NR467 Sec.15, [3.10] \\
\hline

19 & Gear Tooth Strength Calculations & \fa{} & 
\begin{itemize}[leftmargin=*,nosep]
\item Tooth bending stress (ISO 6336)
\item Contact stress
\item Pitting resistance
\item Scuffing analysis
\item Efficiency calculations
\end{itemize} & NR467 Sec.15, [3.10] \\
\hline

20 & Gear Lubrication System & \fa{} & 
\begin{itemize}[leftmargin=*,nosep]
\item Oil type and capacity
\item Circulation system (forced/splash)
\item Cooling arrangement
\item Filtration system
\item Oil level and temperature monitoring
\end{itemize} & NR467 Sec.15, [3.10] \\
\hline

\end{longtable}

\subsection{Slewing (Azimuth) System}

\begin{longtable}{|p{0.5cm}|p{5cm}|p{1.5cm}|p{5.5cm}|p{2cm}|}
\caption{Azimuth Thrusters - Slewing System} \label{tab:azimuth-slewing} \\
\hline
\rowcolor{bvblue}
\textcolor{white}{\textbf{\#}} & \textcolor{white}{\textbf{Drawing/Document}} & \textcolor{white}{\textbf{Category}} & \textcolor{white}{\textbf{Content Requirements}} & \textcolor{white}{\textbf{Reference}} \\
\hline
\endfirsthead

\multicolumn{5}{c}%
{{\tablename\ \thetable{} -- continued from previous page}} \\
\hline
\rowcolor{bvblue}
\textcolor{white}{\textbf{\#}} & \textcolor{white}{\textbf{Drawing/Document}} & \textcolor{white}{\textbf{Category}} & \textcolor{white}{\textbf{Content Requirements}} & \textcolor{white}{\textbf{Reference}} \\
\hline
\endhead

\hline \multicolumn{5}{|r|}{{Continued on next page}} \\ \hline
\endfoot

\hline
\endlastfoot

21 & Slewing Bearing Design & \fa{} & 
\begin{itemize}[leftmargin=*,nosep]
\item Bearing type and manufacturer
\item Internal/external gear (if applicable)
\item Load capacity (axial, radial, moment)
\item Mounting arrangement
\item Lubrication provisions
\end{itemize} & NR467 Sec.15, [3.11] \\
\hline

22 & Slewing Bearing Calculations & \fa{} & 
\begin{itemize}[leftmargin=*,nosep]
\item Load analysis (static and dynamic)
\item Bearing life calculation
\item Bolt preload calculations
\item Overturning moment analysis
\end{itemize} & NR467 Sec.15, [3.11] \\
\hline

23 & Azimuth Drive Mechanism & \fa{} & 
\begin{itemize}[leftmargin=*,nosep]
\item Drive motor (electric or hydraulic)
\item Gear/pinion arrangement
\item Drive ratio
\item Brake system
\item Locking mechanism (if required)
\end{itemize} & NR467 Sec.15, [3.12] \\
\hline

24 & Azimuth Drive Calculations & \fa{} & 
\begin{itemize}[leftmargin=*,nosep]
\item Torque requirements
\item Acceleration/deceleration analysis
\item Motor sizing verification
\item Brake sizing
\item Response time calculations
\end{itemize} & NR467 Sec.15, [3.12] \\
\hline

25 & Position Feedback System & \fa{} & 
\begin{itemize}[leftmargin=*,nosep]
\item Position sensor type and location
\item Accuracy and resolution
\item Signal processing
\item Redundancy (if required for DP/steering)
\end{itemize} & NR467 Sec.15, [3.13] \\
\hline

\end{longtable}

\subsection{Power Transmission}

\subsubsection{For Electric Drive}

\begin{longtable}{|p{0.5cm}|p{5cm}|p{1.5cm}|p{5.5cm}|p{2cm}|}
\caption{Azimuth Thrusters - Electric Power Transmission} \label{tab:azimuth-electric} \\
\hline
\rowcolor{bvblue}
\textcolor{white}{\textbf{\#}} & \textcolor{white}{\textbf{Drawing/Document}} & \textcolor{white}{\textbf{Category}} & \textcolor{white}{\textbf{Content Requirements}} & \textcolor{white}{\textbf{Reference}} \\
\hline
\endfirsthead

\multicolumn{5}{c}%
{{\tablename\ \thetable{} -- continued from previous page}} \\
\hline
\rowcolor{bvblue}
\textcolor{white}{\textbf{\#}} & \textcolor{white}{\textbf{Drawing/Document}} & \textcolor{white}{\textbf{Category}} & \textcolor{white}{\textbf{Content Requirements}} & \textcolor{white}{\textbf{Reference}} \\
\hline
\endhead

\hline \multicolumn{5}{|r|}{{Continued on next page}} \\ \hline
\endfoot

\hline
\endlastfoot

26 & Electric Motor Specification & \fa{} & 
\begin{itemize}[leftmargin=*,nosep]
\item Motor type and rating
\item Voltage and frequency
\item Cooling method
\item Insulation and protection class
\item Mounting arrangement
\end{itemize} & NR467 Sec.15, [3.14]; IEC 60092-301 \\
\hline

27 & Power Cable Routing & \fa{} & 
\begin{itemize}[leftmargin=*,nosep]
\item Cable route from switchboard
\item Cable support and protection
\item Penetrations through bulkheads
\item Cable sizing calculations
\end{itemize} & NR467 Sec.15, [3.14]; IEC 60092-352 \\
\hline

28 & Slip Ring or Rotary Joint & \fa{} & 
\begin{itemize}[leftmargin=*,nosep]
\item Type and manufacturer
\item Current rating
\item Maintenance requirements
\item Cooling provisions
\end{itemize} & NR467 Sec.15, [3.14] \\
\hline

\end{longtable}

\subsubsection{For Hydraulic Drive}

\begin{longtable}{|p{0.5cm}|p{5cm}|p{1.5cm}|p{5.5cm}|p{2cm}|}
\caption{Azimuth Thrusters - Hydraulic Power Transmission} \label{tab:azimuth-hydraulic} \\
\hline
\rowcolor{bvblue}
\textcolor{white}{\textbf{\#}} & \textcolor{white}{\textbf{Drawing/Document}} & \textcolor{white}{\textbf{Category}} & \textcolor{white}{\textbf{Content Requirements}} & \textcolor{white}{\textbf{Reference}} \\
\hline
\endfirsthead

\multicolumn{5}{c}%
{{\tablename\ \thetable{} -- continued from previous page}} \\
\hline
\rowcolor{bvblue}
\textcolor{white}{\textbf{\#}} & \textcolor{white}{\textbf{Drawing/Document}} & \textcolor{white}{\textbf{Category}} & \textcolor{white}{\textbf{Content Requirements}} & \textcolor{white}{\textbf{Reference}} \\
\hline
\endhead

\hline \multicolumn{5}{|r|}{{Continued on next page}} \\ \hline
\endfoot

\hline
\endlastfoot

29 & Hydraulic Motor Specification & \fa{} & 
\begin{itemize}[leftmargin=*,nosep]
\item Motor type and displacement
\item Rated pressure and speed
\item Efficiency curves
\item Mounting details
\end{itemize} & NR467 Sec.15, [3.15] \\
\hline

30 & Hydraulic System Schematic & \fa{} & 
\begin{itemize}[leftmargin=*,nosep]
\item Complete hydraulic circuit
\item HPU specifications
\item Piping arrangement
\item Control valve system
\item See Chapter \ref{ch:hydraulics} for detailed requirements
\end{itemize} & NR467 Sec.15, [3.15]; Ch.\ref{ch:hydraulics} \\
\hline

\end{longtable}

\subsection{Control System}

\begin{longtable}{|p{0.5cm}|p{5cm}|p{1.5cm}|p{5.5cm}|p{2cm}|}
\caption{Azimuth Thrusters - Control System} \label{tab:azimuth-control} \\
\hline
\rowcolor{bvblue}
\textcolor{white}{\textbf{\#}} & \textcolor{white}{\textbf{Drawing/Document}} & \textcolor{white}{\textbf{Category}} & \textcolor{white}{\textbf{Content Requirements}} & \textcolor{white}{\textbf{Reference}} \\
\hline
\endfirsthead

\multicolumn{5}{c}%
{{\tablename\ \thetable{} -- continued from previous page}} \\
\hline
\rowcolor{bvblue}
\textcolor{white}{\textbf{\#}} & \textcolor{white}{\textbf{Drawing/Document}} & \textcolor{white}{\textbf{Category}} & \textcolor{white}{\textbf{Content Requirements}} & \textcolor{white}{\textbf{Reference}} \\
\hline
\endhead

\hline \multicolumn{5}{|r|}{{Continued on next page}} \\ \hline
\endfoot

\hline
\endlastfoot

31 & Control System Architecture & \fa{} & 
\begin{itemize}[leftmargin=*,nosep]
\item Overall control system topology
\item Propulsion control
\item Azimuth control
\item Pitch control (CP thrusters)
\item Interface with ship systems
\end{itemize} & NR467 Sec.15, [3.16]; Part C, Ch.4 \\
\hline

32 & Bridge Control Console & \fa{} & 
\begin{itemize}[leftmargin=*,nosep]
\item Control panel layout
\item Joystick/lever arrangement
\item Display and indication
\item Alarm system
\item Emergency stop provisions
\end{itemize} & NR467 Sec.15, [3.16] \\
\hline

33 & Local Control Station & \fa{} & 
\begin{itemize}[leftmargin=*,nosep]
\item Local control panel design
\item Control mode selection
\item Indication and monitoring
\item Interlock with bridge control
\end{itemize} & NR467 Sec.15, [3.16] \\
\hline

34 & Control Software Description & \fa{} & 
\begin{itemize}[leftmargin=*,nosep]
\item Software architecture
\item Control algorithms
\item Safety functions
\item Software validation documentation
\item Version control
\end{itemize} & NR467 Part C, Ch.4 \\
\hline

35 & Electrical Single-Line Diagram & \fa{} & 
\begin{itemize}[leftmargin=*,nosep]
\item Power supply arrangement
\item Circuit breakers and protection
\item Emergency power provisions
\item Monitoring and alarm circuits
\end{itemize} & NR467 Sec.15, [3.16]; IEC 60092-504 \\
\hline

\end{longtable}

\subsection{Additional Calculations \& Documentation}

\begin{longtable}{|p{0.5cm}|p{5cm}|p{1.5cm}|p{5.5cm}|p{2cm}|}
\caption{Azimuth Thrusters - Additional Calculations} \label{tab:azimuth-calc} \\
\hline
\rowcolor{bvblue}
\textcolor{white}{\textbf{\#}} & \textcolor{white}{\textbf{Drawing/Document}} & \textcolor{white}{\textbf{Category}} & \textcolor{white}{\textbf{Content Requirements}} & \textcolor{white}{\textbf{Reference}} \\
\hline
\endfirsthead

\multicolumn{5}{c}%
{{\tablename\ \thetable{} -- continued from previous page}} \\
\hline
\rowcolor{bvblue}
\textcolor{white}{\textbf{\#}} & \textcolor{white}{\textbf{Drawing/Document}} & \textcolor{white}{\textbf{Category}} & \textcolor{white}{\textbf{Content Requirements}} & \textcolor{white}{\textbf{Reference}} \\
\hline
\endhead

\hline \multicolumn{5}{|r|}{{Continued on next page}} \\ \hline
\endfoot

\hline
\endlastfoot

36 & Thrust and Torque Calculations & \fa{} & 
\begin{itemize}[leftmargin=*,nosep]
\item Propeller thrust curves
\item Torque requirements
\item Power absorption
\item Efficiency analysis
\item Bollard pull prediction
\end{itemize} & NR467 Sec.15, [3.17] \\
\hline

37 & Vibration Analysis & \fa{} & 
\begin{itemize}[leftmargin=*,nosep]
\item Natural frequency calculations
\item Forced vibration analysis
\item Resonance avoidance
\item Foundation vibration isolation
\item Acceptable vibration levels
\end{itemize} & NR467 Sec.15, [3.18] \\
\hline

38 & Noise Analysis & \fa{} & 
\begin{itemize}[leftmargin=*,nosep]
\item Airborne noise prediction
\item Structure-borne noise analysis
\item Compliance with noise regulations
\item Noise reduction measures
\end{itemize} & NR467 Sec.15, [3.18] \\
\hline

39 & Hydrodynamic Analysis & \fa{} & 
\begin{itemize}[leftmargin=*,nosep]
\item Propeller hydrodynamic design
\item Cavitation analysis
\item Efficiency optimization
\item Interaction effects with hull
\end{itemize} & NR467 Sec.15, [3.19] \\
\hline

40 & Thermal Analysis & \fa{} & 
\begin{itemize}[leftmargin=*,nosep]
\item Heat generation in gears, bearings, seals
\item Cooling system capacity
\item Temperature rise predictions
\item Ventilation requirements
\end{itemize} & NR467 Sec.15, [3.20] \\
\hline

\end{longtable}

\newpage

\section{Podded Propulsion Units}

\textbf{Podded propulsion units include all requirements from Section 3.1 (Azimuth Thrusters), with the following ADDITIONAL or MODIFIED requirements:}

\subsection{Pod-Specific Structural Components}

\begin{longtable}{|p{0.5cm}|p{5cm}|p{1.5cm}|p{5.5cm}|p{2cm}|}
\caption{Podded Propulsion - Pod Structure} \label{tab:pod-structure} \\
\hline
\rowcolor{bvblue}
\textcolor{white}{\textbf{\#}} & \textcolor{white}{\textbf{Drawing/Document}} & \textcolor{white}{\textbf{Category}} & \textcolor{white}{\textbf{Content Requirements}} & \textcolor{white}{\textbf{Reference}} \\
\hline
\endfirsthead

\multicolumn{5}{c}%
{{\tablename\ \thetable{} -- continued from previous page}} \\
\hline
\rowcolor{bvblue}
\textcolor{white}{\textbf{\#}} & \textcolor{white}{\textbf{Drawing/Document}} & \textcolor{white}{\textbf{Category}} & \textcolor{white}{\textbf{Content Requirements}} & \textcolor{white}{\textbf{Reference}} \\
\hline
\endhead

\hline \multicolumn{5}{|r|}{{Continued on next page}} \\ \hline
\endfoot

\hline
\endlastfoot

41 & Pod Housing Design & \fa{} & 
\begin{itemize}[leftmargin=*,nosep]
\item Pod shell structure
\item Material specification
\item Hydrodynamic fairing
\item Structural strength analysis
\item Corrosion protection system
\end{itemize} & NR467 Sec.15, [3.21] \\
\hline

42 & Pod Housing Calculations & \fa{} & 
\begin{itemize}[leftmargin=*,nosep]
\item Pressure loading analysis
\item Structural stress analysis
\item Fatigue assessment
\item Impact resistance
\item FEA results
\end{itemize} & NR467 Sec.15, [3.21] \\
\hline

43 & Strut Structure Design & \fa{} & 
\begin{itemize}[leftmargin=*,nosep]
\item Strut dimensions and material
\item Hydrodynamic profile
\item Connection to pod and hull
\item Internal cable/pipe routing
\item Access provisions
\end{itemize} & NR467 Sec.15, [3.22] \\
\hline

44 & Strut Structural Calculations & \fa{} & 
\begin{itemize}[leftmargin=*,nosep]
\item Bending stress analysis
\item Torsional stress
\item Combined loading
\item Fatigue analysis
\item Vibration analysis
\end{itemize} & NR467 Sec.15, [3.22] \\
\hline

\end{longtable}

\subsection{Electric Motor in Pod}

\begin{longtable}{|p{0.5cm}|p{5cm}|p{1.5cm}|p{5.5cm}|p{2cm}|}
\caption{Podded Propulsion - In-Pod Motor} \label{tab:pod-motor} \\
\hline
\rowcolor{bvblue}
\textcolor{white}{\textbf{\#}} & \textcolor{white}{\textbf{Drawing/Document}} & \textcolor{white}{\textbf{Category}} & \textcolor{white}{\textbf{Content Requirements}} & \textcolor{white}{\textbf{Reference}} \\
\hline
\endfirsthead

\multicolumn{5}{c}%
{{\tablename\ \thetable{} -- continued from previous page}} \\
\hline
\rowcolor{bvblue}
\textcolor{white}{\textbf{\#}} & \textcolor{white}{\textbf{Drawing/Document}} & \textcolor{white}{\textbf{Category}} & \textcolor{white}{\textbf{Content Requirements}} & \textcolor{white}{\textbf{Reference}} \\
\hline
\endhead

\hline \multicolumn{5}{|r|}{{Continued on next page}} \\ \hline
\endfoot

\hline
\endlastfoot

45 & In-Pod Motor Design & \fa{} & 
\begin{itemize}[leftmargin=*,nosep]
\item Motor type (permanent magnet, induction)
\item Rated power and speed
\item Voltage and frequency
\item Insulation system (marine environment)
\item Rotor and stator assembly
\item Bearing arrangement
\end{itemize} & NR467 Sec.15, [3.23]; IEC 60092-301 \\
\hline

46 & Motor Thermal Design & \fa{} & 
\begin{itemize}[leftmargin=*,nosep]
\item Heat generation calculations
\item Cooling system design
\item Temperature limits
\item Thermal monitoring
\item Insulation life assessment
\end{itemize} & NR467 Sec.15, [3.23] \\
\hline

47 & Motor Cooling System & \fa{} & 
\begin{itemize}[leftmargin=*,nosep]
\item Cooling method (water-cooled, oil-cooled)
\item Heat exchanger design
\item Coolant circulation system
\item Pump specifications
\item Temperature monitoring
\end{itemize} & NR467 Sec.15, [3.24] \\
\hline

48 & Motor Mounting \& Support & \fa{} & 
\begin{itemize}[leftmargin=*,nosep]
\item Motor mounting arrangement in pod
\item Alignment provisions
\item Vibration isolation
\item Thermal expansion accommodation
\end{itemize} & NR467 Sec.15, [3.23] \\
\hline

\end{longtable}

\subsection{Power Supply Through Strut}

\begin{longtable}{|p{0.5cm}|p{5cm}|p{1.5cm}|p{5.5cm}|p{2cm}|}
\caption{Podded Propulsion - Power Supply} \label{tab:pod-power} \\
\hline
\rowcolor{bvblue}
\textcolor{white}{\textbf{\#}} & \textcolor{white}{\textbf{Drawing/Document}} & \textcolor{white}{\textbf{Category}} & \textcolor{white}{\textbf{Content Requirements}} & \textcolor{white}{\textbf{Reference}} \\
\hline
\endfirsthead

\multicolumn{5}{c}%
{{\tablename\ \thetable{} -- continued from previous page}} \\
\hline
\rowcolor{bvblue}
\textcolor{white}{\textbf{\#}} & \textcolor{white}{\textbf{Drawing/Document}} & \textcolor{white}{\textbf{Category}} & \textcolor{white}{\textbf{Content Requirements}} & \textcolor{white}{\textbf{Reference}} \\
\hline
\endhead

\hline \multicolumn{5}{|r|}{{Continued on next page}} \\ \hline
\endfoot

\hline
\endlastfoot

49 & Power Cable Routing & \fa{} & 
\begin{itemize}[leftmargin=*,nosep]
\item Cable route through strut
\item Cable support and protection
\item Cable entry into pod
\item Sealing arrangements
\item Cable sizing calculations
\end{itemize} & NR467 Sec.15, [3.25]; IEC 60092-352 \\
\hline

50 & Slip Ring or Rotary Transformer & \fa{} & 
\begin{itemize}[leftmargin=*,nosep]
\item Type and manufacturer
\item Current and voltage rating
\item Cooling provisions
\item Maintenance access
\item Reliability data
\end{itemize} & NR467 Sec.15, [3.25] \\
\hline

51 & Power Distribution in Pod & \fa{} & 
\begin{itemize}[leftmargin=*,nosep]
\item Internal switchgear/junction boxes
\item Protection devices
\item Monitoring systems
\item Watertight/waterproof rating
\end{itemize} & NR467 Sec.15, [3.25] \\
\hline

\end{longtable}

\subsection{Propeller Direct Drive}

\begin{longtable}{|p{0.5cm}|p{5cm}|p{1.5cm}|p{5.5cm}|p{2cm}|}
\caption{Podded Propulsion - Direct Drive} \label{tab:pod-drive} \\
\hline
\rowcolor{bvblue}
\textcolor{white}{\textbf{\#}} & \textcolor{white}{\textbf{Drawing/Document}} & \textcolor{white}{\textbf{Category}} & \textcolor{white}{\textbf{Content Requirements}} & \textcolor{white}{\textbf{Reference}} \\
\hline
\endfirsthead

\multicolumn{5}{c}%
{{\tablename\ \thetable{} -- continued from previous page}} \\
\hline
\rowcolor{bvblue}
\textcolor{white}{\textbf{\#}} & \textcolor{white}{\textbf{Drawing/Document}} & \textcolor{white}{\textbf{Category}} & \textcolor{white}{\textbf{Content Requirements}} & \textcolor{white}{\textbf{Reference}} \\
\hline
\endhead

\hline \multicolumn{5}{|r|}{{Continued on next page}} \\ \hline
\endfoot

\hline
\endlastfoot

52 & Direct Drive Arrangement & \fa{} & 
\begin{itemize}[leftmargin=*,nosep]
\item Motor shaft to propeller connection
\item Coupling or direct mounting
\item Bearing arrangement
\item Seal system
\item Alignment procedure
\end{itemize} & NR467 Sec.15, [3.26] \\
\hline

53 & Shaft Seal System & \fa{} & 
\begin{itemize}[leftmargin=*,nosep]
\item Seal type and material
\item Seal arrangement
\item Pressure rating
\item Leakage monitoring
\item Emergency sealing provisions
\end{itemize} & NR467 Sec.15, [3.26] \\
\hline

\end{longtable}

\subsection{Pod-Specific Testing}

\begin{longtable}{|p{0.5cm}|p{5cm}|p{1.5cm}|p{5.5cm}|p{2cm}|}
\caption{Podded Propulsion - Testing Requirements} \label{tab:pod-testing} \\
\hline
\rowcolor{bvblue}
\textcolor{white}{\textbf{\#}} & \textcolor{white}{\textbf{Drawing/Document}} & \textcolor{white}{\textbf{Category}} & \textcolor{white}{\textbf{Content Requirements}} & \textcolor{white}{\textbf{Reference}} \\
\hline
\endfirsthead

\multicolumn{5}{c}%
{{\tablename\ \thetable{} -- continued from previous page}} \\
\hline
\rowcolor{bvblue}
\textcolor{white}{\textbf{\#}} & \textcolor{white}{\textbf{Drawing/Document}} & \textcolor{white}{\textbf{Category}} & \textcolor{white}{\textbf{Content Requirements}} & \textcolor{white}{\textbf{Reference}} \\
\hline
\endhead

\hline \multicolumn{5}{|r|}{{Continued on next page}} \\ \hline
\endfoot

\hline
\endlastfoot

54 & Pressure Testing Procedure & \fa{} & 
\begin{itemize}[leftmargin=*,nosep]
\item Test pressure and duration
\item Test setup
\item Acceptance criteria
\item Leakage limits
\end{itemize} & NR467 Sec.15, [3.27] \\
\hline

55 & Motor Insulation Testing & \fa{} & 
\begin{itemize}[leftmargin=*,nosep]
\item Insulation resistance tests
\item High voltage tests
\item Partial discharge tests
\item Test conditions and acceptance criteria
\end{itemize} & IEC 60092-301 \\
\hline

\end{longtable}

\newpage

% Chapter 4: Water-Jets
\chapter{Water-Jet Propulsion Systems}
\label{ch:waterjets}

\section{General Arrangement}

\begin{longtable}{|p{0.5cm}|p{5cm}|p{1.5cm}|p{5.5cm}|p{2cm}|}
\caption{Water-Jets - General Arrangement} \label{tab:waterjet-general} \\
\hline
\rowcolor{bvblue}
\textcolor{white}{\textbf{\#}} & \textcolor{white}{\textbf{Drawing/Document}} & \textcolor{white}{\textbf{Category}} & \textcolor{white}{\textbf{Content Requirements}} & \textcolor{white}{\textbf{Reference}} \\
\hline
\endfirsthead

\multicolumn{5}{c}%
{{\tablename\ \thetable{} -- continued from previous page}} \\
\hline
\rowcolor{bvblue}
\textcolor{white}{\textbf{\#}} & \textcolor{white}{\textbf{Drawing/Document}} & \textcolor{white}{\textbf{Category}} & \textcolor{white}{\textbf{Content Requirements}} & \textcolor{white}{\textbf{Reference}} \\
\hline
\endhead

\hline \multicolumn{5}{|r|}{{Continued on next page}} \\ \hline
\endfoot

\hline
\endlastfoot

1 & Water-Jet General Arrangement & \fa{} & 
\begin{itemize}[leftmargin=*,nosep]
\item Installation location on vessel
\item Inlet and outlet positions
\item Nozzle orientation
\item Access for maintenance
\item Clearances
\end{itemize} & NR467 Sec.15, [4.1] \\
\hline

2 & Hull Integration Drawing & \fa{} & 
\begin{itemize}[leftmargin=*,nosep]
\item Inlet duct design
\item Hull penetration details
\item Structural integration
\item Hydrodynamic considerations
\item Reinforcement arrangements
\end{itemize} & NR467 Sec.15, [4.2]; Part B, Ch.5 \\
\hline

\end{longtable}

\section{Water-Jet Unit Components}

\begin{longtable}{|p{0.5cm}|p{5cm}|p{1.5cm}|p{5.5cm}|p{2cm}|}
\caption{Water-Jets - Unit Components} \label{tab:waterjet-components} \\
\hline
\rowcolor{bvblue}
\textcolor{white}{\textbf{\#}} & \textcolor{white}{\textbf{Drawing/Document}} & \textcolor{white}{\textbf{Category}} & \textcolor{white}{\textbf{Content Requirements}} & \textcolor{white}{\textbf{Reference}} \\
\hline
\endfirsthead

\multicolumn{5}{c}%
{{\tablename\ \thetable{} -- continued from previous page}} \\
\hline
\rowcolor{bvblue}
\textcolor{white}{\textbf{\#}} & \textcolor{white}{\textbf{Drawing/Document}} & \textcolor{white}{\textbf{Category}} & \textcolor{white}{\textbf{Content Requirements}} & \textcolor{white}{\textbf{Reference}} \\
\hline
\endhead

\hline \multicolumn{5}{|r|}{{Continued on next page}} \\ \hline
\endfoot

\hline
\endlastfoot

3 & Water-Jet Unit Assembly & \fa{} & 
\begin{itemize}[leftmargin=*,nosep]
\item Complete unit cross-section
\item Pump assembly
\item Inlet duct
\item Nozzle and steering/reversing mechanism
\item Material specifications
\end{itemize} & NR467 Sec.15, [4.3] \\
\hline

4 & Impeller Design & \fa{} & 
\begin{itemize}[leftmargin=*,nosep]
\item Impeller geometry
\item Number of blades
\item Material specification
\item Blade profile
\item Hub design
\end{itemize} & NR467 Sec.15, [4.4] \\
\hline

5 & Impeller Strength Calculations & \fa{} & 
\begin{itemize}[leftmargin=*,nosep]
\item Blade stress analysis
\item Hub stress analysis
\item Centrifugal forces
\item Hydrodynamic forces
\item Fatigue assessment
\item Cavitation analysis
\end{itemize} & NR467 Sec.15, [4.4] \\
\hline

6 & Pump Shaft Design & \fa{} & 
\begin{itemize}[leftmargin=*,nosep]
\item Shaft dimensions and material
\item Bearing locations
\item Seal arrangements
\item Coupling details
\end{itemize} & NR467 Sec.15, [4.5] \\
\hline

7 & Pump Shaft Calculations & \fa{} & 
\begin{itemize}[leftmargin=*,nosep]
\item Torsional strength
\item Bending stress
\item Combined stress analysis
\item Deflection
\item Critical speed
\item Fatigue assessment
\end{itemize} & NR467 Sec.15, [4.5]; Sec.7 \\
\hline

8 & Bearing Arrangement & \fa{} & 
\begin{itemize}[leftmargin=*,nosep]
\item Bearing types and locations
\item Bearing housing
\item Lubrication system
\item Seal design
\item Cooling provisions
\end{itemize} & NR467 Sec.15, [4.6] \\
\hline

9 & Bearing Load Calculations & \fa{} & 
\begin{itemize}[leftmargin=*,nosep]
\item Radial and axial loads
\item Bearing life calculation
\item Lubrication requirements
\item Temperature analysis
\end{itemize} & NR467 Sec.15, [4.6] \\
\hline

\end{longtable}

\section{Inlet Duct System}

\begin{longtable}{|p{0.5cm}|p{5cm}|p{1.5cm}|p{5.5cm}|p{2cm}|}
\caption{Water-Jets - Inlet Duct} \label{tab:waterjet-inlet} \\
\hline
\rowcolor{bvblue}
\textcolor{white}{\textbf{\#}} & \textcolor{white}{\textbf{Drawing/Document}} & \textcolor{white}{\textbf{Category}} & \textcolor{white}{\textbf{Content Requirements}} & \textcolor{white}{\textbf{Reference}} \\
\hline
\endfirsthead

\multicolumn{5}{c}%
{{\tablename\ \thetable{} -- continued from previous page}} \\
\hline
\rowcolor{bvblue}
\textcolor{white}{\textbf{\#}} & \textcolor{white}{\textbf{Drawing/Document}} & \textcolor{white}{\textbf{Category}} & \textcolor{white}{\textbf{Content Requirements}} & \textcolor{white}{\textbf{Reference}} \\
\hline
\endhead

\hline \multicolumn{5}{|r|}{{Continued on next page}} \\ \hline
\endfoot

\hline
\endlastfoot

10 & Inlet Duct Design & \fa{} & 
\begin{itemize}[leftmargin=*,nosep]
\item Duct geometry and dimensions
\item Material specification
\item Structural design
\item Hydrodynamic profile
\item Debris protection (grating/screen)
\end{itemize} & NR467 Sec.15, [4.7] \\
\hline

11 & Inlet Duct Calculations & \fa{} & 
\begin{itemize}[leftmargin=*,nosep]
\item Structural strength analysis
\item Pressure loading
\item Flow analysis
\item Cavitation assessment
\end{itemize} & NR467 Sec.15, [4.7] \\
\hline

12 & Hull Penetration \& Reinforcement & \fa{} & 
\begin{itemize}[leftmargin=*,nosep]
\item Opening dimensions
\item Reinforcement design
\item Watertight integrity
\item Structural attachment
\end{itemize} & NR467 Sec.15, [4.7]; Part B, Ch.5 \\
\hline

\end{longtable}

\section{Steering and Reversing Mechanism}

\begin{longtable}{|p{0.5cm}|p{5cm}|p{1.5cm}|p{5.5cm}|p{2cm}|}
\caption{Water-Jets - Steering \& Reversing} \label{tab:waterjet-steering} \\
\hline
\rowcolor{bvblue}
\textcolor{white}{\textbf{\#}} & \textcolor{white}{\textbf{Drawing/Document}} & \textcolor{white}{\textbf{Category}} & \textcolor{white}{\textbf{Content Requirements}} & \textcolor{white}{\textbf{Reference}} \\
\hline
\endfirsthead

\multicolumn{5}{c}%
{{\tablename\ \thetable{} -- continued from previous page}} \\
\hline
\rowcolor{bvblue}
\textcolor{white}{\textbf{\#}} & \textcolor{white}{\textbf{Drawing/Document}} & \textcolor{white}{\textbf{Category}} & \textcolor{white}{\textbf{Content Requirements}} & \textcolor{white}{\textbf{Reference}} \\
\hline
\endhead

\hline \multicolumn{5}{|r|}{{Continued on next page}} \\ \hline
\endfoot

\hline
\endlastfoot

13 & Steering Nozzle Design & \fa{} & 
\begin{itemize}[leftmargin=*,nosep]
\item Nozzle geometry
\item Deflection range
\item Actuator arrangement
\item Structural design
\item Material specification
\end{itemize} & NR467 Sec.15, [4.8] \\
\hline

14 & Reversing Bucket Design & \fa{} & 
\begin{itemize}[leftmargin=*,nosep]
\item Bucket geometry
\item Deployment mechanism
\item Actuator system
\item Structural strength
\item Material specification
\end{itemize} & NR467 Sec.15, [4.9] \\
\hline

15 & Steering/Reversing Actuators & \fa{} & 
\begin{itemize}[leftmargin=*,nosep]
\item Hydraulic cylinder specifications
\item Operating pressure
\item Force requirements
\item Position feedback
\item Emergency provisions
\end{itemize} & NR467 Sec.15, [4.10] \\
\hline

16 & Steering/Reversing Calculations & \fa{} & 
\begin{itemize}[leftmargin=*,nosep]
\item Actuator force requirements
\item Structural strength of nozzle/bucket
\item Response time analysis
\item Hydraulic system sizing
\end{itemize} & NR467 Sec.15, [4.10] \\
\hline

\end{longtable}

\section{Drive System}

\subsection{For Direct Drive from Main Engine}

\begin{longtable}{|p{0.5cm}|p{5cm}|p{1.5cm}|p{5.5cm}|p{2cm}|}
\caption{Water-Jets - Direct Drive} \label{tab:waterjet-direct} \\
\hline
\rowcolor{bvblue}
\textcolor{white}{\textbf{\#}} & \textcolor{white}{\textbf{Drawing/Document}} & \textcolor{white}{\textbf{Category}} & \textcolor{white}{\textbf{Content Requirements}} & \textcolor{white}{\textbf{Reference}} \\
\hline
\endfirsthead

\multicolumn{5}{c}%
{{\tablename\ \thetable{} -- continued from previous page}} \\
\hline
\rowcolor{bvblue}
\textcolor{white}{\textbf{\#}} & \textcolor{white}{\textbf{Drawing/Document}} & \textcolor{white}{\textbf{Category}} & \textcolor{white}{\textbf{Content Requirements}} & \textcolor{white}{\textbf{Reference}} \\
\hline
\endhead

\hline \multicolumn{5}{|r|}{{Continued on next page}} \\ \hline
\endfoot

\hline
\endlastfoot

17 & Drive Shaft Arrangement & \fa{} & 
\begin{itemize}[leftmargin=*,nosep]
\item Shaft line from engine to water-jet
\item Bearing locations
\item Coupling arrangements
\item Alignment provisions
\end{itemize} & NR467 Sec.15, [4.11]; Sec.7 \\
\hline

18 & Drive Shaft Calculations & \fa{} & 
\begin{itemize}[leftmargin=*,nosep]
\item Torsional strength
\item Bending stress
\item Critical speed analysis
\item Alignment calculations
\item See Chapter \ref{ch:shafting} for full requirements
\end{itemize} & NR467 Sec.15, [4.11]; Sec.7 \\
\hline

\end{longtable}

\subsection{For Gearbox Drive}

\begin{longtable}{|p{0.5cm}|p{5cm}|p{1.5cm}|p{5.5cm}|p{2cm}|}
\caption{Water-Jets - Gearbox Drive} \label{tab:waterjet-gearbox} \\
\hline
\rowcolor{bvblue}
\textcolor{white}{\textbf{\#}} & \textcolor{white}{\textbf{Drawing/Document}} & \textcolor{white}{\textbf{Category}} & \textcolor{white}{\textbf{Content Requirements}} & \textcolor{white}{\textbf{Reference}} \\
\hline
\endfirsthead

\multicolumn{5}{c}%
{{\tablename\ \thetable{} -- continued from previous page}} \\
\hline
\rowcolor{bvblue}
\textcolor{white}{\textbf{\#}} & \textcolor{white}{\textbf{Drawing/Document}} & \textcolor{white}{\textbf{Category}} & \textcolor{white}{\textbf{Content Requirements}} & \textcolor{white}{\textbf{Reference}} \\
\hline
\endhead

\hline \multicolumn{5}{|r|}{{Continued on next page}} \\ \hline
\endfoot

\hline
\endlastfoot

19 & Gearbox Assembly & \fa{} & 
\begin{itemize}[leftmargin=*,nosep]
\item Gear arrangement
\item Gear ratios
\item Casing design
\item Lubrication system
\item Mounting arrangement
\end{itemize} & NR467 Sec.15, [4.12] \\
\hline

20 & Gear Calculations & \fa{} & 
\begin{itemize}[leftmargin=*,nosep]
\item Gear tooth strength
\item Contact stress
\item Bending stress
\item Lubrication analysis
\item Efficiency
\end{itemize} & NR467 Sec.15, [4.12] \\
\hline

\end{longtable}

\section{Hydraulic System (for steering/reversing)}

\begin{longtable}{|p{0.5cm}|p{5cm}|p{1.5cm}|p{5.5cm}|p{2cm}|}
\caption{Water-Jets - Hydraulic System} \label{tab:waterjet-hydraulic} \\
\hline
\rowcolor{bvblue}
\textcolor{white}{\textbf{\#}} & \textcolor{white}{\textbf{Drawing/Document}} & \textcolor{white}{\textbf{Category}} & \textcolor{white}{\textbf{Content Requirements}} & \textcolor{white}{\textbf{Reference}} \\
\hline
\endfirsthead

\multicolumn{5}{c}%
{{\tablename\ \thetable{} -- continued from previous page}} \\
\hline
\rowcolor{bvblue}
\textcolor{white}{\textbf{\#}} & \textcolor{white}{\textbf{Drawing/Document}} & \textcolor{white}{\textbf{Category}} & \textcolor{white}{\textbf{Content Requirements}} & \textcolor{white}{\textbf{Reference}} \\
\hline
\endhead

\hline \multicolumn{5}{|r|}{{Continued on next page}} \\ \hline
\endfoot

\hline
\endlastfoot

21 & Hydraulic System Schematic & \fa{} & 
\begin{itemize}[leftmargin=*,nosep]
\item Complete hydraulic circuit
\item HPU specifications
\item Control valves
\item Piping arrangement
\item See Chapter \ref{ch:hydraulics} for detailed requirements
\end{itemize} & NR467 Sec.15, [4.13]; Ch.\ref{ch:hydraulics} \\
\hline

\end{longtable}

\section{Control System}

\begin{longtable}{|p{0.5cm}|p{5cm}|p{1.5cm}|p{5.5cm}|p{2cm}|}
\caption{Water-Jets - Control System} \label{tab:waterjet-control} \\
\hline
\rowcolor{bvblue}
\textcolor{white}{\textbf{\#}} & \textcolor{white}{\textbf{Drawing/Document}} & \textcolor{white}{\textbf{Category}} & \textcolor{white}{\textbf{Content Requirements}} & \textcolor{white}{\textbf{Reference}} \\
\hline
\endfirsthead

\multicolumn{5}{c}%
{{\tablename\ \thetable{} -- continued from previous page}} \\
\hline
\rowcolor{bvblue}
\textcolor{white}{\textbf{\#}} & \textcolor{white}{\textbf{Drawing/Document}} & \textcolor{white}{\textbf{Category}} & \textcolor{white}{\textbf{Content Requirements}} & \textcolor{white}{\textbf{Reference}} \\
\hline
\endhead

\hline \multicolumn{5}{|r|}{{Continued on next page}} \\ \hline
\endfoot

\hline
\endlastfoot

22 & Control System Architecture & \fa{} & 
\begin{itemize}[leftmargin=*,nosep]
\item Control system topology
\item Steering control
\item Reversing control
\item Engine/pump speed control
\item Interface with ship systems
\end{itemize} & NR467 Sec.15, [4.14]; Part C, Ch.4 \\
\hline

23 & Bridge Control Console & \fa{} & 
\begin{itemize}[leftmargin=*,nosep]
\item Control panel layout
\item Joystick/lever arrangement
\item Display and indication
\item Alarm system
\item Emergency stop
\end{itemize} & NR467 Sec.15, [4.14] \\
\hline

24 & Electrical Single-Line Diagram & \fa{} & 
\begin{itemize}[leftmargin=*,nosep]
\item Power supply
\item Protection devices
\item Monitoring circuits
\item Emergency power provisions
\end{itemize} & NR467 Sec.15, [4.14]; IEC 60092-504 \\
\hline

\end{longtable}

\section{Performance Calculations}

\begin{longtable}{|p{0.5cm}|p{5cm}|p{1.5cm}|p{5.5cm}|p{2cm}|}
\caption{Water-Jets - Performance Calculations} \label{tab:waterjet-performance} \\
\hline
\rowcolor{bvblue}
\textcolor{white}{\textbf{\#}} & \textcolor{white}{\textbf{Drawing/Document}} & \textcolor{white}{\textbf{Category}} & \textcolor{white}{\textbf{Content Requirements}} & \textcolor{white}{\textbf{Reference}} \\
\hline
\endfirsthead

\multicolumn{5}{c}%
{{\tablename\ \thetable{} -- continued from previous page}} \\
\hline
\rowcolor{bvblue}
\textcolor{white}{\textbf{\#}} & \textcolor{white}{\textbf{Drawing/Document}} & \textcolor{white}{\textbf{Category}} & \textcolor{white}{\textbf{Content Requirements}} & \textcolor{white}{\textbf{Reference}} \\
\hline
\endhead

\hline \multicolumn{5}{|r|}{{Continued on next page}} \\ \hline
\endfoot

\hline
\endlastfoot

25 & Thrust and Performance Calculations & \fa{} & 
\begin{itemize}[leftmargin=*,nosep]
\item Thrust curves
\item Power absorption
\item Efficiency analysis
\item Speed predictions
\item Bollard pull
\end{itemize} & NR467 Sec.15, [4.15] \\
\hline

26 & Cavitation Analysis & \fa{} & 
\begin{itemize}[leftmargin=*,nosep]
\item Cavitation inception speed
\item Cavitation patterns
\item Erosion risk assessment
\item Mitigation measures
\end{itemize} & NR467 Sec.15, [4.15] \\
\hline

27 & Vibration Analysis & \fa{} & 
\begin{itemize}[leftmargin=*,nosep]
\item Natural frequency calculations
\item Forced vibration analysis
\item Resonance avoidance
\item Foundation design
\end{itemize} & NR467 Sec.15, [4.16] \\
\hline

28 & Noise Analysis & \fa{} & 
\begin{itemize}[leftmargin=*,nosep]
\item Airborne noise prediction
\item Structure-borne noise
\item Compliance with regulations
\item Noise reduction measures
\end{itemize} & NR467 Sec.15, [4.16] \\
\hline

\end{longtable}

\newpage

% Chapter: Voith-Schneider / Cycloidal
\chapter{Voith-Schneider (Cycloidal) Propulsion Systems}
\label{ch:vsp}

\section{Description \& Classification Basis}

The Voith-Schneider Propeller (VSP), or cycloidal propeller, is a vertical-axis
propulsor in which a set of vertical blades projects below a rotating casing
(rotor). Each blade oscillates about its own vertical axis through a kinematic
control gear linked to a movable control (steering) centre. By displacing the
control centre, the unit produces thrust of \emph{any magnitude in any
horizontal direction} at constant rotor speed. A VSP therefore combines
\textbf{propulsion and steering in a single unit}; the steering function makes
\textbf{NR467 Part C, Ch 1, Sec 14 (Steering Gear) and SOLAS II-1 Reg.\,29
requirements applicable} (see Chapter \ref{ch:steering}). Typical applications:
Voith Water Tractors, double-ended ferries, buoy tenders, mine-countermeasure
vessels.

VSP units are assessed under \textbf{NR467 Pt C, Ch 1, Sec 15} for the
mechanical thruster elements, \textbf{Sec 14} for the steering/control function,
and \textbf{Sec 7} where the drive forms part of the propulsion shaft line.

\section{VSP-Specific Structural \& Mechanical Components}

\begin{longtable}{|p{0.5cm}|p{5cm}|p{1.5cm}|p{5.5cm}|p{2cm}|}
\caption{Voith-Schneider Propeller --- Mechanical \& Structural Requirements} \label{tab:vsp-mech} \\
\hline
\rowcolor{bvblue}
\textcolor{white}{\textbf{\#}} & \textcolor{white}{\textbf{Drawing/Document}} & \textcolor{white}{\textbf{Category}} & \textcolor{white}{\textbf{Content Requirements}} & \textcolor{white}{\textbf{Reference}} \\
\hline
\endfirsthead
\multicolumn{5}{c}%
{{\tablename\ \thetable{} -- continued from previous page}} \\
\hline
\rowcolor{bvblue}
\textcolor{white}{\textbf{\#}} & \textcolor{white}{\textbf{Drawing/Document}} & \textcolor{white}{\textbf{Category}} & \textcolor{white}{\textbf{Content Requirements}} & \textcolor{white}{\textbf{Reference}} \\
\hline
\endhead
\hline \multicolumn{5}{|r|}{{Continued on next page}} \\ \hline
\endfoot
\hline
\endlastfoot

1 & VSP Unit General Assembly & \fa{} &
\begin{itemize}[leftmargin=*,nosep]
\item Complete unit cross-section (rotor casing, blades, gear)
\item Number of blades and blade circle diameter
\item Rotor speed and orbit geometry
\item Material list and lifting arrangement
\end{itemize} & NR467 Sec.15, [5.1] \\
\hline

2 & Blade (Wing) Design Drawing & \fa{} &
\begin{itemize}[leftmargin=*,nosep]
\item Blade profile, chord, span
\item Blade shaft dimensions and material
\item Blade root / shaft connection
\item Surface finish and corrosion protection
\end{itemize} & NR467 Sec.15, [5.2] \\
\hline

3 & Blade Shaft \& Blade Strength Calculations & \fa{} &
\begin{itemize}[leftmargin=*,nosep]
\item Hydrodynamic blade loads (thrust + side force)
\item Bending + torsion of blade shaft
\item Centrifugal loads at max rotor speed
\item Fatigue assessment (reversing loads)
\end{itemize} & NR467 Sec.15, [5.2]; Sec.7 \\
\hline

4 & Blade Bearing Arrangement & \fa{} &
\begin{itemize}[leftmargin=*,nosep]
\item Upper/lower blade bearings, seals
\item Bearing load and L10 life
\item Lubrication path from oil system
\end{itemize} & NR467 Sec.15, [5.3] \\
\hline

5 & Kinematic Control Gear (Steering Mechanism) & \fa{} &
\begin{itemize}[leftmargin=*,nosep]
\item Control-centre / eccentric arrangement
\item Control rods, links, crank geometry
\item Servo actuator interface
\item Strength of linkage members and pins
\item Fail-safe / neutral-thrust behaviour
\end{itemize} & NR467 Sec.15, [5.4]; Sec.14 \\
\hline

6 & Rotor Casing \& Main Bevel Gear & \fa{} &
\begin{itemize}[leftmargin=*,nosep]
\item Rotor casing structure and thrust path
\item Input bevel gear (horizontal drive $\rightarrow$ vertical rotor)
\item Gear tooth calculations (ISO 6336)
\item Main thrust bearing design and life
\item Lubrication and oil monitoring
\end{itemize} & NR467 Sec.15, [5.5] \\
\hline

7 & Bottom Cover / Guard Plate Structure & \fa{} &
\begin{itemize}[leftmargin=*,nosep]
\item Guard plate below blade circle
\item Structural strength (grounding/impact)
\item Attachment to unit and hull
\item Detachment / fail-safe considerations
\end{itemize} & NR467 Sec.15, [5.6]; Part B, Ch.5 \\
\hline

8 & Hull Seat, Foundation \& Well Structure & \fa{} &
\begin{itemize}[leftmargin=*,nosep]
\item Foundation ring and bolting
\item Local hull reinforcement / FEA
\item Watertight integrity of installation well
\item Load transfer (thrust, steering moment)
\end{itemize} & NR467 Sec.15, [5.6]; Part B, Ch.5 \\
\hline

9 & Seal System (Rotor \& Blade Shafts) & \fa{} &
\begin{itemize}[leftmargin=*,nosep]
\item Seal type, arrangement and material
\item Leakage monitoring / oil-in-water prevention
\item EAL fluid if CLEANSHIP (see Ch.\ref{ch:notations})
\end{itemize} & NR467 Sec.15, [5.7] \\
\hline

\end{longtable}

\section{VSP Drive, Servo Control \& Steering Function}

\begin{longtable}{|p{0.5cm}|p{5cm}|p{1.5cm}|p{5.5cm}|p{2cm}|}
\caption{Voith-Schneider Propeller --- Drive \& Control} \label{tab:vsp-ctrl} \\
\hline
\rowcolor{bvblue}
\textcolor{white}{\textbf{\#}} & \textcolor{white}{\textbf{Drawing/Document}} & \textcolor{white}{\textbf{Category}} & \textcolor{white}{\textbf{Content Requirements}} & \textcolor{white}{\textbf{Reference}} \\
\hline
\endfirsthead
\multicolumn{5}{c}%
{{\tablename\ \thetable{} -- continued from previous page}} \\
\hline
\rowcolor{bvblue}
\textcolor{white}{\textbf{\#}} & \textcolor{white}{\textbf{Drawing/Document}} & \textcolor{white}{\textbf{Category}} & \textcolor{white}{\textbf{Content Requirements}} & \textcolor{white}{\textbf{Reference}} \\
\hline
\endhead
\hline \multicolumn{5}{|r|}{{Continued on next page}} \\ \hline
\endfoot
\hline
\endlastfoot

10 & Prime Mover Interface & \fa{} &
\begin{itemize}[leftmargin=*,nosep]
\item Diesel engine or electric motor input (see Ch.\ref{ch:primemover})
\item Input coupling / clutch / cardan shaft
\item Torsional vibration calculation of the complete drive train
\end{itemize} & NR467 Sec.15; Sec.7 \\
\hline

11 & Servo / Hydraulic Control Unit & \fa{} &
\begin{itemize}[leftmargin=*,nosep]
\item Control (steering) hydraulic schematic
\item Two independent power/servo systems where used for steering (redundancy)
\item Pressure ratings, relief, accumulators
\item See Chapter \ref{ch:hydraulics}
\end{itemize} & NR467 Sec.14; Ch.\ref{ch:hydraulics} \\
\hline

12 & Steering Control System \& Bridge Console & \fa{} &
\begin{itemize}[leftmargin=*,nosep]
\item Control transmitter (tiller/joystick) to control centre
\item Thrust magnitude \& direction indication
\item Follow-up / non-follow-up control
\item Alarms, emergency stop, mode changeover
\end{itemize} & NR467 Sec.14; Part C, Ch.4 \\
\hline

13 & Steering Function Compliance Package & \fa{} &
\begin{itemize}[leftmargin=*,nosep]
\item FMEA of steering function
\item Single-failure / duplication analysis
\item Response time, hard-over to hard-over
\item SOLAS II-1 Reg.29 demonstration (see Ch.\ref{ch:steering})
\end{itemize} & NR467 Sec.14; SOLAS II-1/29 \\
\hline

\end{longtable}

\noindent\textbf{Surveyor note (VSP):} A VSP is simultaneously a propeller, a
gear unit \emph{and} a steering gear. Do not approve the mechanical package in
isolation --- confirm that the steering-function documentation of
Chapter \ref{ch:steering} (redundancy, control stations, alarms, response time,
SOLAS compliance) has also been submitted. The blade guard/bottom plate and its
hull attachment are frequently under-documented; require structural
substantiation against grounding loads.

\newpage

% Chapter: Other Thruster Types
\chapter{Other Thruster Types (Rim-Driven, Pump-Jet, Gill/Jet, Steering Grid)}
\label{ch:othertypes}

This chapter addresses less-common propulsor arrangements recognised by BV. Each
follows the \textbf{general submissions of Chapter \ref{ch:general}} plus the
requirements of the \emph{nearest analogous conventional type}, with the
additional type-specific items below. Confirm the intended notation and service
(propulsion, take-me-home, DP, steering) before fixing the package.

\section{Rim-Driven Thrusters (Integrated Motor Propulsors)}

Rim-driven / hubless thrusters integrate a permanent-magnet (or induction) motor
whose rotor forms the propeller-blade tip ring inside the duct. There is no
mechanical shaft line or reduction gear; excitation and speed are set by a
frequency converter. Assessment shifts from gearing toward
\textbf{electrical machine, converter and duct-bearing} substantiation.

\begin{longtable}{|p{0.5cm}|p{5cm}|p{1.5cm}|p{5.5cm}|p{2cm}|}
\caption{Rim-Driven Thruster --- Type-Specific Requirements} \label{tab:rdt} \\
\hline
\rowcolor{bvblue}
\textcolor{white}{\textbf{\#}} & \textcolor{white}{\textbf{Drawing/Document}} & \textcolor{white}{\textbf{Category}} & \textcolor{white}{\textbf{Content Requirements}} & \textcolor{white}{\textbf{Reference}} \\
\hline
\endfirsthead
\multicolumn{5}{c}%
{{\tablename\ \thetable{} -- continued from previous page}} \\
\hline
\rowcolor{bvblue}
\textcolor{white}{\textbf{\#}} & \textcolor{white}{\textbf{Drawing/Document}} & \textcolor{white}{\textbf{Category}} & \textcolor{white}{\textbf{Content Requirements}} & \textcolor{white}{\textbf{Reference}} \\
\hline
\endhead
\hline \multicolumn{5}{|r|}{{Continued on next page}} \\ \hline
\endfoot
\hline
\endlastfoot

1 & Integrated Motor \& Duct Assembly & \fa{} &
\begin{itemize}[leftmargin=*,nosep]
\item Motor type (PM / induction), pole/winding data
\item Rotor rim / blade-tip ring construction
\item Duct structure and stator support
\item Cooling (seawater / freshwater path)
\end{itemize} & NR467 Sec.15; Pt C, Ch.2; IEC 60092-301 \\
\hline

2 & Frequency Converter / Drive Package & \fa{} &
\begin{itemize}[leftmargin=*,nosep]
\item Converter rating and topology
\item Harmonic / power-quality analysis
\item Protection, cooling, EMC
\item See Chapter \ref{ch:primemover} (electric drive)
\end{itemize} & NR467 Pt C, Ch.2; IEC 60092-101 \\
\hline

3 & Bearing \& Seal Arrangement & \fa{} &
\begin{itemize}[leftmargin=*,nosep]
\item Water-lubricated / magnetic bearing design
\item Bearing load and wear-life basis
\item Gap / air-gap control, seal design
\end{itemize} & NR467 Sec.15, [6.1] \\
\hline

4 & Blade \& Rim Strength Calculations & \fa{} &
\begin{itemize}[leftmargin=*,nosep]
\item Blade + rim stress (hydro + magnetic + centrifugal)
\item Fatigue and corrosion allowance
\item Cavitation assessment
\end{itemize} & NR467 Sec.15, [6.1]; Sec.7 \\
\hline

\end{longtable}

\section{Pump-Jet \& Ducted Rotor-Stator Propulsors}

Pump-jet units (ducted rotor with stator vanes, often for low-signature or
shallow-draught craft) are treated as a ducted propeller + shaft line. Submit the
\textbf{duct/stator structure and strength, rotor blade design and strength,
shaft line (Chapter \ref{ch:shafting}), bearings, seals} and the relevant drive
package. Where the unit also steers (deflector/vectoring), add the steering
function package of Chapter \ref{ch:steering}.

\section{Gill / Hydro-Jet Transverse Thrusters (Pump-Type)}

Pump-type transverse thrusters draw water through a hull inlet and discharge
through athwartship outlets, providing side thrust without an open tunnel. Submit
as for a tunnel thruster (Chapter \ref{ch:transverse}) but replace the tunnel
structure items with: \textbf{inlet grating and duct, internal pump/impeller
design and strength, discharge outlet and directional-valve arrangement, hull
penetration and watertight integrity}, plus the pump drive package.

\section{Steering Grid / Steering-Nozzle \& Rudder-Propeller Variants}

Fixed-nozzle (Kort-type) units with a steering grid or a steerable nozzle behave
as combined propulsion + steering devices. Submit the \textbf{nozzle/grid
structure and strength, propeller in nozzle, shaft line, steering actuator and
control}, and the full steering-gear compliance package of
Chapter \ref{ch:steering} (SOLAS II-1 Reg.29 where the device is the ship's
means of steering).

\noindent\textbf{Surveyor note (other types):} For any novel or hybrid
propulsor, apply the \emph{function-based} approach --- decompose the unit into
propeller/impeller, structure, shaft line, gear (if any), drive
(electric/diesel/hydraulic), bearings/seals and control --- and require the
corresponding chapter's submissions for each function. Novel technology may
additionally require an Approval-in-Principle (AiP) or risk-based/alternative-design
assessment.

\newpage

% Chapter: Prime Mover Considerations
\chapter{Prime Mover Considerations --- Diesel-Engine vs.\ Electric-Motor Driven Thrusters}
\label{ch:primemover}

\section{Purpose}

The preceding type chapters split the drive into \emph{electric} or
\emph{hydraulic}. This chapter consolidates the \textbf{additional package
elements driven by the choice of prime mover} so the surveyor can confirm the
correct drive-train documentation is present. A thruster may be driven by:
\begin{itemize}[leftmargin=*]
\item an \textbf{electric motor} (constant- or variable-speed, converter-fed), typical of diesel-electric ships, DP vessels, cruise ships and pods; or
\item a \textbf{diesel engine} (direct/geared mechanical drive), typical of tugs, OSVs and workboats with inboard engines driving Z-drive azimuth or tunnel thrusters via shafting, clutches and gearboxes; or
\item a \textbf{hydraulic motor} (covered per type chapter and Chapter \ref{ch:hydraulics}), common on retractable and smaller units.
\end{itemize}

\section{Electric-Motor-Driven Thrusters --- Additional Submissions}

\begin{longtable}{|p{0.5cm}|p{5cm}|p{1.5cm}|p{5.5cm}|p{2cm}|}
\caption{Electric-Motor Prime Mover --- Additional Requirements} \label{tab:pm-electric} \\
\hline
\rowcolor{bvblue}
\textcolor{white}{\textbf{\#}} & \textcolor{white}{\textbf{Drawing/Document}} & \textcolor{white}{\textbf{Category}} & \textcolor{white}{\textbf{Content Requirements}} & \textcolor{white}{\textbf{Reference}} \\
\hline
\endfirsthead
\multicolumn{5}{c}%
{{\tablename\ \thetable{} -- continued from previous page}} \\
\hline
\rowcolor{bvblue}
\textcolor{white}{\textbf{\#}} & \textcolor{white}{\textbf{Drawing/Document}} & \textcolor{white}{\textbf{Category}} & \textcolor{white}{\textbf{Content Requirements}} & \textcolor{white}{\textbf{Reference}} \\
\hline
\endhead
\hline \multicolumn{5}{|r|}{{Continued on next page}} \\ \hline
\endfoot
\hline
\endlastfoot

1 & Propulsion/Drive Motor Data Sheet & \fa{} &
\begin{itemize}[leftmargin=*,nosep]
\item Type (induction / PM / synchronous)
\item Rated \& overload power, torque-speed curve
\item Voltage, frequency, insulation \& IP class
\item Cooling method, temperature rise (IEC 60034)
\end{itemize} & NR467 Pt C, Ch.2; IEC 60092-301 \\
\hline

2 & Frequency Converter (VFD) Specification \& Arrangement & \fa{} &
\begin{itemize}[leftmargin=*,nosep]
\item Converter topology and rating
\item Output torque ripple at thruster
\item Protection, redundancy, bypass
\item Cooling and heat load to space
\end{itemize} & NR467 Pt C, Ch.2 \\
\hline

3 & Harmonic \& Power-Quality Analysis & \fa{} &
\begin{itemize}[leftmargin=*,nosep]
\item THD at main switchboard (limits per rules)
\item Effect on generators/other consumers
\item Mitigation (multi-pulse, active front end, filters)
\end{itemize} & NR467 Pt C, Ch.2 \\
\hline

4 & Supply Transformer \& Switchboard Interface & \fa{} &
\begin{itemize}[leftmargin=*,nosep]
\item Transformer rating, vector group, cooling
\item Feeder protection and discrimination
\item Short-circuit \& earth-fault coordination
\end{itemize} & NR467 Pt C, Ch.2; IEC 60092-201 \\
\hline

5 & Power Cable Sizing \& Routing & \fa{} &
\begin{itemize}[leftmargin=*,nosep]
\item Current rating, voltage drop, derating
\item Routing, support, bulkhead penetrations
\item Segregation (DP redundancy where relevant)
\end{itemize} & IEC 60092-352 / -352 \\
\hline

6 & Motor Protection, Monitoring \& Cooling & \fa{} &
\begin{itemize}[leftmargin=*,nosep]
\item Winding/bearing temperature monitoring
\item Overload, earth-fault, differential protection
\item Space heaters, forced cooling, condensation control
\end{itemize} & NR467 Pt C, Ch.2 / Ch.3 \\
\hline

7 & Torque-Ripple / Torsional Interaction Note & \fa{} &
\begin{itemize}[leftmargin=*,nosep]
\item Converter harmonics vs.\ shaft-line torsionals
\item Barred-speed / avoidance ranges if any
\item Cross-reference shaft TVC (Chapter \ref{ch:shafting})
\end{itemize} & NR467 Sec.7 \\
\hline

8 & Integration with PMS / Blackout Recovery & \fic{} &
\begin{itemize}[leftmargin=*,nosep]
\item Load-dependent start/stop, load shedding
\item Blackout and thruster restart sequence
\end{itemize} & NR467 Pt C, Ch.3 \\
\hline

\end{longtable}

\section{Diesel-Engine-Driven Thrusters --- Additional Submissions}

Where a thruster (typically a Z-drive azimuth on a tug/OSV, or a directly driven
tunnel/CP unit) is driven by a diesel engine, the \textbf{engine and its
auxiliary systems become part of the approvable package} under
\textbf{NR467 Pt C, Ch 1, Sec 2 (Diesel Engines)} in addition to the thruster
under Sec 15. A \textbf{torsional vibration calculation (TVC) of the complete
engine--clutch--shafting--thruster system is mandatory.}

\begin{longtable}{|p{0.5cm}|p{5cm}|p{1.5cm}|p{5.5cm}|p{2cm}|}
\caption{Diesel-Engine Prime Mover --- Additional Requirements} \label{tab:pm-diesel} \\
\hline
\rowcolor{bvblue}
\textcolor{white}{\textbf{\#}} & \textcolor{white}{\textbf{Drawing/Document}} & \textcolor{white}{\textbf{Category}} & \textcolor{white}{\textbf{Content Requirements}} & \textcolor{white}{\textbf{Reference}} \\
\hline
\endfirsthead
\multicolumn{5}{c}%
{{\tablename\ \thetable{} -- continued from previous page}} \\
\hline
\rowcolor{bvblue}
\textcolor{white}{\textbf{\#}} & \textcolor{white}{\textbf{Drawing/Document}} & \textcolor{white}{\textbf{Category}} & \textcolor{white}{\textbf{Content Requirements}} & \textcolor{white}{\textbf{Reference}} \\
\hline
\endhead
\hline \multicolumn{5}{|r|}{{Continued on next page}} \\ \hline
\endfoot
\hline
\endlastfoot

1 & Diesel Engine Approval / Type Data & \fa{} &
\begin{itemize}[leftmargin=*,nosep]
\item Engine general arrangement \& rating (MCR)
\item Type Approval / unit certification status
\item Crankshaft \& main running-gear data
\item Safety devices (overspeed, LO low-pressure shutdown)
\end{itemize} & NR467 Pt C, Ch.1, Sec.2 \\
\hline

2 & Torsional Vibration Calculation (TVC) & \fa{} &
\begin{itemize}[leftmargin=*,nosep]
\item Complete engine--coupling--gear--thruster model
\item Barred/critical speed ranges
\item Stresses in shafts, coupling, gear at criticals
\item Clutch-in/out and misfiring conditions
\end{itemize} & NR467 Sec.2; Sec.7 \\
\hline

3 & Clutch \& Flexible Coupling & \fa{} &
\begin{itemize}[leftmargin=*,nosep]
\item Clutch type, torque \& engagement control
\item Flexible coupling stiffness/damping data
\item Alignment and vibration decoupling
\end{itemize} & NR467 Sec.2; Sec.7 \\
\hline

4 & Engine Foundation \& Chocking & \fa{} &
\begin{itemize}[leftmargin=*,nosep]
\item Seating structure and holding-down bolts
\item Resilient mounts / chocking arrangement
\item Alignment provisions and thrust transmission
\end{itemize} & NR467 Sec.2; Part B, Ch.5 \\
\hline

5 & Fuel Oil System & \fa{} &
\begin{itemize}[leftmargin=*,nosep]
\item FO service/settling tanks, transfer, filtration
\item Piping, valves, quick-closing valves
\item Leakage containment, high-pressure line shielding
\end{itemize} & NR467 Pt C, Ch.1 (Piping) \\
\hline

6 & Lubricating Oil System & \fa{} &
\begin{itemize}[leftmargin=*,nosep]
\item LO pumps, cooler, filters, sump
\item Pressure/temperature monitoring \& alarms
\end{itemize} & NR467 Pt C, Ch.1 \\
\hline

7 & Cooling Water System & \fa{} &
\begin{itemize}[leftmargin=*,nosep]
\item Jacket \& charge-air cooling circuits
\item Sea-water / central cooling, pumps, expansion
\end{itemize} & NR467 Pt C, Ch.1 \\
\hline

8 & Starting \& Control Air / Starting System & \fa{} &
\begin{itemize}[leftmargin=*,nosep]
\item Air receivers, compressors, capacity for starts
\item Or electric-start battery capacity
\end{itemize} & NR467 Pt C, Ch.1 \\
\hline

9 & Exhaust Gas System & \fa{} &
\begin{itemize}[leftmargin=*,nosep]
\item Piping, expansion, lagging, silencer
\item Back-pressure, boundary insulation (fire)
\item SCR/EGR arrangement if fitted
\end{itemize} & NR467 Pt C, Ch.1; SOLAS II-2 \\
\hline

10 & Engine Control, Safety \& Monitoring & \fa{} &
\begin{itemize}[leftmargin=*,nosep]
\item Governor / speed control, overspeed trip
\item Alarms \& safeguards, local/remote control
\item Integration with thruster \& bridge controls
\end{itemize} & NR467 Pt C, Ch.1 / Ch.3 \\
\hline

11 & Emissions / NOx Documentation & \fic{} &
\begin{itemize}[leftmargin=*,nosep]
\item EIAPP certificate \& NOx Technical File
\item MARPOL Annex VI Tier compliance
\end{itemize} & MARPOL Annex VI \\
\hline

12 & Engine-Room Ventilation \& Air Supply & \fic{} &
\begin{itemize}[leftmargin=*,nosep]
\item Combustion air and space ventilation
\item Heat balance / temperature rise
\end{itemize} & NR467 Pt C, Ch.1 \\
\hline

\end{longtable}

\noindent\textbf{Surveyor note (prime mover):} The single most common gap for
diesel-driven Z-drive tugs is a \textbf{missing or incomplete TVC for the whole
train} (engine + clutch + cardan/shaft + azimuth gears + propeller). For
electric drives, the recurring gaps are the \textbf{harmonic/power-quality study}
and \textbf{cable-sizing/segregation}. Hybrid PTI/PTO arrangements must document
\emph{every} driving path and the mode-changeover logic. Cross-reference the
shaft-line torsional/alignment package in Chapter \ref{ch:shafting}.

\newpage

% Chapter 5: Steering Gear
\chapter{Thrusters as Part of Steering Gear}
\label{ch:steering}

\section{Applicability}

\textbf{\textcolor{warningred}{CRITICAL:}} When thrusters (azimuth, podded, or water-jets) are used as the PRIMARY or AUXILIARY steering system, they MUST comply with ALL steering gear requirements in addition to thruster requirements.

This applies to:
\begin{itemize}[leftmargin=*]
\item Azimuth thrusters used for steering
\item Podded propulsion units providing steering
\item Water-jet systems providing steering
\item Any thruster configuration that replaces or supplements rudder steering
\end{itemize}

\section{Additional Steering Gear Requirements}

\textbf{All applicable thruster drawings from Chapters \ref{ch:transverse}-\ref{ch:waterjets} apply, PLUS the following steering-specific submissions:}

\subsection{Steering System Design}

\begin{longtable}{|p{0.5cm}|p{5cm}|p{1.5cm}|p{5.5cm}|p{2cm}|}
\caption{Steering Gear - System Design} \label{tab:steering-design} \\
\hline
\rowcolor{bvblue}
\textcolor{white}{\textbf{\#}} & \textcolor{white}{\textbf{Drawing/Document}} & \textcolor{white}{\textbf{Category}} & \textcolor{white}{\textbf{Content Requirements}} & \textcolor{white}{\textbf{Reference}} \\
\hline
\endfirsthead

\multicolumn{5}{c}%
{{\tablename\ \thetable{} -- continued from previous page}} \\
\hline
\rowcolor{bvblue}
\textcolor{white}{\textbf{\#}} & \textcolor{white}{\textbf{Drawing/Document}} & \textcolor{white}{\textbf{Category}} & \textcolor{white}{\textbf{Content Requirements}} & \textcolor{white}{\textbf{Reference}} \\
\hline
\endhead

\hline \multicolumn{5}{|r|}{{Continued on next page}} \\ \hline
\endfoot

\hline
\endlastfoot

1 & Steering System General Arrangement & \fa{} & 
\begin{itemize}[leftmargin=*,nosep]
\item Complete steering system layout
\item Main and auxiliary steering provisions
\item Control stations
\item Power supply arrangements
\item Emergency steering provisions
\end{itemize} & NR467 Sec.14, [1] \\
\hline

2 & Steering Gear Capacity Calculations & \fa{} & 
\begin{itemize}[leftmargin=*,nosep]
\item Maximum steering torque
\item Time to move from 35° port to 35° starboard
\item Time to move from 35° to 30° on opposite side
\item Compliance with SOLAS requirements
\item Redundancy analysis
\end{itemize} & NR467 Sec.14, [2]; SOLAS Ch.II-1, Reg.29 \\
\hline

3 & Steering Control System & \fa{} & 
\begin{itemize}[leftmargin=*,nosep]
\item Control system architecture
\item Main control station (bridge)
\item Auxiliary control provisions
\item Emergency control
\item Autopilot integration
\item Manual override provisions
\end{itemize} & NR467 Sec.14, [3] \\
\hline

\end{longtable}

\subsection{Redundancy \& Reliability}

\begin{longtable}{|p{0.5cm}|p{5cm}|p{1.5cm}|p{5.5cm}|p{2cm}|}
\caption{Steering Gear - Redundancy Requirements} \label{tab:steering-redundancy} \\
\hline
\rowcolor{bvblue}
\textcolor{white}{\textbf{\#}} & \textcolor{white}{\textbf{Drawing/Document}} & \textcolor{white}{\textbf{Category}} & \textcolor{white}{\textbf{Content Requirements}} & \textcolor{white}{\textbf{Reference}} \\
\hline
\endfirsthead

\multicolumn{5}{c}%
{{\tablename\ \thetable{} -- continued from previous page}} \\
\hline
\rowcolor{bvblue}
\textcolor{white}{\textbf{\#}} & \textcolor{white}{\textbf{Drawing/Document}} & \textcolor{white}{\textbf{Category}} & \textcolor{white}{\textbf{Content Requirements}} & \textcolor{white}{\textbf{Reference}} \\
\hline
\endhead

\hline \multicolumn{5}{|r|}{{Continued on next page}} \\ \hline
\endfoot

\hline
\endlastfoot

4 & Main and Auxiliary Steering Systems & \fa{} & 
\begin{itemize}[leftmargin=*,nosep]
\item Main steering arrangement
\item Auxiliary steering arrangement
\item Independence of systems
\item Changeover provisions
\item Simultaneous operation capability
\end{itemize} & NR467 Sec.14, [4]; SOLAS Ch.II-1, Reg.29 \\
\hline

5 & Failure Mode and Effects Analysis (FMEA) & \fa{} & 
\begin{itemize}[leftmargin=*,nosep]
\item All potential failure modes
\item Effects on steering capability
\item Single failure tolerance
\item Redundancy verification
\item Mitigation measures
\end{itemize} & NR467 Sec.14, [5] \\
\hline

6 & Reliability Analysis & \fa{} & 
\begin{itemize}[leftmargin=*,nosep]
\item System reliability calculations
\item Component failure rates
\item Mean time between failures (MTBF)
\item Availability analysis
\item Compliance with steering gear standards
\end{itemize} & NR467 Sec.14, [5] \\
\hline

\end{longtable}

\subsection{Power Supply}

\begin{longtable}{|p{0.5cm}|p{5cm}|p{1.5cm}|p{5.5cm}|p{2cm}|}
\caption{Steering Gear - Power Supply} \label{tab:steering-power} \\
\hline
\rowcolor{bvblue}
\textcolor{white}{\textbf{\#}} & \textcolor{white}{\textbf{Drawing/Document}} & \textcolor{white}{\textbf{Category}} & \textcolor{white}{\textbf{Content Requirements}} & \textcolor{white}{\textbf{Reference}} \\
\hline
\endfirsthead

\multicolumn{5}{c}%
{{\tablename\ \thetable{} -- continued from previous page}} \\
\hline
\rowcolor{bvblue}
\textcolor{white}{\textbf{\#}} & \textcolor{white}{\textbf{Drawing/Document}} & \textcolor{white}{\textbf{Category}} & \textcolor{white}{\textbf{Content Requirements}} & \textcolor{white}{\textbf{Reference}} \\
\hline
\endhead

\hline \multicolumn{5}{|r|}{{Continued on next page}} \\ \hline
\endfoot

\hline
\endlastfoot

7 & Steering Gear Power Supply & \fa{} & 
\begin{itemize}[leftmargin=*,nosep]
\item Main power supply arrangement
\item Emergency power supply
\item Independent power sources for main/auxiliary
\item Switchboard arrangements
\item Circuit protection
\end{itemize} & NR467 Sec.14, [6]; SOLAS Ch.II-1, Reg.29 \\
\hline

8 & Power Supply Calculations & \fa{} & 
\begin{itemize}[leftmargin=*,nosep]
\item Load analysis
\item Cable sizing
\item Protection device coordination
\item Voltage drop calculations
\item Short circuit analysis
\end{itemize} & NR467 Sec.14, [6]; IEC 60092-352 \\
\hline

\end{longtable}

\subsection{Control Stations}

\begin{longtable}{|p{0.5cm}|p{5cm}|p{1.5cm}|p{5.5cm}|p{2cm}|}
\caption{Steering Gear - Control Stations} \label{tab:steering-control} \\
\hline
\rowcolor{bvblue}
\textcolor{white}{\textbf{\#}} & \textcolor{white}{\textbf{Drawing/Document}} & \textcolor{white}{\textbf{Category}} & \textcolor{white}{\textbf{Content Requirements}} & \textcolor{white}{\textbf{Reference}} \\
\hline
\endfirsthead

\multicolumn{5}{c}%
{{\tablename\ \thetable{} -- continued from previous page}} \\
\hline
\rowcolor{bvblue}
\textcolor{white}{\textbf{\#}} & \textcolor{white}{\textbf{Drawing/Document}} & \textcolor{white}{\textbf{Category}} & \textcolor{white}{\textbf{Content Requirements}} & \textcolor{white}{\textbf{Reference}} \\
\hline
\endhead

\hline \multicolumn{5}{|r|}{{Continued on next page}} \\ \hline
\endfoot

\hline
\endlastfoot

9 & Bridge Control Console & \fa{} & 
\begin{itemize}[leftmargin=*,nosep]
\item Main steering control layout
\item Steering wheel or joystick
\item Rudder angle indicator (or thruster angle)
\item Rate of turn indicator
\item Alarm and monitoring displays
\item Autopilot controls
\end{itemize} & NR467 Sec.14, [7] \\
\hline

10 & Emergency Steering Station & \fa{} & 
\begin{itemize}[leftmargin=*,nosep]
\item Emergency control location
\item Control equipment
\item Communication with bridge
\item Indication and monitoring
\item Access provisions
\end{itemize} & NR467 Sec.14, [8]; SOLAS Ch.II-1, Reg.29 \\
\hline

11 & Local Control Station & \fa{} & 
\begin{itemize}[leftmargin=*,nosep]
\item Local control panel (if applicable)
\item Control mode selection
\item Indication and monitoring
\item Interlock with other control stations
\end{itemize} & NR467 Sec.14, [9] \\
\hline

\end{longtable}

\subsection{Alarms and Monitoring}

\begin{longtable}{|p{0.5cm}|p{5cm}|p{1.5cm}|p{5.5cm}|p{2cm}|}
\caption{Steering Gear - Alarms \& Monitoring} \label{tab:steering-alarms} \\
\hline
\rowcolor{bvblue}
\textcolor{white}{\textbf{\#}} & \textcolor{white}{\textbf{Drawing/Document}} & \textcolor{white}{\textbf{Category}} & \textcolor{white}{\textbf{Content Requirements}} & \textcolor{white}{\textbf{Reference}} \\
\hline
\endfirsthead

\multicolumn{5}{c}%
{{\tablename\ \thetable{} -- continued from previous page}} \\
\hline
\rowcolor{bvblue}
\textcolor{white}{\textbf{\#}} & \textcolor{white}{\textbf{Drawing/Document}} & \textcolor{white}{\textbf{Category}} & \textcolor{white}{\textbf{Content Requirements}} & \textcolor{white}{\textbf{Reference}} \\
\hline
\endhead

\hline \multicolumn{5}{|r|}{{Continued on next page}} \\ \hline
\endfoot

\hline
\endlastfoot

12 & Alarm System Design & \fa{} & 
\begin{itemize}[leftmargin=*,nosep]
\item Power supply failure alarms
\item Control system failure alarms
\item Actuator failure alarms
\item Position deviation alarms
\item Alarm display on bridge
\end{itemize} & NR467 Sec.14, [10] \\
\hline

13 & Monitoring System & \fa{} & 
\begin{itemize}[leftmargin=*,nosep]
\item Position monitoring
\item Rate monitoring
\item Torque/thrust monitoring
\item Temperature monitoring
\item Data logging
\end{itemize} & NR467 Sec.14, [10] \\
\hline

\end{longtable}

\subsection{Hydraulic Systems (if applicable)}

\begin{longtable}{|p{0.5cm}|p{5cm}|p{1.5cm}|p{5.5cm}|p{2cm}|}
\caption{Steering Gear - Hydraulic System} \label{tab:steering-hydraulic} \\
\hline
\rowcolor{bvblue}
\textcolor{white}{\textbf{\#}} & \textcolor{white}{\textbf{Drawing/Document}} & \textcolor{white}{\textbf{Category}} & \textcolor{white}{\textbf{Content Requirements}} & \textcolor{white}{\textbf{Reference}} \\
\hline
\endfirsthead

\multicolumn{5}{c}%
{{\tablename\ \thetable{} -- continued from previous page}} \\
\hline
\rowcolor{bvblue}
\textcolor{white}{\textbf{\#}} & \textcolor{white}{\textbf{Drawing/Document}} & \textcolor{white}{\textbf{Category}} & \textcolor{white}{\textbf{Content Requirements}} & \textcolor{white}{\textbf{Reference}} \\
\hline
\endhead

\hline \multicolumn{5}{|r|}{{Continued on next page}} \\ \hline
\endfoot

\hline
\endlastfoot

14 & Steering Gear Hydraulic System & \fa{} & 
\begin{itemize}[leftmargin=*,nosep]
\item Main hydraulic system
\item Auxiliary hydraulic system
\item Independence of systems
\item Accumulator provisions
\item See Chapter \ref{ch:hydraulics} for detailed requirements
\end{itemize} & NR467 Sec.14, [11]; Ch.\ref{ch:hydraulics} \\
\hline

\end{longtable}

\subsection{Testing \& Commissioning}

\begin{longtable}{|p{0.5cm}|p{5cm}|p{1.5cm}|p{5.5cm}|p{2cm}|}
\caption{Steering Gear - Testing Requirements} \label{tab:steering-testing} \\
\hline
\rowcolor{bvblue}
\textcolor{white}{\textbf{\#}} & \textcolor{white}{\textbf{Drawing/Document}} & \textcolor{white}{\textbf{Category}} & \textcolor{white}{\textbf{Content Requirements}} & \textcolor{white}{\textbf{Reference}} \\
\hline
\endfirsthead

\multicolumn{5}{c}%
{{\tablename\ \thetable{} -- continued from previous page}} \\
\hline
\rowcolor{bvblue}
\textcolor{white}{\textbf{\#}} & \textcolor{white}{\textbf{Drawing/Document}} & \textcolor{white}{\textbf{Category}} & \textcolor{white}{\textbf{Content Requirements}} & \textcolor{white}{\textbf{Reference}} \\
\hline
\endhead

\hline \multicolumn{5}{|r|}{{Continued on next page}} \\ \hline
\endfoot

\hline
\endlastfoot

15 & Steering Gear Testing Procedures & \fa{} & 
\begin{itemize}[leftmargin=*,nosep]
\item Factory acceptance tests
\item Installation tests
\item Sea trial procedures
\item Performance verification
\item Redundancy testing
\item Emergency steering tests
\end{itemize} & NR467 Sec.14, [12]; SOLAS Ch.II-1, Reg.29 \\
\hline

16 & Steering Gear Test Records & \fic{} & 
\begin{itemize}[leftmargin=*,nosep]
\item FAT reports
\item Installation test reports
\item Sea trial results
\item Performance verification
\item As-built documentation
\end{itemize} & NR467 Sec.14, [12] \\
\hline

\end{longtable}

\subsection{Operating Manuals}

\begin{longtable}{|p{0.5cm}|p{5cm}|p{1.5cm}|p{5.5cm}|p{2cm}|}
\caption{Steering Gear - Operating Manuals} \label{tab:steering-manuals} \\
\hline
\rowcolor{bvblue}
\textcolor{white}{\textbf{\#}} & \textcolor{white}{\textbf{Drawing/Document}} & \textcolor{white}{\textbf{Category}} & \textcolor{white}{\textbf{Content Requirements}} & \textcolor{white}{\textbf{Reference}} \\
\hline
\endfirsthead

\multicolumn{5}{c}%
{{\tablename\ \thetable{} -- continued from previous page}} \\
\hline
\rowcolor{bvblue}
\textcolor{white}{\textbf{\#}} & \textcolor{white}{\textbf{Drawing/Document}} & \textcolor{white}{\textbf{Category}} & \textcolor{white}{\textbf{Content Requirements}} & \textcolor{white}{\textbf{Reference}} \\
\hline
\endhead

\hline \multicolumn{5}{|r|}{{Continued on next page}} \\ \hline
\endfoot

\hline
\endlastfoot

17 & Steering Gear Operating Manual & \fic{} & 
\begin{itemize}[leftmargin=*,nosep]
\item Normal operating procedures
\item Changeover procedures (main/auxiliary)
\item Emergency procedures
\item Troubleshooting guide
\item Safety precautions
\end{itemize} & NR467 Sec.14, [13] \\
\hline

18 & Maintenance Manual & \fic{} & 
\begin{itemize}[leftmargin=*,nosep]
\item Maintenance schedule
\item Inspection procedures
\item Lubrication requirements
\item Spare parts list
\item Overhaul procedures
\end{itemize} & NR467 Sec.14, [13] \\
\hline

\end{longtable}

\subsection{SOLAS Compliance Documentation}

\begin{longtable}{|p{0.5cm}|p{5cm}|p{1.5cm}|p{5.5cm}|p{2cm}|}
\caption{Steering Gear - SOLAS Compliance} \label{tab:steering-solas} \\
\hline
\rowcolor{bvblue}
\textcolor{white}{\textbf{\#}} & \textcolor{white}{\textbf{Drawing/Document}} & \textcolor{white}{\textbf{Category}} & \textcolor{white}{\textbf{Content Requirements}} & \textcolor{white}{\textbf{Reference}} \\
\hline
\endfirsthead

\multicolumn{5}{c}%
{{\tablename\ \thetable{} -- continued from previous page}} \\
\hline
\rowcolor{bvblue}
\textcolor{white}{\textbf{\#}} & \textcolor{white}{\textbf{Drawing/Document}} & \textcolor{white}{\textbf{Category}} & \textcolor{white}{\textbf{Content Requirements}} & \textcolor{white}{\textbf{Reference}} \\
\hline
\endhead

\hline \multicolumn{5}{|r|}{{Continued on next page}} \\ \hline
\endfoot

\hline
\endlastfoot

19 & SOLAS Compliance Statement & \fa{} & 
\begin{itemize}[leftmargin=*,nosep]
\item Demonstration of compliance with SOLAS Ch.II-1, Reg.29
\item Steering gear capacity verification
\item Redundancy verification
\item Emergency steering capability
\item Testing and commissioning compliance
\end{itemize} & NR467 Sec.14, [14]; SOLAS Ch.II-1, Reg.29 \\
\hline

\end{longtable}

\newpage

% Chapter 6: Main Shafting
\chapter{Thrusters as Part of Main Shafting}
\label{ch:shafting}

\section{Applicability}

When thruster shafts (propeller shafts, intermediate shafts, vertical shafts in azimuth thrusters, or water-jet drive shafts) are considered part of the main propulsion shafting system, they MUST comply with shafting requirements in NR467 Section 7.

This typically applies to:
\begin{itemize}[leftmargin=*]
\item Azimuth thruster vertical and horizontal shafts
\item Podded propulsion shafts
\item Water-jet drive shafts
\item Tunnel thruster propeller shafts (for high-power units)
\item Any shaft transmitting significant propulsion power
\end{itemize}

\section{Overlap with Thruster Requirements}

\textbf{NOTE:} Many shaft-related drawings and calculations are already covered in Chapters \ref{ch:transverse}-\ref{ch:waterjets}. This chapter identifies ADDITIONAL requirements or ENHANCED detail levels required when shafts are classified as "main shafting."

\subsection{Shaft Design Drawings}

\begin{longtable}{|p{0.5cm}|p{5cm}|p{1.5cm}|p{5.5cm}|p{2cm}|}
\caption{Main Shafting - Design Drawings} \label{tab:shafting-design} \\
\hline
\rowcolor{bvblue}
\textcolor{white}{\textbf{\#}} & \textcolor{white}{\textbf{Drawing/Document}} & \textcolor{white}{\textbf{Category}} & \textcolor{white}{\textbf{Content Requirements}} & \textcolor{white}{\textbf{Reference}} \\
\hline
\endfirsthead

\multicolumn{5}{c}%
{{\tablename\ \thetable{} -- continued from previous page}} \\
\hline
\rowcolor{bvblue}
\textcolor{white}{\textbf{\#}} & \textcolor{white}{\textbf{Drawing/Document}} & \textcolor{white}{\textbf{Category}} & \textcolor{white}{\textbf{Content Requirements}} & \textcolor{white}{\textbf{Reference}} \\
\hline
\endhead

\hline \multicolumn{5}{|r|}{{Continued on next page}} \\ \hline
\endfoot

\hline
\endlastfoot

1 & Shaft Line General Arrangement & \fa{} & 
\begin{itemize}[leftmargin=*,nosep]
\item Complete shaft line layout
\item All shaft sections identified
\item Bearing locations
\item Coupling locations
\item Seal locations
\item Coordinate system
\end{itemize} & NR467 Sec.7, [1] \\
\hline

2 & Shaft Detail Drawings & \fa{} & 
\begin{itemize}[leftmargin=*,nosep]
\item Dimensional drawings of each shaft section
\item Diameter changes and transitions
\item Keyway details
\item Thread details
\item Surface finish requirements
\item Material specifications
\item Heat treatment requirements
\end{itemize} & NR467 Sec.7, [2] \\
\hline

3 & Shaft End Connections & \fa{} & 
\begin{itemize}[leftmargin=*,nosep]
\item Flange designs
\item Coupling designs
\item Bolt specifications
\item Fit tolerances
\item Assembly procedures
\end{itemize} & NR467 Sec.7, [3] \\
\hline

\end{longtable}

\subsection{Shaft Strength Calculations}

\begin{longtable}{|p{0.5cm}|p{5cm}|p{1.5cm}|p{5.5cm}|p{2cm}|}
\caption{Main Shafting - Strength Calculations} \label{tab:shafting-strength} \\
\hline
\rowcolor{bvblue}
\textcolor{white}{\textbf{\#}} & \textcolor{white}{\textbf{Drawing/Document}} & \textcolor{white}{\textbf{Category}} & \textcolor{white}{\textbf{Content Requirements}} & \textcolor{white}{\textbf{Reference}} \\
\hline
\endfirsthead

\multicolumn{5}{c}%
{{\tablename\ \thetable{} -- continued from previous page}} \\
\hline
\rowcolor{bvblue}
\textcolor{white}{\textbf{\#}} & \textcolor{white}{\textbf{Drawing/Document}} & \textcolor{white}{\textbf{Category}} & \textcolor{white}{\textbf{Content Requirements}} & \textcolor{white}{\textbf{Reference}} \\
\hline
\endhead

\hline \multicolumn{5}{|r|}{{Continued on next page}} \\ \hline
\endfoot

\hline
\endlastfoot

4 & Shaft Diameter Calculations & \fa{} & 
\begin{itemize}[leftmargin=*,nosep]
\item Minimum diameter calculations per BV formula
\item Torsional strength verification
\item Bending strength verification
\item Combined stress analysis
\item Safety factors applied
\end{itemize} & NR467 Sec.7, [4] \\
\hline

5 & Shaft Stress Analysis & \fa{} & 
\begin{itemize}[leftmargin=*,nosep]
\item Detailed stress calculations at critical sections
\item Stress concentration factors
\item Keyway effects
\item Thread effects
\item FEA results (if applicable)
\end{itemize} & NR467 Sec.7, [4] \\
\hline

6 & Fatigue Analysis & \fa{} & 
\begin{itemize}[leftmargin=*,nosep]
\item Fatigue life calculations
\item S-N curves for materials
\item Stress range analysis
\item Cumulative damage assessment
\item Compliance with fatigue criteria
\end{itemize} & NR467 Sec.7, [5] \\
\hline

7 & Deflection Calculations & \fa{} & 
\begin{itemize}[leftmargin=*,nosep]
\item Shaft deflection under load
\item Bearing alignment effects
\item Seal clearance verification
\item Acceptable deflection limits
\end{itemize} & NR467 Sec.7, [6] \\
\hline

\end{longtable}

\subsection{Critical Speed Analysis}

\begin{longtable}{|p{0.5cm}|p{5cm}|p{1.5cm}|p{5.5cm}|p{2cm}|}
\caption{Main Shafting - Critical Speed Analysis} \label{tab:shafting-critical} \\
\hline
\rowcolor{bvblue}
\textcolor{white}{\textbf{\#}} & \textcolor{white}{\textbf{Drawing/Document}} & \textcolor{white}{\textbf{Category}} & \textcolor{white}{\textbf{Content Requirements}} & \textcolor{white}{\textbf{Reference}} \\
\hline
\endfirsthead

\multicolumn{5}{c}%
{{\tablename\ \thetable{} -- continued from previous page}} \\
\hline
\rowcolor{bvblue}
\textcolor{white}{\textbf{\#}} & \textcolor{white}{\textbf{Drawing/Document}} & \textcolor{white}{\textbf{Category}} & \textcolor{white}{\textbf{Content Requirements}} & \textcolor{white}{\textbf{Reference}} \\
\hline
\endhead

\hline \multicolumn{5}{|r|}{{Continued on next page}} \\ \hline
\endfoot

\hline
\endlastfoot

8 & Torsional Vibration Analysis & \fa{} & 
\begin{itemize}[leftmargin=*,nosep]
\item Natural frequencies calculation
\item Forced vibration analysis
\item Resonance identification
\item Barred speed ranges
\item Damping analysis
\item Compliance with acceptable vibration levels
\end{itemize} & NR467 Sec.7, [7] \\
\hline

9 & Lateral Vibration Analysis & \fa{} & 
\begin{itemize}[leftmargin=*,nosep]
\item Critical speed calculations
\item Campbell diagram
\item Operating speed verification
\item Margin to critical speeds
\item Bearing stiffness effects
\end{itemize} & NR467 Sec.7, [7] \\
\hline

10 & Whirling Analysis & \fa{} & 
\begin{itemize}[leftmargin=*,nosep]
\item Whirling speed calculations
\item Forward and backward whirl
\item Operating range verification
\item Stability analysis
\end{itemize} & NR467 Sec.7, [7] \\
\hline

\end{longtable}

\subsection{Bearing Design}

\begin{longtable}{|p{0.5cm}|p{5cm}|p{1.5cm}|p{5.5cm}|p{2cm}|}
\caption{Main Shafting - Bearing Design} \label{tab:shafting-bearing} \\
\hline
\rowcolor{bvblue}
\textcolor{white}{\textbf{\#}} & \textcolor{white}{\textbf{Drawing/Document}} & \textcolor{white}{\textbf{Category}} & \textcolor{white}{\textbf{Content Requirements}} & \textcolor{white}{\textbf{Reference}} \\
\hline
\endfirsthead

\multicolumn{5}{c}%
{{\tablename\ \thetable{} -- continued from previous page}} \\
\hline
\rowcolor{bvblue}
\textcolor{white}{\textbf{\#}} & \textcolor{white}{\textbf{Drawing/Document}} & \textcolor{white}{\textbf{Category}} & \textcolor{white}{\textbf{Content Requirements}} & \textcolor{white}{\textbf{Reference}} \\
\hline
\endhead

\hline \multicolumn{5}{|r|}{{Continued on next page}} \\ \hline
\endfoot

\hline
\endlastfoot

11 & Bearing Arrangement Drawing & \fa{} & 
\begin{itemize}[leftmargin=*,nosep]
\item Bearing types and locations
\item Bearing housing design
\item Lubrication system
\item Seal arrangements
\item Cooling provisions
\item Mounting details
\end{itemize} & NR467 Sec.7, [8] \\
\hline

12 & Bearing Load Calculations & \fa{} & 
\begin{itemize}[leftmargin=*,nosep]
\item Radial loads
\item Axial loads (thrust bearings)
\item Moment loads
\item Load distribution among bearings
\item Dynamic load factors
\end{itemize} & NR467 Sec.7, [8] \\
\hline

13 & Bearing Life Calculations & \fa{} & 
\begin{itemize}[leftmargin=*,nosep]
\item L10 life calculations
\item Operating conditions
\item Load and speed effects
\item Lubrication effects
\item Minimum acceptable life
\end{itemize} & NR467 Sec.7, [8] \\
\hline

14 & Bearing Lubrication System & \fa{} & 
\begin{itemize}[leftmargin=*,nosep]
\item Lubrication method (oil bath, forced circulation, grease)
\item Oil type and properties
\item Oil flow rate (forced systems)
\item Cooling arrangements
\item Filtration system
\item Oil level and temperature monitoring
\end{itemize} & NR467 Sec.7, [9] \\
\hline

\end{longtable}

\subsection{Couplings}

\begin{longtable}{|p{0.5cm}|p{5cm}|p{1.5cm}|p{5.5cm}|p{2cm}|}
\caption{Main Shafting - Couplings} \label{tab:shafting-coupling} \\
\hline
\rowcolor{bvblue}
\textcolor{white}{\textbf{\#}} & \textcolor{white}{\textbf{Drawing/Document}} & \textcolor{white}{\textbf{Category}} & \textcolor{white}{\textbf{Content Requirements}} & \textcolor{white}{\textbf{Reference}} \\
\hline
\endfirsthead

\multicolumn{5}{c}%
{{\tablename\ \thetable{} -- continued from previous page}} \\
\hline
\rowcolor{bvblue}
\textcolor{white}{\textbf{\#}} & \textcolor{white}{\textbf{Drawing/Document}} & \textcolor{white}{\textbf{Category}} & \textcolor{white}{\textbf{Content Requirements}} & \textcolor{white}{\textbf{Reference}} \\
\hline
\endhead

\hline \multicolumn{5}{|r|}{{Continued on next page}} \\ \hline
\endfoot

\hline
\endlastfoot

15 & Coupling Design Drawing & \fa{} & 
\begin{itemize}[leftmargin=*,nosep]
\item Coupling type (rigid, flexible)
\item Dimensions and material
\item Bolt specifications
\item Fit tolerances
\item Assembly procedure
\end{itemize} & NR467 Sec.7, [10] \\
\hline

16 & Coupling Strength Calculations & \fa{} & 
\begin{itemize}[leftmargin=*,nosep]
\item Torque capacity
\item Bolt strength verification
\item Flange strength
\item Fit stress analysis
\item Fatigue assessment
\end{itemize} & NR467 Sec.7, [10] \\
\hline

17 & Flexible Coupling Characteristics & \fa{} & 
\begin{itemize}[leftmargin=*,nosep]
\item Torsional stiffness
\item Damping properties
\item Misalignment capacity
\item Effect on torsional vibration
\end{itemize} & NR467 Sec.7, [10] \\
\hline

\end{longtable}

\subsection{Seals}

\begin{longtable}{|p{0.5cm}|p{5cm}|p{1.5cm}|p{5.5cm}|p{2cm}|}
\caption{Main Shafting - Seals} \label{tab:shafting-seal} \\
\hline
\rowcolor{bvblue}
\textcolor{white}{\textbf{\#}} & \textcolor{white}{\textbf{Drawing/Document}} & \textcolor{white}{\textbf{Category}} & \textcolor{white}{\textbf{Content Requirements}} & \textcolor{white}{\textbf{Reference}} \\
\hline
\endfirsthead

\multicolumn{5}{c}%
{{\tablename\ \thetable{} -- continued from previous page}} \\
\hline
\rowcolor{bvblue}
\textcolor{white}{\textbf{\#}} & \textcolor{white}{\textbf{Drawing/Document}} & \textcolor{white}{\textbf{Category}} & \textcolor{white}{\textbf{Content Requirements}} & \textcolor{white}{\textbf{Reference}} \\
\hline
\endhead

\hline \multicolumn{5}{|r|}{{Continued on next page}} \\ \hline
\endfoot

\hline
\endlastfoot

18 & Seal System Design & \fa{} & 
\begin{itemize}[leftmargin=*,nosep]
\item Seal type and material
\item Seal arrangement
\item Pressure rating
\item Leakage monitoring provisions
\item Emergency sealing (if applicable)
\end{itemize} & NR467 Sec.7, [11] \\
\hline

19 & Seal Performance Data & \fa{} & 
\begin{itemize}[leftmargin=*,nosep]
\item Operating pressure and temperature
\item Shaft speed limitations
\item Expected leakage rates
\item Service life
\item Maintenance requirements
\end{itemize} & NR467 Sec.7, [11] \\
\hline

\end{longtable}

\subsection{Alignment}

\begin{longtable}{|p{0.5cm}|p{5cm}|p{1.5cm}|p{5.5cm}|p{2cm}|}
\caption{Main Shafting - Alignment} \label{tab:shafting-alignment} \\
\hline
\rowcolor{bvblue}
\textcolor{white}{\textbf{\#}} & \textcolor{white}{\textbf{Drawing/Document}} & \textcolor{white}{\textbf{Category}} & \textcolor{white}{\textbf{Content Requirements}} & \textcolor{white}{\textbf{Reference}} \\
\hline
\endfirsthead

\multicolumn{5}{c}%
{{\tablename\ \thetable{} -- continued from previous page}} \\
\hline
\rowcolor{bvblue}
\textcolor{white}{\textbf{\#}} & \textcolor{white}{\textbf{Drawing/Document}} & \textcolor{white}{\textbf{Category}} & \textcolor{white}{\textbf{Content Requirements}} & \textcolor{white}{\textbf{Reference}} \\
\hline
\endhead

\hline \multicolumn{5}{|r|}{{Continued on next page}} \\ \hline
\endfoot

\hline
\endlastfoot

20 & Shaft Alignment Calculations & \fa{} & 
\begin{itemize}[leftmargin=*,nosep]
\item Alignment procedure
\item Bearing reactions
\item Shaft sag calculations
\item Thermal growth considerations
\item Acceptable alignment tolerances
\end{itemize} & NR467 Sec.7, [12] \\
\hline

21 & Alignment Procedure & \fic{} & 
\begin{itemize}[leftmargin=*,nosep]
\item Step-by-step alignment procedure
\item Measurement methods
\item Adjustment provisions
\item Verification methods
\item Documentation requirements
\end{itemize} & NR467 Sec.7, [12] \\
\hline

\end{longtable}

\subsection{Material Specifications}

\begin{longtable}{|p{0.5cm}|p{5cm}|p{1.5cm}|p{5.5cm}|p{2cm}|}
\caption{Main Shafting - Material Specifications} \label{tab:shafting-material} \\
\hline
\rowcolor{bvblue}
\textcolor{white}{\textbf{\#}} & \textcolor{white}{\textbf{Drawing/Document}} & \textcolor{white}{\textbf{Category}} & \textcolor{white}{\textbf{Content Requirements}} & \textcolor{white}{\textbf{Reference}} \\
\hline
\endfirsthead

\multicolumn{5}{c}%
{{\tablename\ \thetable{} -- continued from previous page}} \\
\hline
\rowcolor{bvblue}
\textcolor{white}{\textbf{\#}} & \textcolor{white}{\textbf{Drawing/Document}} & \textcolor{white}{\textbf{Category}} & \textcolor{white}{\textbf{Content Requirements}} & \textcolor{white}{\textbf{Reference}} \\
\hline
\endhead

\hline \multicolumn{5}{|r|}{{Continued on next page}} \\ \hline
\endfoot

\hline
\endlastfoot

22 & Shaft Material Specifications & \fa{} & 
\begin{itemize}[leftmargin=*,nosep]
\item Material grades (steel, stainless steel, etc.)
\item Mechanical properties
\item Heat treatment requirements
\item Surface treatment (if applicable)
\item Corrosion protection
\item Compliance with BV approved materials
\end{itemize} & NR467 Sec.7, [13]; Part B, Ch.3 \\
\hline

23 & Material Certificates & \fa{} & 
\begin{itemize}[leftmargin=*,nosep]
\item Material test certificates (during construction)
\item Chemical composition
\item Mechanical properties
\item Heat treatment records
\item Traceability
\end{itemize} & NR467 Part B, Ch.3 \\
\hline

\end{longtable}

\subsection{Non-Destructive Testing (NDT)}

\begin{longtable}{|p{0.5cm}|p{5cm}|p{1.5cm}|p{5.5cm}|p{2cm}|}
\caption{Main Shafting - NDT Requirements} \label{tab:shafting-ndt} \\
\hline
\rowcolor{bvblue}
\textcolor{white}{\textbf{\#}} & \textcolor{white}{\textbf{Drawing/Document}} & \textcolor{white}{\textbf{Category}} & \textcolor{white}{\textbf{Content Requirements}} & \textcolor{white}{\textbf{Reference}} \\
\hline
\endfirsthead

\multicolumn{5}{c}%
{{\tablename\ \thetable{} -- continued from previous page}} \\
\hline
\rowcolor{bvblue}
\textcolor{white}{\textbf{\#}} & \textcolor{white}{\textbf{Drawing/Document}} & \textcolor{white}{\textbf{Category}} & \textcolor{white}{\textbf{Content Requirements}} & \textcolor{white}{\textbf{Reference}} \\
\hline
\endhead

\hline \multicolumn{5}{|r|}{{Continued on next page}} \\ \hline
\endfoot

\hline
\endlastfoot

24 & NDT Plan & \fa{} & 
\begin{itemize}[leftmargin=*,nosep]
\item NDT methods (UT, MT, PT, RT)
\item Extent of testing
\item Acceptance criteria
\item Inspector qualifications
\item Testing stages
\end{itemize} & NR467 Sec.7, [14]; Part B, Ch.4 \\
\hline

25 & NDT Reports & \fic{} & 
\begin{itemize}[leftmargin=*,nosep]
\item NDT test reports (during construction)
\item Test results
\item Defect identification and disposition
\item Inspector certification
\end{itemize} & NR467 Part B, Ch.4 \\
\hline

\end{longtable}

\subsection{Testing \& Commissioning}

\begin{longtable}{|p{0.5cm}|p{5cm}|p{1.5cm}|p{5.5cm}|p{2cm}|}
\caption{Main Shafting - Testing \& Commissioning} \label{tab:shafting-testing} \\
\hline
\rowcolor{bvblue}
\textcolor{white}{\textbf{\#}} & \textcolor{white}{\textbf{Drawing/Document}} & \textcolor{white}{\textbf{Category}} & \textcolor{white}{\textbf{Content Requirements}} & \textcolor{white}{\textbf{Reference}} \\
\hline
\endfirsthead

\multicolumn{5}{c}%
{{\tablename\ \thetable{} -- continued from previous page}} \\
\hline
\rowcolor{bvblue}
\textcolor{white}{\textbf{\#}} & \textcolor{white}{\textbf{Drawing/Document}} & \textcolor{white}{\textbf{Category}} & \textcolor{white}{\textbf{Content Requirements}} & \textcolor{white}{\textbf{Reference}} \\
\hline
\endhead

\hline \multicolumn{5}{|r|}{{Continued on next page}} \\ \hline
\endfoot

\hline
\endlastfoot

26 & Shaft Testing Procedures & \fa{} & 
\begin{itemize}[leftmargin=*,nosep]
\item Dimensional verification
\item Material testing
\item NDT
\item Balance testing (if applicable)
\item Assembly testing
\item Acceptance criteria
\end{itemize} & NR467 Sec.7, [15] \\
\hline

27 & Vibration Monitoring Plan & \fa{} & 
\begin{itemize}[leftmargin=*,nosep]
\item Vibration measurement locations
\item Measurement methods
\item Acceptance criteria
\item Monitoring during sea trials
\item Long-term monitoring provisions
\end{itemize} & NR467 Sec.7, [16] \\
\hline

28 & Test Reports & \fic{} &
\begin{itemize}[leftmargin=*,nosep]
\item Factory test reports
\item Installation test reports
\item Sea trial vibration measurements
\item Alignment verification
\item As-built documentation
\end{itemize} & NR467 Sec.7, [15] \\
\hline

\end{longtable}

\section{MON-SHAFT \& Shaft-Line Condition-Monitoring Notations}
\label{sec:monshaft}

Where the thruster forms part of the propulsion shaft line (azimuth, podded,
water-jet or high-power tunnel units with oil-lubricated stern-tube/lower-unit
bearings), the owner may request the \textbf{MON-SHAFT} additional class
notation. MON-SHAFT allows the classical periodic \emph{tailshaft withdrawal /
complete shaft survey} to be replaced by a survey regime based on
\textbf{continuous monitoring and condition assessment} of the oil-lubricated
shaft/bearing/seal arrangement. It therefore \emph{adds documentation} to the
shaft-line package and interacts directly with the thruster's lower-unit
bearing, seal and oil systems.

\begin{longtable}{|p{0.5cm}|p{5cm}|p{1.5cm}|p{5.5cm}|p{2cm}|}
\caption{MON-SHAFT --- Additional Shaft-Line Monitoring Submissions} \label{tab:monshaft} \\
\hline
\rowcolor{bvblue}
\textcolor{white}{\textbf{\#}} & \textcolor{white}{\textbf{Drawing/Document}} & \textcolor{white}{\textbf{Category}} & \textcolor{white}{\textbf{Content Requirements}} & \textcolor{white}{\textbf{Reference}} \\
\hline
\endfirsthead
\multicolumn{5}{c}%
{{\tablename\ \thetable{} -- continued from previous page}} \\
\hline
\rowcolor{bvblue}
\textcolor{white}{\textbf{\#}} & \textcolor{white}{\textbf{Drawing/Document}} & \textcolor{white}{\textbf{Category}} & \textcolor{white}{\textbf{Content Requirements}} & \textcolor{white}{\textbf{Reference}} \\
\hline
\endhead
\hline \multicolumn{5}{|r|}{{Continued on next page}} \\ \hline
\endfoot
\hline
\endlastfoot

1 & Oil-Lubricated Bearing \& Seal Arrangement & \fa{} &
\begin{itemize}[leftmargin=*,nosep]
\item Approved stern-tube / lower-unit oil-lubricated bearing \& seal design
\item Weardown measurement provisions
\item Compatibility with monitoring notation
\end{itemize} & NR467 Sec.7; Pt E (MON-SHAFT) \\
\hline

2 & Lube-Oil Monitoring \& Analysis Programme & \fa{} &
\begin{itemize}[leftmargin=*,nosep]
\item Oil temperature (and pressure where applicable) monitoring
\item Periodic oil-analysis programme (contamination, wear metals, water)
\item Sampling points and acceptance/alert criteria
\end{itemize} & NR467 Pt E (MON-SHAFT) \\
\hline

3 & Bearing Weardown / Sensor Instrumentation & \fa{} &
\begin{itemize}[leftmargin=*,nosep]
\item Bearing temperature sensors and locations
\item Weardown/liner-wear measurement method \& trending
\item Alarm set-points and data logging
\end{itemize} & NR467 Pt E; Pt C, Ch.3 \\
\hline

4 & Survey Arrangement \& Data-Management Plan & \fic{} &
\begin{itemize}[leftmargin=*,nosep]
\item Monitoring-based survey scheme replacing shaft withdrawal
\item Record keeping, trend review, thresholds for action
\item Interface with planned maintenance / class survey
\end{itemize} & NR467 Pt A (surveys); Pt E \\
\hline

\end{longtable}

\noindent\textbf{Surveyor note (MON-SHAFT):} MON-SHAFT is only compatible with
\textbf{oil-lubricated} shaft/bearing arrangements of an approved type. For
azimuth and podded units the ``shaft line'' extends into the lower gear/pod;
confirm the seal, bearing and oil-monitoring documentation of the thruster
lower unit (Chapters \ref{ch:azimuth}) is consistent with the MON-SHAFT scheme,
and that the oil-analysis and weardown-trending programme, sensors and
acceptance criteria are all submitted. Related shaft-line notations (e.g.\
oil-lubricated stern-tube monitoring / \texttt{TIGHT}-type arrangements) follow
the same principle --- verify the exact notation and its rule article in the
current NR467 Part E.

\newpage

% Chapter 7: Hydraulics
\chapter{Hydraulic Systems \& Pitch Control Mechanisms}
\label{ch:hydraulics}

\section{Applicability}

Hydraulic systems are used in thrusters for:
\begin{itemize}[leftmargin=*]
\item Hydraulic motor drive (alternative to electric motor)
\item Retractable thruster lifting/lowering mechanism
\item Controllable pitch propeller pitch control
\item Azimuth slewing drive (if hydraulic)
\item Water-jet steering and reversing actuation
\item Voith-Schneider control (steering) servo system
\end{itemize}

\noindent\textbf{Rule basis:} Thruster hydraulic systems are assessed under
\textbf{NR467 Pt C, Ch 1} --- the thruster/water-jet requirements of Sec.\,15,
the \textbf{piping-system rules} (design pressure, materials, testing, relief),
and, where the hydraulics actuate a \emph{steering} function, the more stringent
\textbf{steering-gear requirements of Sec.\,14} (duplication, single-failure,
power supply, response time). Recognised standard \textbf{ISO 4413} applies to
general hydraulic-system design. Steering-critical hydraulics attract the
redundancy and control provisions of Chapter \ref{ch:steering}; pitch-control
hydraulics must define the \emph{fail-safe} pitch position.

\section{Hydraulic Power Unit (HPU)}

\begin{longtable}{|p{0.5cm}|p{5cm}|p{1.5cm}|p{5.5cm}|p{2cm}|}
\caption{Hydraulic Systems - HPU} \label{tab:hydraulic-hpu} \\
\hline
\rowcolor{bvblue}
\textcolor{white}{\textbf{\#}} & \textcolor{white}{\textbf{Drawing/Document}} & \textcolor{white}{\textbf{Category}} & \textcolor{white}{\textbf{Content Requirements}} & \textcolor{white}{\textbf{Reference}} \\
\hline
\endfirsthead

\multicolumn{5}{c}%
{{\tablename\ \thetable{} -- continued from previous page}} \\
\hline
\rowcolor{bvblue}
\textcolor{white}{\textbf{\#}} & \textcolor{white}{\textbf{Drawing/Document}} & \textcolor{white}{\textbf{Category}} & \textcolor{white}{\textbf{Content Requirements}} & \textcolor{white}{\textbf{Reference}} \\
\hline
\endhead

\hline \multicolumn{5}{|r|}{{Continued on next page}} \\ \hline
\endfoot

\hline
\endlastfoot

1 & HPU General Arrangement & \fa{} & 
\begin{itemize}[leftmargin=*,nosep]
\item Complete HPU layout
\item Pump arrangement
\item Motor/drive arrangement
\item Reservoir design
\item Valve manifold
\item Accumulator location
\item Filtration system
\item Cooling system
\item Mounting arrangement
\end{itemize} & NR467 Sec.15 \\
\hline

2 & Hydraulic Pump Specification & \fa{} & 
\begin{itemize}[leftmargin=*,nosep]
\item Pump type (gear, piston, vane)
\item Displacement
\item Rated pressure and flow
\item Drive motor specifications
\item Efficiency curves
\item Manufacturer and model
\end{itemize} & NR467 Sec.15 \\
\hline

3 & Hydraulic Reservoir Design & \fa{} & 
\begin{itemize}[leftmargin=*,nosep]
\item Reservoir capacity
\item Material and construction
\item Baffle arrangement
\item Oil level indication
\item Breather and filler
\item Drain provisions
\item Heating elements (if applicable)
\end{itemize} & NR467 Sec.15 \\
\hline

4 & HPU Calculations & \fa{} & 
\begin{itemize}[leftmargin=*,nosep]
\item Flow requirements
\item Pressure requirements
\item Pump sizing verification
\item Reservoir sizing
\item Heat generation and cooling capacity
\item Accumulator sizing (if applicable)
\end{itemize} & NR467 Sec.15 \\
\hline

\end{longtable}

\section{Hydraulic Circuit Design}

\begin{longtable}{|p{0.5cm}|p{5cm}|p{1.5cm}|p{5.5cm}|p{2cm}|}
\caption{Hydraulic Systems - Circuit Design} \label{tab:hydraulic-circuit} \\
\hline
\rowcolor{bvblue}
\textcolor{white}{\textbf{\#}} & \textcolor{white}{\textbf{Drawing/Document}} & \textcolor{white}{\textbf{Category}} & \textcolor{white}{\textbf{Content Requirements}} & \textcolor{white}{\textbf{Reference}} \\
\hline
\endfirsthead

\multicolumn{5}{c}%
{{\tablename\ \thetable{} -- continued from previous page}} \\
\hline
\rowcolor{bvblue}
\textcolor{white}{\textbf{\#}} & \textcolor{white}{\textbf{Drawing/Document}} & \textcolor{white}{\textbf{Category}} & \textcolor{white}{\textbf{Content Requirements}} & \textcolor{white}{\textbf{Reference}} \\
\hline
\endhead

\hline \multicolumn{5}{|r|}{{Continued on next page}} \\ \hline
\endfoot

\hline
\endlastfoot

5 & Hydraulic System Schematic & \fa{} & 
\begin{itemize}[leftmargin=*,nosep]
\item Complete hydraulic circuit
\item All components identified with symbols
\item Pressure ratings
\item Flow ratings
\item Control logic
\item Safety devices
\item Monitoring points
\end{itemize} & NR467 Sec.15 \\
\hline

6 & Piping and Instrumentation Diagram (P\&ID) & \fa{} & 
\begin{itemize}[leftmargin=*,nosep]
\item Detailed piping layout
\item Pipe sizes and materials
\item Valve locations and specifications
\item Instrument locations
\item Pressure and temperature monitoring
\item Drain and vent points
\end{itemize} & NR467 Sec.15 \\
\hline

7 & Hydraulic Piping Arrangement & \fa{} & 
\begin{itemize}[leftmargin=*,nosep]
\item 3D piping layout (or isometric)
\item Pipe routing
\item Pipe supports
\item Flexibility analysis
\item Penetrations through bulkheads
\item Accessibility for maintenance
\end{itemize} & NR467 Sec.15 \\
\hline

\end{longtable}

\section{Hydraulic Components}

\begin{longtable}{|p{0.5cm}|p{5cm}|p{1.5cm}|p{5.5cm}|p{2cm}|}
\caption{Hydraulic Systems - Components} \label{tab:hydraulic-components} \\
\hline
\rowcolor{bvblue}
\textcolor{white}{\textbf{\#}} & \textcolor{white}{\textbf{Drawing/Document}} & \textcolor{white}{\textbf{Category}} & \textcolor{white}{\textbf{Content Requirements}} & \textcolor{white}{\textbf{Reference}} \\
\hline
\endfirsthead

\multicolumn{5}{c}%
{{\tablename\ \thetable{} -- continued from previous page}} \\
\hline
\rowcolor{bvblue}
\textcolor{white}{\textbf{\#}} & \textcolor{white}{\textbf{Drawing/Document}} & \textcolor{white}{\textbf{Category}} & \textcolor{white}{\textbf{Content Requirements}} & \textcolor{white}{\textbf{Reference}} \\
\hline
\endhead

\hline \multicolumn{5}{|r|}{{Continued on next page}} \\ \hline
\endfoot

\hline
\endlastfoot

8 & Hydraulic Motor Specification & \fa{} & 
\begin{itemize}[leftmargin=*,nosep]
\item Motor type and displacement
\item Rated pressure and speed
\item Torque and power curves
\item Efficiency
\item Mounting details
\item Manufacturer and model
\end{itemize} & NR467 Sec.15 \\
\hline

9 & Hydraulic Cylinder Specification & \fa{} & 
\begin{itemize}[leftmargin=*,nosep]
\item Cylinder bore and stroke
\item Operating pressure
\item Rod diameter and material
\item Seal specifications
\item Mounting arrangement
\item Position sensors
\item Manufacturer and model
\end{itemize} & NR467 Sec.15 \\
\hline

10 & Control Valve Specifications & \fa{} & 
\begin{itemize}[leftmargin=*,nosep]
\item Valve types (directional, pressure, flow)
\item Valve ratings
\item Actuation method (solenoid, manual, pilot)
\item Response time
\item Manufacturer and model
\end{itemize} & NR467 Sec.15 \\
\hline

11 & Pressure Relief Valve & \fa{} & 
\begin{itemize}[leftmargin=*,nosep]
\item Set pressure
\item Flow capacity
\item Type (direct-acting, pilot-operated)
\item Mounting location
\item Manufacturer and model
\end{itemize} & NR467 Sec.15 \\
\hline

12 & Accumulator Specification & \fa{} & 
\begin{itemize}[leftmargin=*,nosep]
\item Accumulator type (bladder, piston, diaphragm)
\item Volume
\item Pre-charge pressure
\item Operating pressure range
\item Gas type (nitrogen)
\item Mounting arrangement
\item Safety provisions
\end{itemize} & NR467 Sec.15 \\
\hline

13 & Accumulator Sizing Calculations & \fa{} & 
\begin{itemize}[leftmargin=*,nosep]
\item Volume requirements
\item Pre-charge pressure calculation
\item Energy storage capacity
\item Response time analysis
\end{itemize} & NR467 Sec.15 \\
\hline

\end{longtable}

\section{Filtration System}

\begin{longtable}{|p{0.5cm}|p{5cm}|p{1.5cm}|p{5.5cm}|p{2cm}|}
\caption{Hydraulic Systems - Filtration} \label{tab:hydraulic-filtration} \\
\hline
\rowcolor{bvblue}
\textcolor{white}{\textbf{\#}} & \textcolor{white}{\textbf{Drawing/Document}} & \textcolor{white}{\textbf{Category}} & \textcolor{white}{\textbf{Content Requirements}} & \textcolor{white}{\textbf{Reference}} \\
\hline
\endfirsthead

\multicolumn{5}{c}%
{{\tablename\ \thetable{} -- continued from previous page}} \\
\hline
\rowcolor{bvblue}
\textcolor{white}{\textbf{\#}} & \textcolor{white}{\textbf{Drawing/Document}} & \textcolor{white}{\textbf{Category}} & \textcolor{white}{\textbf{Content Requirements}} & \textcolor{white}{\textbf{Reference}} \\
\hline
\endhead

\hline \multicolumn{5}{|r|}{{Continued on next page}} \\ \hline
\endfoot

\hline
\endlastfoot

14 & Filtration System Design & \fa{} & 
\begin{itemize}[leftmargin=*,nosep]
\item Filter types and locations
\item Filtration rating (micron)
\item Flow capacity
\item Pressure drop
\item Filter change indicators
\item Bypass provisions
\end{itemize} & NR467 Sec.15 \\
\hline

15 & Filter Specifications & \fa{} & 
\begin{itemize}[leftmargin=*,nosep]
\item Filter element type
\item Filtration rating
\item Collapse pressure
\item Manufacturer and model
\item Replacement schedule
\end{itemize} & NR467 Sec.15 \\
\hline

\end{longtable}

\section{Cooling System}

\begin{longtable}{|p{0.5cm}|p{5cm}|p{1.5cm}|p{5.5cm}|p{2cm}|}
\caption{Hydraulic Systems - Cooling} \label{tab:hydraulic-cooling} \\
\hline
\rowcolor{bvblue}
\textcolor{white}{\textbf{\#}} & \textcolor{white}{\textbf{Drawing/Document}} & \textcolor{white}{\textbf{Category}} & \textcolor{white}{\textbf{Content Requirements}} & \textcolor{white}{\textbf{Reference}} \\
\hline
\endfirsthead

\multicolumn{5}{c}%
{{\tablename\ \thetable{} -- continued from previous page}} \\
\hline
\rowcolor{bvblue}
\textcolor{white}{\textbf{\#}} & \textcolor{white}{\textbf{Drawing/Document}} & \textcolor{white}{\textbf{Category}} & \textcolor{white}{\textbf{Content Requirements}} & \textcolor{white}{\textbf{Reference}} \\
\hline
\endhead

\hline \multicolumn{5}{|r|}{{Continued on next page}} \\ \hline
\endfoot

\hline
\endlastfoot

16 & Hydraulic Oil Cooling System & \fa{} & 
\begin{itemize}[leftmargin=*,nosep]
\item Cooling method (air, water)
\item Heat exchanger specifications
\item Cooling capacity
\item Temperature control
\item Monitoring provisions
\end{itemize} & NR467 Sec.15 \\
\hline

17 & Thermal Analysis & \fa{} & 
\begin{itemize}[leftmargin=*,nosep]
\item Heat generation calculations
\item Heat dissipation calculations
\item Temperature rise predictions
\item Cooling capacity verification
\item Operating temperature range
\end{itemize} & NR467 Sec.15 \\
\hline

\end{longtable}

\section{Hydraulic Fluid}

\begin{longtable}{|p{0.5cm}|p{5cm}|p{1.5cm}|p{5.5cm}|p{2cm}|}
\caption{Hydraulic Systems - Fluid Specification} \label{tab:hydraulic-fluid} \\
\hline
\rowcolor{bvblue}
\textcolor{white}{\textbf{\#}} & \textcolor{white}{\textbf{Drawing/Document}} & \textcolor{white}{\textbf{Category}} & \textcolor{white}{\textbf{Content Requirements}} & \textcolor{white}{\textbf{Reference}} \\
\hline
\endfirsthead

\multicolumn{5}{c}%
{{\tablename\ \thetable{} -- continued from previous page}} \\
\hline
\rowcolor{bvblue}
\textcolor{white}{\textbf{\#}} & \textcolor{white}{\textbf{Drawing/Document}} & \textcolor{white}{\textbf{Category}} & \textcolor{white}{\textbf{Content Requirements}} & \textcolor{white}{\textbf{Reference}} \\
\hline
\endhead

\hline \multicolumn{5}{|r|}{{Continued on next page}} \\ \hline
\endfoot

\hline
\endlastfoot

18 & Hydraulic Fluid Specification & \fa{} & 
\begin{itemize}[leftmargin=*,nosep]
\item Fluid type (mineral oil, synthetic, biodegradable)
\item Viscosity grade
\item Operating temperature range
\item Fire resistance properties (if applicable)
\item Compatibility with system materials
\item Approved manufacturer and grade
\end{itemize} & NR467 Sec.15 \\
\hline

\end{longtable}

\section{Piping Design}

\begin{longtable}{|p{0.5cm}|p{5cm}|p{1.5cm}|p{5.5cm}|p{2cm}|}
\caption{Hydraulic Systems - Piping} \label{tab:hydraulic-piping} \\
\hline
\rowcolor{bvblue}
\textcolor{white}{\textbf{\#}} & \textcolor{white}{\textbf{Drawing/Document}} & \textcolor{white}{\textbf{Category}} & \textcolor{white}{\textbf{Content Requirements}} & \textcolor{white}{\textbf{Reference}} \\
\hline
\endfirsthead

\multicolumn{5}{c}%
{{\tablename\ \thetable{} -- continued from previous page}} \\
\hline
\rowcolor{bvblue}
\textcolor{white}{\textbf{\#}} & \textcolor{white}{\textbf{Drawing/Document}} & \textcolor{white}{\textbf{Category}} & \textcolor{white}{\textbf{Content Requirements}} & \textcolor{white}{\textbf{Reference}} \\
\hline
\endhead

\hline \multicolumn{5}{|r|}{{Continued on next page}} \\ \hline
\endfoot

\hline
\endlastfoot

19 & Pipe Sizing Calculations & \fa{} & 
\begin{itemize}[leftmargin=*,nosep]
\item Flow velocity calculations
\item Pressure drop calculations
\item Pipe diameter selection
\item Compliance with velocity limits
\end{itemize} & NR467 Sec.15 \\
\hline

20 & Pipe Material Specifications & \fa{} & 
\begin{itemize}[leftmargin=*,nosep]
\item Pipe material (steel, stainless steel, copper-nickel)
\item Wall thickness
\item Pressure rating
\item Corrosion protection
\item Compliance with BV approved materials
\end{itemize} & NR467 Sec.15; Part B, Ch.3 \\
\hline

21 & Pipe Support Design & \fa{} & 
\begin{itemize}[leftmargin=*,nosep]
\item Support locations and types
\item Support spacing
\item Vibration isolation
\item Thermal expansion accommodation
\item Load calculations
\end{itemize} & NR467 Sec.15 \\
\hline

22 & Piping Stress Analysis & \fa{} & 
\begin{itemize}[leftmargin=*,nosep]
\item Pressure stress
\item Thermal stress
\item Vibration analysis
\item Support reaction forces
\item Compliance with stress limits
\end{itemize} & NR467 Sec.15 \\
\hline

\end{longtable}

\section{Hose Assemblies}

\begin{longtable}{|p{0.5cm}|p{5cm}|p{1.5cm}|p{5.5cm}|p{2cm}|}
\caption{Hydraulic Systems - Hoses} \label{tab:hydraulic-hoses} \\
\hline
\rowcolor{bvblue}
\textcolor{white}{\textbf{\#}} & \textcolor{white}{\textbf{Drawing/Document}} & \textcolor{white}{\textbf{Category}} & \textcolor{white}{\textbf{Content Requirements}} & \textcolor{white}{\textbf{Reference}} \\
\hline
\endfirsthead

\multicolumn{5}{c}%
{{\tablename\ \thetable{} -- continued from previous page}} \\
\hline
\rowcolor{bvblue}
\textcolor{white}{\textbf{\#}} & \textcolor{white}{\textbf{Drawing/Document}} & \textcolor{white}{\textbf{Category}} & \textcolor{white}{\textbf{Content Requirements}} & \textcolor{white}{\textbf{Reference}} \\
\hline
\endhead

\hline \multicolumn{5}{|r|}{{Continued on next page}} \\ \hline
\endfoot

\hline
\endlastfoot

23 & Hydraulic Hose Specifications & \fa{} & 
\begin{itemize}[leftmargin=*,nosep]
\item Hose type and construction
\item Pressure rating
\item Temperature rating
\item End fittings
\item Length and routing
\item Fire resistance (if required)
\item Manufacturer and model
\end{itemize} & NR467 Sec.15 \\
\hline

24 & Hose Installation Drawing & \fa{} & 
\begin{itemize}[leftmargin=*,nosep]
\item Hose routing
\item Bend radius verification
\item Support and protection
\item Accessibility for inspection/replacement
\end{itemize} & NR467 Sec.15 \\
\hline

\end{longtable}

\section{Safety Devices}

\begin{longtable}{|p{0.5cm}|p{5cm}|p{1.5cm}|p{5.5cm}|p{2cm}|}
\caption{Hydraulic Systems - Safety Devices} \label{tab:hydraulic-safety} \\
\hline
\rowcolor{bvblue}
\textcolor{white}{\textbf{\#}} & \textcolor{white}{\textbf{Drawing/Document}} & \textcolor{white}{\textbf{Category}} & \textcolor{white}{\textbf{Content Requirements}} & \textcolor{white}{\textbf{Reference}} \\
\hline
\endfirsthead

\multicolumn{5}{c}%
{{\tablename\ \thetable{} -- continued from previous page}} \\
\hline
\rowcolor{bvblue}
\textcolor{white}{\textbf{\#}} & \textcolor{white}{\textbf{Drawing/Document}} & \textcolor{white}{\textbf{Category}} & \textcolor{white}{\textbf{Content Requirements}} & \textcolor{white}{\textbf{Reference}} \\
\hline
\endhead

\hline \multicolumn{5}{|r|}{{Continued on next page}} \\ \hline
\endfoot

\hline
\endlastfoot

25 & Pressure Relief System & \fa{} & 
\begin{itemize}[leftmargin=*,nosep]
\item Relief valve locations
\item Set pressures
\item Discharge routing
\item Capacity verification
\end{itemize} & NR467 Sec.15 \\
\hline

26 & Emergency Shutdown System & \fa{} & 
\begin{itemize}[leftmargin=*,nosep]
\item Emergency stop provisions
\item Valve actuation
\item System depressurization
\item Safe state definition
\end{itemize} & NR467 Sec.15 \\
\hline

27 & Leak Detection System & \fa{} & 
\begin{itemize}[leftmargin=*,nosep]
\item Leak detection method
\item Alarm provisions
\item Automatic shutdown (if required)
\end{itemize} & NR467 Sec.15 \\
\hline

\end{longtable}

\section{Control System}

\begin{longtable}{|p{0.5cm}|p{5cm}|p{1.5cm}|p{5.5cm}|p{2cm}|}
\caption{Hydraulic Systems - Control} \label{tab:hydraulic-control} \\
\hline
\rowcolor{bvblue}
\textcolor{white}{\textbf{\#}} & \textcolor{white}{\textbf{Drawing/Document}} & \textcolor{white}{\textbf{Category}} & \textcolor{white}{\textbf{Content Requirements}} & \textcolor{white}{\textbf{Reference}} \\
\hline
\endfirsthead

\multicolumn{5}{c}%
{{\tablename\ \thetable{} -- continued from previous page}} \\
\hline
\rowcolor{bvblue}
\textcolor{white}{\textbf{\#}} & \textcolor{white}{\textbf{Drawing/Document}} & \textcolor{white}{\textbf{Category}} & \textcolor{white}{\textbf{Content Requirements}} & \textcolor{white}{\textbf{Reference}} \\
\hline
\endhead

\hline \multicolumn{5}{|r|}{{Continued on next page}} \\ \hline
\endfoot

\hline
\endlastfoot

28 & Hydraulic Control System & \fa{} & 
\begin{itemize}[leftmargin=*,nosep]
\item Control architecture
\item Manual/automatic control modes
\item Valve actuation signals
\item Feedback systems
\item Safety interlocks
\end{itemize} & NR467 Sec.15; Part C, Ch.4 \\
\hline

29 & Control Panel Layout & \fa{} & 
\begin{itemize}[leftmargin=*,nosep]
\item Panel arrangement
\item Control devices
\item Indication and monitoring
\item Alarm displays
\item Emergency stop
\end{itemize} & NR467 Sec.15 \\
\hline

\end{longtable}

\section{Monitoring \& Instrumentation}

\begin{longtable}{|p{0.5cm}|p{5cm}|p{1.5cm}|p{5.5cm}|p{2cm}|}
\caption{Hydraulic Systems - Monitoring} \label{tab:hydraulic-monitoring} \\
\hline
\rowcolor{bvblue}
\textcolor{white}{\textbf{\#}} & \textcolor{white}{\textbf{Drawing/Document}} & \textcolor{white}{\textbf{Category}} & \textcolor{white}{\textbf{Content Requirements}} & \textcolor{white}{\textbf{Reference}} \\
\hline
\endfirsthead

\multicolumn{5}{c}%
{{\tablename\ \thetable{} -- continued from previous page}} \\
\hline
\rowcolor{bvblue}
\textcolor{white}{\textbf{\#}} & \textcolor{white}{\textbf{Drawing/Document}} & \textcolor{white}{\textbf{Category}} & \textcolor{white}{\textbf{Content Requirements}} & \textcolor{white}{\textbf{Reference}} \\
\hline
\endhead

\hline \multicolumn{5}{|r|}{{Continued on next page}} \\ \hline
\endfoot

\hline
\endlastfoot

30 & Instrumentation List & \fa{} & 
\begin{itemize}[leftmargin=*,nosep]
\item Pressure gauges/transmitters
\item Temperature sensors
\item Flow meters (if applicable)
\item Level indicators
\item Position sensors
\item Filter condition indicators
\item Specifications for each instrument
\end{itemize} & NR467 Sec.15 \\
\hline

31 & Alarm and Monitoring System & \fa{} & 
\begin{itemize}[leftmargin=*,nosep]
\item Parameters monitored
\item Alarm setpoints
\item Alarm display location
\item Data logging (if applicable)
\end{itemize} & NR467 Sec.15 \\
\hline

\end{longtable}

\section{Controllable Pitch Mechanism (CP Propellers)}

\begin{longtable}{|p{0.5cm}|p{5cm}|p{1.5cm}|p{5.5cm}|p{2cm}|}
\caption{Hydraulic Systems - CP Mechanism} \label{tab:hydraulic-cp} \\
\hline
\rowcolor{bvblue}
\textcolor{white}{\textbf{\#}} & \textcolor{white}{\textbf{Drawing/Document}} & \textcolor{white}{\textbf{Category}} & \textcolor{white}{\textbf{Content Requirements}} & \textcolor{white}{\textbf{Reference}} \\
\hline
\endfirsthead

\multicolumn{5}{c}%
{{\tablename\ \thetable{} -- continued from previous page}} \\
\hline
\rowcolor{bvblue}
\textcolor{white}{\textbf{\#}} & \textcolor{white}{\textbf{Drawing/Document}} & \textcolor{white}{\textbf{Category}} & \textcolor{white}{\textbf{Content Requirements}} & \textcolor{white}{\textbf{Reference}} \\
\hline
\endhead

\hline \multicolumn{5}{|r|}{{Continued on next page}} \\ \hline
\endfoot

\hline
\endlastfoot

32 & CP Mechanism Assembly & \fa{} & 
\begin{itemize}[leftmargin=*,nosep]
\item Hub internal mechanism
\item Pitch control rod/linkage
\item Hydraulic cylinder (in hub or external)
\item Blade retention system
\item Bearing arrangement
\item Seal design
\end{itemize} & NR467 Sec.15, [3.4] \\
\hline

33 & CP Mechanism Calculations & \fa{} & 
\begin{itemize}[leftmargin=*,nosep]
\item Pitch control forces
\item Hydraulic cylinder sizing
\item Linkage strength
\item Blade retention strength
\item Seal pressure rating
\end{itemize} & NR467 Sec.15, [3.4] \\
\hline

34 & Pitch Control Hydraulic System & \fa{} & 
\begin{itemize}[leftmargin=*,nosep]
\item Hydraulic supply to hub
\item Rotary joint/slip ring design
\item Control valve system
\item Pitch feedback system
\item Emergency pitch setting
\end{itemize} & NR467 Sec.15, [3.4] \\
\hline

35 & Pitch Feedback System & \fa{} & 
\begin{itemize}[leftmargin=*,nosep]
\item Position sensor type and location
\item Signal transmission (rotary joint)
\item Accuracy and resolution
\item Redundancy (if required)
\end{itemize} & NR467 Sec.15, [3.4] \\
\hline

\end{longtable}

\section{Testing \& Commissioning}

\begin{longtable}{|p{0.5cm}|p{5cm}|p{1.5cm}|p{5.5cm}|p{2cm}|}
\caption{Hydraulic Systems - Testing} \label{tab:hydraulic-testing} \\
\hline
\rowcolor{bvblue}
\textcolor{white}{\textbf{\#}} & \textcolor{white}{\textbf{Drawing/Document}} & \textcolor{white}{\textbf{Category}} & \textcolor{white}{\textbf{Content Requirements}} & \textcolor{white}{\textbf{Reference}} \\
\hline
\endfirsthead

\multicolumn{5}{c}%
{{\tablename\ \thetable{} -- continued from previous page}} \\
\hline
\rowcolor{bvblue}
\textcolor{white}{\textbf{\#}} & \textcolor{white}{\textbf{Drawing/Document}} & \textcolor{white}{\textbf{Category}} & \textcolor{white}{\textbf{Content Requirements}} & \textcolor{white}{\textbf{Reference}} \\
\hline
\endhead

\hline \multicolumn{5}{|r|}{{Continued on next page}} \\ \hline
\endfoot

\hline
\endlastfoot

36 & Hydraulic System Test Procedures & \fa{} & 
\begin{itemize}[leftmargin=*,nosep]
\item Pressure testing
\item Functional testing
\item Leak testing
\item Performance verification
\item Safety device testing
\item Acceptance criteria
\end{itemize} & NR467 Sec.15 \\
\hline

37 & Hydraulic System Test Reports & \fic{} & 
\begin{itemize}[leftmargin=*,nosep]
\item Pressure test results
\item Functional test results
\item Leak test results
\item Performance data
\item As-built documentation
\end{itemize} & NR467 Sec.15 \\
\hline

\end{longtable}

\section{Operating \& Maintenance Manuals}

\begin{longtable}{|p{0.5cm}|p{5cm}|p{1.5cm}|p{5.5cm}|p{2cm}|}
\caption{Hydraulic Systems - Manuals} \label{tab:hydraulic-manuals} \\
\hline
\rowcolor{bvblue}
\textcolor{white}{\textbf{\#}} & \textcolor{white}{\textbf{Drawing/Document}} & \textcolor{white}{\textbf{Category}} & \textcolor{white}{\textbf{Content Requirements}} & \textcolor{white}{\textbf{Reference}} \\
\hline
\endfirsthead

\multicolumn{5}{c}%
{{\tablename\ \thetable{} -- continued from previous page}} \\
\hline
\rowcolor{bvblue}
\textcolor{white}{\textbf{\#}} & \textcolor{white}{\textbf{Drawing/Document}} & \textcolor{white}{\textbf{Category}} & \textcolor{white}{\textbf{Content Requirements}} & \textcolor{white}{\textbf{Reference}} \\
\hline
\endhead

\hline \multicolumn{5}{|r|}{{Continued on next page}} \\ \hline
\endfoot

\hline
\endlastfoot

38 & Hydraulic System Operating Manual & \fic{} & 
\begin{itemize}[leftmargin=*,nosep]
\item System description
\item Operating procedures
\item Normal operation
\item Emergency procedures
\item Troubleshooting guide
\item Safety precautions
\end{itemize} & NR467 Sec.15 \\
\hline

39 & Hydraulic System Maintenance Manual & \fic{} & 
\begin{itemize}[leftmargin=*,nosep]
\item Maintenance schedule
\item Inspection procedures
\item Fluid sampling and analysis
\item Filter replacement
\item Component overhaul
\item Spare parts list
\end{itemize} & NR467 Sec.15 \\
\hline

\end{longtable}

\newpage

% Due to length, I'll provide the remaining chapters in a condensed format
% Continue with Chapter 8 (DP), Chapter 9 (Testing), Chapter 10 (Checklists),
% Chapter 11 (Errors), Chapter 12 (Quick Reference), Chapter 13 (Regulatory References)

% Chapter 8: DP Thrusters (condensed)
\chapter{Dynamic Positioning (DP) Thrusters - Additional Requirements}
\label{ch:dp}

\section{Applicability}

Thrusters that are part of a Dynamic Positioning (DP) system must comply with ALL applicable requirements from Chapters \ref{ch:transverse}-\ref{ch:hydraulics}, PLUS the additional DP-specific requirements in this chapter.

DP systems are classified into three classes:
\begin{itemize}[leftmargin=*]
\item \textbf{DP-1}: Basic DP capability, no redundancy required
\item \textbf{DP-2}: Redundant systems, single failure tolerance
\item \textbf{DP-3}: Enhanced redundancy with fire/flood separation
\end{itemize}

The DP-specific submissions below are \emph{in addition to} the full type
package for each thruster used in the DP system (structure, propeller, shaft,
gear, drive, control). They focus on redundancy, power segregation, control and
failure analysis, and are assessed principally under \textbf{NR467 Part F}
(Additional Class Notation \textbf{DYNAPOS}).

\textbf{Key DP-Specific Documents Required (\fa{}):}
\begin{enumerate}
\item DP System General Arrangement
\item Thruster Configuration Drawing
\item DP Capability Plot
\item FMEA Documentation (comprehensive)
\item Redundancy Concept (DP-2/DP-3)
\item Power Supply Segregation
\item Fire Zone Arrangement (DP-3)
\item Watertight Subdivision (DP-3)
\item DP Control System Architecture
\item Thruster Allocation Logic
\item Position Reference System Configuration
\item Environmental Sensor Configuration
\item Power Management System (PMS) Configuration
\item UPS Configuration
\item DP FMEA Proving Trials Plan
\item DP Annual Trials Plan
\end{enumerate}

\textbf{Key DP-Specific Documents Required (\fic{}):}
\begin{enumerate}
\item DP Operations Manual
\item DP Maintenance Manual
\item DP Trials Reports
\item DP Training Requirements
\end{enumerate}

\section{Reference}
All DP (DYNAPOS) requirements per \textbf{NR467 Part F (Additional Class Notations)}; see also IMO MSC.1/Circ.1580 and IMCA guidance.

\newpage

% Chapter: Ice Class
\chapter{Ice Class \& Cold-Climate (Winterisation) Requirements}
\label{ch:ice}

\section{Applicability}

When a vessel carries a BV ice-class or polar notation, thrusters and their
drive trains must be strengthened and winterised in addition to all base
requirements of the preceding chapters. Two families apply:
\begin{itemize}[leftmargin=*]
\item \textbf{Baltic / Finnish-Swedish Ice Classes} --- notations
\textbf{ICE CLASS IA SUPER, IA, IB, IC} (and \textbf{ID}), based on the
Finnish-Swedish Ice Class Rules (FSICR / TRAFICOM), adopted in
\textbf{NR467 Part E}. Intended for first-year Baltic-type ice with icebreaker
assistance.
\item \textbf{Polar Classes} --- notations \textbf{POLAR CLASS PC1--PC7}
(IACS Polar Class / IMO Polar Code), addressed in \textbf{NR467 Part E} and
Rule Note \textbf{NR527}. Machinery strengthening follows \textbf{IACS UR I3}.
\end{itemize}
Cold-climate service without ice loads is covered by winterisation notations such
as \textbf{COLD} (e.g.\ \texttt{COLD(H\,$t_{DH}$, E\,$t_{DE}$)}), addressing
design ambient temperatures rather than ice impact.

\textbf{Key principle:} ice class primarily drives \emph{ice-load strengthening}
of the propeller, shaft line, gears and thruster body, plus \emph{winterisation}
of exposed and seawater systems. The surveyor overlays these on the type package.

\section{Ice-Strengthening Submissions}

\begin{longtable}{|p{0.5cm}|p{5cm}|p{1.5cm}|p{5.5cm}|p{2cm}|}
\caption{Ice Class --- Additional Thruster Requirements} \label{tab:ice} \\
\hline
\rowcolor{bvblue}
\textcolor{white}{\textbf{\#}} & \textcolor{white}{\textbf{Drawing/Document}} & \textcolor{white}{\textbf{Category}} & \textcolor{white}{\textbf{Content Requirements}} & \textcolor{white}{\textbf{Reference}} \\
\hline
\endfirsthead
\multicolumn{5}{c}%
{{\tablename\ \thetable{} -- continued from previous page}} \\
\hline
\rowcolor{bvblue}
\textcolor{white}{\textbf{\#}} & \textcolor{white}{\textbf{Drawing/Document}} & \textcolor{white}{\textbf{Category}} & \textcolor{white}{\textbf{Content Requirements}} & \textcolor{white}{\textbf{Reference}} \\
\hline
\endhead
\hline \multicolumn{5}{|r|}{{Continued on next page}} \\ \hline
\endfoot
\hline
\endlastfoot

1 & Ice Class Notation Basis \& Design Conditions & \fa{} &
\begin{itemize}[leftmargin=*,nosep]
\item Declared notation (IA Super/IA/\ldots\ or PC1--PC7)
\item Design ice thickness / ice-load basis
\item Applicable rule edition and lower design temperature
\end{itemize} & NR467 Part F; NR527 \\
\hline

2 & Propeller Ice-Load Strength Calculation & \fa{} &
\begin{itemize}[leftmargin=*,nosep]
\item Ice interaction loads (milling, impact) on blades
\item Blade root and blade tip strength vs.\ ice bending
\item Blade material and thickness increase
\item Open vs.\ ducted / nozzle ice loads
\end{itemize} & NR467 Part E; IACS UR I3 \\
\hline

3 & Shaft Line \& Ice Torque Verification & \fa{} &
\begin{itemize}[leftmargin=*,nosep]
\item Maximum ice-induced torque and thrust
\item Propeller shaft, intermediate shaft, couplings
\item Overload/peak-torque capacity of gears
\item Blade-loss / one-blade-milling scenarios
\end{itemize} & NR467 Part E; Sec.7; IACS UR I3 \\
\hline

4 & Gear \& Azimuth/Nozzle Ice-Load Substantiation & \fa{} &
\begin{itemize}[leftmargin=*,nosep]
\item Bevel/reduction gears vs.\ ice peak torque
\item Azimuth slewing bearing/gear vs.\ ice-block impact on lower unit
\item Nozzle/steering-grid ice-impact strength
\end{itemize} & NR467 Part E; Sec.15 \\
\hline

5 & Thruster Body / Lower Unit Ice Protection & \fa{} &
\begin{itemize}[leftmargin=*,nosep]
\item Lower gearbox/pod shape vs.\ ice interaction
\item Ice knife / deflector where required
\item Local structural strengthening
\end{itemize} & NR467 Part E \\
\hline

6 & Tunnel / Inlet Ice \& Grid Arrangement & \fa{} &
\begin{itemize}[leftmargin=*,nosep]
\item Tunnel thruster grid ice/debris protection
\item De-icing / anti-blocking provisions
\item Ice-load on grid bars
\end{itemize} & NR467 Part E; Sec.15 \\
\hline

7 & Winterisation of Systems (Hydraulic, Cooling, Seals) & \fa{} &
\begin{itemize}[leftmargin=*,nosep]
\item Low-temperature grades of oils/greases/seals
\item Heating/insulation of exposed hydraulics \& sensors
\item Sea-chest/cooling anti-icing (recirculation)
\item Materials for low-temperature toughness (Charpy)
\end{itemize} & NR467 Part E (COLD); Part B, Ch.3 \\
\hline

8 & Materials Low-Temperature Certification & \fa{} &
\begin{itemize}[leftmargin=*,nosep]
\item Impact testing at design temperature
\item Grade selection for exposed \& load-bearing parts
\end{itemize} & NR467 Part B, Ch.3; Part E \\
\hline

9 & Retraction/Deployment in Ice (if retractable) & \fic{} &
\begin{itemize}[leftmargin=*,nosep]
\item Operating limits and procedures in ice
\item Protection of opening/seals from ice ingress
\end{itemize} & NR467 Part E \\
\hline

\end{longtable}

\noindent\textbf{Surveyor note (ice):} The governing case for ice-classed
propulsion is usually the \textbf{ice peak torque / blade-milling load}
propagating through propeller $\rightarrow$ shaft $\rightarrow$ gear
$\rightarrow$ prime mover, \emph{not} the open-water rating. Verify the ice-load
calculation explicitly and check that it is carried consistently into the shaft
(Chapter \ref{ch:shafting}) and gear strength checks. For azimuth/podded units in
ice, confirm lower-unit and slewing-gear substantiation against ice-block impact
--- a frequent omission. Match the rule edition to the declared notation.

\newpage

% Chapter: Additional Notations
\chapter{Additional BV Class Notations Affecting Thruster Packages}
\label{ch:notations}

\section{Purpose}

Beyond ice class, several BV \textbf{additional class notations} and
\textbf{service notations} impose extra requirements on the thruster package.
This chapter summarises the ones most often encountered so the surveyor can
confirm the client has submitted the corresponding documents. Requirements here
\emph{add to} the type package; they do not replace it.

\section{Notation Impact Matrix}

\begin{longtable}{|p{2.8cm}|p{3.7cm}|p{5.8cm}|p{1.9cm}|}
\caption{Additional Notations --- Impact on Thruster Submissions} \label{tab:notations} \\
\hline
\rowcolor{bvblue}
\textcolor{white}{\textbf{Notation (family)}} & \textcolor{white}{\textbf{Concern}} & \textcolor{white}{\textbf{Additional thruster submissions}} & \textcolor{white}{\textbf{Reference}} \\
\hline
\endfirsthead
\multicolumn{4}{c}%
{{\tablename\ \thetable{} -- continued from previous page}} \\
\hline
\rowcolor{bvblue}
\textcolor{white}{\textbf{Notation (family)}} & \textcolor{white}{\textbf{Concern}} & \textcolor{white}{\textbf{Additional thruster submissions}} & \textcolor{white}{\textbf{Reference}} \\
\hline
\endhead
\hline \multicolumn{4}{|r|}{{Continued on next page}} \\ \hline
\endfoot
\hline
\endlastfoot

\textbf{DYNAPOS} (AM/AT/ATR, SAM/SAT); \textbf{DP-1/2/3} & Dynamic positioning &
Full DP package --- see Chapter \ref{ch:dp}: FMEA, redundancy, power segregation,
capability plot, thruster allocation, proving trials. & NR467 Pt F (DYNAPOS) \\
\hline

\textbf{AUT-UMS, AUT-CCS, AUT-PORT} & Automated / unattended machinery &
Extended alarm \& monitoring of thruster (bearing temps, oil, vibration), remote
\& safe control, fail-safe on power/control loss, integration with alarm system.
& NR467 Pt C, Ch.3; Pt E \\
\hline

\textbf{COMF-NOISE / COMF-VIB} (Comfort) & Noise \& vibration on board &
Structure-borne \& airborne noise prediction, vibration levels at thruster
excitation (blade-rate), resilient mounting, cavitation-noise control, source
data. & NR467 Pt F (COMF) \\
\hline

\textbf{CLEANSHIP / environmental} & Pollution prevention &
Environmentally Acceptable Lubricants (EAL) for stern-tube/thruster seals \&
hydraulics, leakage monitoring, oil-to-sea interface documentation. & NR467 Pt E
(environmental); VGP/EAL \\
\hline

\textbf{Redundant Propulsion (RP, RP-1, AVM-\ldots)} & Propulsion availability &
Independence/segregation of propulsion trains, single-failure analysis,
take-me-home capability, redundant thruster arrangement. & NR467 Pt E \\
\hline

\textbf{MON-SHAFT} (shaft-line monitoring) & Shaft survey by condition monitoring &
Oil-lubricated bearing/seal of approved type; lube-oil monitoring \& analysis
programme; bearing temperature \& weardown instrumentation; monitoring-based
survey scheme (see Section \ref{sec:monshaft}). & NR467 Pt E; Sec.7 \\
\hline

\textbf{TUG / ESCORT TUG / AHT} (service) & High steering \& towing loads &
Escort/indirect-towing loads on azimuth \& steering system, structural
substantiation of unit \& foundation, stability interaction, uprated steering
duty. & NR467 Pt D \\
\hline

\textbf{SPS} (Special Purpose Ship) & Personnel / mission systems &
Higher survivability \& redundancy expectations may extend to propulsion/steering
arrangement. & NR467 Pt D/E \\
\hline

\textbf{Electric hybrid / battery} & Hybrid power &
Documentation of every driving path (PTI/PTO), energy-storage interface, mode
changeover, converter data. & NR467 Pt E \\
\hline

\end{longtable}

\noindent\textbf{Surveyor note (notations):} Always cross-check the vessel's
\emph{requested class notation string} against the thruster package. A DP or
comfort notation in particular pulls in substantial extra documentation
(FMEA/proving trials; noise \& vibration studies) that clients frequently omit
from the initial thruster submission. Where a notation adds a steering-critical
duty (escort tug), confirm the steering-gear package (Chapter \ref{ch:steering})
reflects the \emph{uprated} loads, not just the open-water condition.

\newpage

% Chapter: Ship-Type / Service Notation
\chapter{Ship-Type \& Service-Notation Specific Requirements}
\label{ch:shiptype}

\section{Purpose}

The base thruster package (Chapters \ref{ch:general}--\ref{ch:othertypes}) and the
prime-mover, ice and notation overlays (Chapters
\ref{ch:primemover}, \ref{ch:ice}, \ref{ch:notations}) define \emph{what a
thruster needs}. This chapter answers the complementary question the surveyor
must also ask: \textbf{``given the ship type / service notation, what
\emph{extra} do I expect in the thruster package, and which \emph{additional ship
rules} apply?''}

BV assigns each ship a \textbf{service notation} (NR467 \textbf{Part D}) ---
e.g.\ \texttt{tug}, \texttt{supply}, \texttt{oil tanker}, \texttt{passenger
ship} --- plus \textbf{additional class notations} (Part E). Together these pull
in requirements beyond NR467 Sec.\,15, and sometimes bring in \emph{separate
rule sets} (offshore units, naval ships, high-speed craft). The matrix below is
the surveyor's ``what else to demand'' guide; the subsequent sections expand the
cases with the most rule impact.

\section{Ship-Type $\rightarrow$ Thruster Expectation Matrix}

\begin{longtable}{|p{2.9cm}|p{3.7cm}|p{5.9cm}|p{1.8cm}|}
\caption{Ship Type / Service --- Typical Arrangement \& Additional Requirements} \label{tab:shiptype} \\
\hline
\rowcolor{bvblue}
\textcolor{white}{\textbf{Ship type / service}} & \textcolor{white}{\textbf{Typical thruster arrangement}} & \textcolor{white}{\textbf{Additional requirements \& notations to expect}} & \textcolor{white}{\textbf{Reference}} \\
\hline
\endfirsthead
\multicolumn{4}{c}%
{{\tablename\ \thetable{} -- continued from previous page}} \\
\hline
\rowcolor{bvblue}
\textcolor{white}{\textbf{Ship type / service}} & \textcolor{white}{\textbf{Typical thruster arrangement}} & \textcolor{white}{\textbf{Additional requirements \& notations to expect}} & \textcolor{white}{\textbf{Reference}} \\
\hline
\endhead
\hline \multicolumn{4}{|r|}{{Continued on next page}} \\ \hline
\endfoot
\hline
\endlastfoot

Harbour / escort tug & Azimuth Z-drive or VSP (main propulsion + steering) &
Escort/indirect-towing steering \& structural loads; uprated slewing duty;
steering redundancy; stability interaction; full steering-gear package. &
NR467 Pt D (tug/escort tug); Sec.14 \\
\hline

Offshore supply / PSV / AHTS & Bow tunnel(s) + retractable + stern azimuth; DP &
DYNAPOS/DP (usually DP-2): FMEA, redundancy, power segregation, capability plot;
FiFi where fitted; station-keeping. & NR467 Pt F (DYNAPOS); Pt E \\
\hline

Dredger (TSHD / CSD) & Bow tunnels; azimuth; sometimes DP &
DREDGING service; DP for some; robust manoeuvring; abrasion/duty considerations.
& NR467 Pt D \\
\hline

Passenger ship / RoPax / ferry & Bow (+ stern) tunnels; azimuth/pods on some &
SOLAS steering (II-1/29); \textbf{Safe Return to Port} survivability \&
redundancy for ships $\geq$120\,m / 3 fire zones; COMF. & SOLAS II-1/II-2; Pt E \\
\hline

Cruise ship & Podded propulsors + multiple bow/stern tunnels &
SRtP; COMF-NOISE/VIB; propulsion/steering redundancy; sometimes DP. &
SOLAS II-2 Reg.21/22; Pt E \\
\hline

Double-ended ferry & Azimuth or VSP at both ends &
Steering redundancy \emph{at each end}; control-station changeover; response
time both directions. & Sec.14; SOLAS II-1/29 \\
\hline

Cable / pipe layer, offshore construction & Retractable + azimuth + tunnels; DP &
DYNAPOS DP-2/DP-3: FMEA, fire/flood separation (DP-3), redundancy, proving
trials. & NR467 Pt F (DYNAPOS) \\
\hline

Drillship / semi-sub (MODU) & Multiple azimuth thrusters; DP-3 &
DP-3 station keeping; offshore-unit rules; fire/flood zone separation;
extensive FMEA. & NR467 Pt F (DYNAPOS); NR445 \\
\hline

Oil / chemical tanker & Bow thruster (manoeuvring) &
\textbf{Hazardous-area / explosion protection} where equipment is in or adjacent
to dangerous zones; thruster-room zoning; certified-safe (Ex) equipment. &
SOLAS II-2; IEC 60079; Pt C \\
\hline

Gas carrier (LNG/LPG) & Bow thruster; sometimes azimuth &
As tanker \emph{plus} IGC Code hazardous-zone requirements; gas-safe equipment. &
IGC Code; IEC 60079 \\
\hline

Ice-class / polar trader & Ice-strengthened tunnel/azimuth &
Ice-load strengthening \& winterisation (see Chapter \ref{ch:ice}). &
NR467 Pt E; NR527 \\
\hline

Research / naval / patrol & Rim-driven / pods / azimuth &
Low underwater-noise \& signature; COMF/silent; naval ships follow dedicated
naval rules. & Naval rules (e.g.\ NR483); Pt E \\
\hline

Large yacht & Bow/stern tunnels; pods &
Comfort (noise/vibration); yacht-specific rules; retractable/stabiliser
interfaces. & Yacht rules; Pt E \\
\hline

Container / bulk / general cargo & Bow tunnel thruster (manoeuvring) &
Base Sec.\,15 package; typically electric drive; no special service overlay. &
NR467 Sec.15; Pt C \\
\hline

\end{longtable}

\section{Passenger Ships \& Safe Return to Port (SRtP)}

For passenger ships to which SOLAS SRtP applies (broadly, ships $\geq$120\,m in
length or having three or more main vertical fire zones), the ship must retain,
after a defined casualty (fire in one zone or flooding of one watertight
compartment), a set of \textbf{essential systems --- including propulsion and
steering}. Where thrusters/pods/steering thrusters form part of the propulsion
or steering means, the surveyor must confirm:
\begin{itemize}[leftmargin=*]
\item redundancy and physical/functional separation so a single casualty does
not disable all propulsion \& steering;
\item power supply, control and cooling for the retained thruster(s) surviving
the casualty (fire zone / flooded compartment);
\item consistency with the ship's SRtP systems-availability documentation and
casualty-threshold analysis.
\end{itemize}
This typically \emph{adds} redundancy/segregation drawings and a systems
availability analysis to the base thruster package. \textit{(Ref.\ SOLAS II-2
Reg.21/22; NR467 passenger-ship requirements.)}

\section{Tankers \& Gas Carriers --- Hazardous Areas \& Explosion Protection}

On oil tankers, chemical tankers and gas carriers, the bow-thruster room and any
electrical/rotating equipment located in or adjacent to \textbf{hazardous
(gas-dangerous) zones} must meet explosion-protection requirements. The surveyor
should verify:
\begin{itemize}[leftmargin=*]
\item hazardous-area classification of the thruster space and boundaries;
\item \textbf{certified safe-type (Ex) equipment} for motors, sensors and
junction boxes where required, with certificates;
\item cable entries, penetrations and ventilation preserving zone integrity;
\item for gas carriers, compliance with the \textbf{IGC Code} zoning in addition
to SOLAS II-2.
\end{itemize}
\textit{(Ref.\ SOLAS II-2; IEC 60079 series; IGC Code; NR467 Pt C electrical in
hazardous areas.)} A tanker bow thruster is often electric; confirm the motor's
location relative to the cargo/hazardous zone \emph{before} accepting standard
(non-Ex) equipment.

\section{Offshore \& DP-Intensive Vessels}

OSVs, construction, cable/pipe-lay units and MODUs are defined by their DP
capability. Beyond the DP package of Chapter \ref{ch:dp}, expect: multiple
independent thruster trains, power-plant segregation, DP-3 fire/flood separation
where notated, worst-case-failure design intent, and comprehensive FMEA and
proving-trials programmes. Mobile offshore units are additionally governed by
\textbf{offshore-unit / MODU rules}, which may modify hull, machinery and
location (hazardous-area) requirements relative to NR467.

\section{Tugs \& Escort Tugs}

For tugs the thruster \emph{is} the main propulsion and steering means (Z-drive
or VSP). Escort tugs generate very high \textbf{indirect steering and towing
loads}; the surveyor must confirm the azimuth/steering system, unit structure
and hull foundation are substantiated for these loads (not merely free-running
thrust), together with the full steering-gear package (Chapter \ref{ch:steering})
and stability interaction. \textit{(Ref.\ NR467 Pt D tug / escort-tug service
notations.)}

\section{Thruster Compartment (All Ship Types)}

Independent of ship type, the thruster room / bow-thruster compartment carries
its own requirements the surveyor should not overlook: watertight boundaries and
subdivision, bilge/drainage, fire protection and boundary insulation, ventilation
and (for tankers/gas carriers) hazardous-area integrity, plus safe access and
escape. Confirm these are reflected in the general arrangement and hull/piping
documentation. \textit{(Ref.\ NR467 Part B, Ch 5; Part C fire \& piping;
SOLAS II-2.)}

\noindent\textbf{Surveyor note (ship type):} Read the \emph{requested notation
string} first --- it tells you the ship type, service and additional notations,
and therefore which overlays of Chapters
\ref{ch:primemover}--\ref{ch:notations} and which \emph{separate rule sets}
(SRtP, hazardous-area/IGC, offshore/MODU, naval, HSC) apply. The most common
ship-type gaps are: SRtP redundancy on passenger ships, Ex/hazardous-area
equipment on tankers and gas carriers, and DP-3 separation on offshore units.

\newpage

% Chapter 9: Testing
\chapter{Testing \& Commissioning Requirements}
\label{ch:testing}

\section{Factory Acceptance Tests (FAT)}

\textbf{Required for All Thrusters:}
\begin{itemize}[leftmargin=*]
\item Dimensional verification
\item Material verification
\item Non-destructive testing
\item Assembly inspection
\item Mechanical running test
\item Hydraulic system test (if applicable)
\item Electric motor test (if applicable)
\item Control system test
\item Retraction mechanism test (retractable)
\item CP mechanism test (CP propellers)
\item Azimuth slewing test (azimuth thrusters)
\end{itemize}

\section{Installation Tests}

\textbf{On-Site Verification:}
\begin{itemize}[leftmargin=*]
\item Foundation inspection
\item Alignment verification
\item Piping inspection
\item Electrical installation inspection
\item Control system integration test
\item Hydraulic system commissioning
\end{itemize}

\section{Dock Trials}

\textbf{Pre-Sea Trials:}
\begin{itemize}[leftmargin=*]
\item Static thrust test
\item No-load running test
\item Control system functional test
\item Retraction/extension test (retractable)
\item Azimuth rotation test (azimuth)
\item CP pitch control test (CP propellers)
\end{itemize}

\section{Sea Trials}

\textbf{Performance Verification:}
\begin{itemize}[leftmargin=*]
\item Performance verification
\item Maneuverability trials
\item Vibration measurements
\item Noise measurements
\item Thermal performance
\item Steering gear trials (if steering function)
\item DP trials (for DP systems)
\item DP FMEA proving trials (for DP systems)
\end{itemize}

\section{Documentation \& Reporting}

\textbf{Required Reports (\fic{}):}
\begin{itemize}[leftmargin=*]
\item FAT Report
\item Installation Report
\item Sea Trial Report
\item DP Trials Report (for DP systems)
\item As-Built Documentation
\end{itemize}

\newpage

% Chapter 10: Surveyor Checklists
\chapter{Surveyor Checklists by Thruster Type}
\label{ch:checklists}

\section{Tunnel Thrusters (Fixed) - Complete Checklist}

\subsection{General Information}
\begin{itemize}[label=$\square$]
\item Thruster power rating >110 kW (Section 15 applies)
\item Thruster manufacturer and model
\item Installation location (bow/stern)
\end{itemize}

\subsection{Structural}
\begin{itemize}[label=$\square$]
\item General Arrangement Drawing (\fa{})
\item Tunnel Structural Arrangement (\fa{})
\item Tunnel Structural Calculations (\fa{})
\item Tunnel End Grating/Grid Design (\fa{})
\item Hull Penetration Details (\fa{})
\end{itemize}

\subsection{Thruster Unit - Mechanical}
\begin{itemize}[label=$\square$]
\item Thruster Unit General Assembly (\fa{})
\item Propeller Design Drawing (\fa{})
\item Propeller Strength Calculations (\fa{})
\item Propeller Shaft Drawing (\fa{})
\item Shaft Strength Calculations (\fa{})
\item Bearing Arrangement Drawing (\fa{})
\item Bearing Load Calculations (\fa{})
\item Gearbox Assembly Drawing (\fa{})
\item Gear Tooth Calculations (\fa{})
\item Gearbox Lubrication System (\fa{})
\end{itemize}

\subsection{Drive System}
\textbf{If Electric:}
\begin{itemize}[label=$\square$]
\item Electric Motor Specification (\fa{})
\item Motor Foundation Drawing (\fa{})
\item Motor-Gearbox Coupling (\fa{})
\end{itemize}

\textbf{If Hydraulic:}
\begin{itemize}[label=$\square$]
\item Hydraulic Motor Specification (\fa{})
\item Hydraulic System Schematic (\fa{})
\item See Chapter \ref{ch:hydraulics} for complete hydraulic requirements
\end{itemize}

\subsection{Control System}
\begin{itemize}[label=$\square$]
\item Control System Block Diagram (\fa{})
\item Control Panel Layout (\fa{})
\item Electrical Single-Line Diagram (\fa{})
\item Control Software Description (\fa{})
\end{itemize}

\subsection{Calculations}
\begin{itemize}[label=$\square$]
\item Thrust Calculation (\fa{})
\item Vibration Analysis (\fa{})
\item Noise Level Prediction (\fa{})
\item Thermal Analysis (\fa{})
\end{itemize}

\subsection{Materials \& Welding}
\begin{itemize}[label=$\square$]
\item Material Specifications (\fa{})
\item Welding Procedures (\fa{})
\item Material Certificates (during construction)
\item NDT Plan (\fa{})
\end{itemize}

\subsection{Testing}
\begin{itemize}[label=$\square$]
\item FAT Procedures (\fa{})
\item Installation Test Procedures (\fa{})
\item Sea Trial Procedures (\fa{})
\item Test Reports (\fic{} - after testing)
\end{itemize}

\subsection{Manuals}
\begin{itemize}[label=$\square$]
\item Operation \& Maintenance Manual (\fic{})
\item Spare Parts List (\fic{})
\item Installation Manual (\fic{})
\end{itemize}

\section{Retractable Thrusters - Complete Checklist}

\textbf{All items from Section 10.1 (Tunnel Thrusters), PLUS:}

\subsection{Retraction Mechanism}
\begin{itemize}[label=$\square$]
\item Retraction Mechanism Assembly (\fa{})
\item Hydraulic Cylinder Specifications (\fa{})
\item Retraction System Calculations (\fa{})
\item Hull Opening \& Closure Design (\fa{})
\item Hull Opening Structural Analysis (\fa{})
\item Watertight Integrity Documentation (\fa{})
\item Locking System Design (\fa{})
\item Position Indication System (\fa{})
\item Emergency Lowering System (\fa{})
\item Hydraulic System for Retraction (\fa{})
\end{itemize}

\subsection{Additional Manuals}
\begin{itemize}[label=$\square$]
\item Retraction/Extension Procedures (\fic{})
\item Emergency Recovery Procedures (\fic{})
\end{itemize}

\subsection{Additional Testing}
\begin{itemize}[label=$\square$]
\item Retraction mechanism FAT
\item Seal integrity testing
\item Emergency lowering test
\item Dock trials - retraction/extension
\item Sea trials - retraction under way (if applicable)
\end{itemize}

\section{Azimuth Thrusters (Fixed Pitch) - Complete Checklist}

\subsection{General \& Structure}
\begin{itemize}[label=$\square$]
\item General Arrangement Drawing incl.\ 360° clearance envelope (\fa{})
\item Azimuth Unit Complete Assembly (\fa{})
\item Hull Penetration \& Foundation (\fa{})
\item Hull Structure Calculations / FEA (\fa{})
\end{itemize}

\subsection{Propeller \& Shaft}
\begin{itemize}[label=$\square$]
\item Propeller Design Drawing (\fa{})
\item Propeller Strength Calculations (\fa{})
\item Vertical Shaft Assembly \& Calculations (\fa{})
\item Thrust Bearing Design \& Calculations (\fa{})
\item Radial Bearing Arrangement (\fa{})
\item Column Structure Design (\fa{})
\item Seal System Design (\fa{})
\end{itemize}

\subsection{Gear \& Slewing System}
\begin{itemize}[label=$\square$]
\item Gear Train Assembly (\fa{})
\item Gear Tooth Strength Calculations --- ISO 6336 (\fa{})
\item Gear Lubrication System (\fa{})
\item Slewing Bearing Design \& Calculations (\fa{})
\item Azimuth Drive Mechanism \& Calculations (\fa{})
\item Position Feedback System (\fa{})
\end{itemize}

\subsection{Drive System \& Prime Mover}
\begin{itemize}[label=$\square$]
\item Electric or hydraulic power-transmission package (\fa{})
\item Slip ring / rotary joint (electric) (\fa{})
\item \textbf{Prime-mover add-on} --- electric or diesel (Chapter \ref{ch:primemover})
\end{itemize}

\subsection{Control, Calculations, Materials \& Testing}
\begin{itemize}[label=$\square$]
\item Control System Architecture \& Bridge Console (\fa{})
\item Local Control Station (\fa{})
\item Thrust / vibration / noise / thermal calculations (\fa{})
\item Material Specifications, Welding, NDT Plan (\fa{})
\item FAT / dock trial (azimuth slewing) / sea-trial procedures (\fa{})
\item O\&M and Installation Manuals (\fic{})
\end{itemize}

\subsection{If Steering / DP / Ice / Notations}
\begin{itemize}[label=$\square$]
\item If steering function $\rightarrow$ Chapter \ref{ch:steering} package (SOLAS II-1/29, FMEA)
\item If part of main shaft line $\rightarrow$ Chapter \ref{ch:shafting} package
\item If DP $\rightarrow$ Section \ref{ch:checklists}.\ (DP add-on below)
\item If ice class $\rightarrow$ Chapter \ref{ch:ice} add-on
\item Overlay applicable notations $\rightarrow$ Chapter \ref{ch:notations}
\end{itemize}

\section{Azimuth Thrusters (Controllable Pitch) - Complete Checklist}

\textbf{All items from Azimuth FP, PLUS:}
\begin{itemize}[label=$\square$]
\item CP Hub Mechanism Assembly (\fa{})
\item Blade Retention Calculations (\fa{})
\item Pitch Control Mechanism (\fa{})
\item Pitch Control Hydraulic System incl.\ rotary joint (\fa{}) --- Chapter \ref{ch:hydraulics}
\item Pitch feedback \& fail-safe position (\fa{})
\end{itemize}

\section{Podded Propulsion Units - Complete Checklist}

\textbf{All items from Azimuth (FP or CP), PLUS:}
\begin{itemize}[label=$\square$]
\item Pod-specific structural components \& strut (\fa{})
\item Electric motor in pod --- specification, cooling, monitoring (\fa{})
\item Power supply through strut / slip-ring (\fa{})
\item Propeller direct-drive arrangement (\fa{})
\item Pod-specific FAT/testing (\fa{})
\item \textbf{Electric prime-mover add-on} (Chapter \ref{ch:primemover})
\end{itemize}

\section{Voith-Schneider (Cycloidal) Units - Complete Checklist}
\begin{itemize}[label=$\square$]
\item VSP Unit General Assembly (\fa{})
\item Blade (wing) design + blade-shaft/blade strength calculations (\fa{})
\item Blade bearing arrangement (\fa{})
\item Kinematic control gear / steering mechanism strength (\fa{})
\item Rotor casing \& main bevel gear + gear calculations (\fa{})
\item Main thrust bearing design \& life (\fa{})
\item Bottom cover / guard plate structure (\fa{})
\item Hull seat, foundation \& well structure / FEA (\fa{})
\item Seal system (rotor \& blade shafts) (\fa{})
\item Servo / hydraulic control unit (\fa{}) --- Chapter \ref{ch:hydraulics}
\item Steering control system \& bridge console (\fa{})
\item \textbf{Steering-function package} --- Chapter \ref{ch:steering} (FMEA, redundancy, SOLAS II-1/29)
\item Prime-mover add-on (Chapter \ref{ch:primemover}) + drive-train TVC
\end{itemize}

\section{Water-Jet Systems - Complete Checklist}
\begin{itemize}[label=$\square$]
\item Water-jet general arrangement \& installation (\fa{})
\item Impeller/rotor design \& strength calculations (\fa{})
\item Shaft line \& bearings (\fa{}) --- Chapter \ref{ch:shafting}
\item Inlet duct \& transom structure / hull penetration (\fa{})
\item Steering \& reversing (bucket) mechanism (\fa{})
\item Steering/reversing hydraulic system (\fa{}) --- Chapter \ref{ch:hydraulics}
\item Control system \& bridge interface (\fa{})
\item Drive package (direct-from-engine or via gearbox) + TVC (\fa{})
\item Performance / thrust calculations (\fa{})
\item Materials, welding, NDT (\fa{})
\item FAT / trials procedures (\fa{}); O\&M manuals (\fic{})
\item If steering means of ship $\rightarrow$ Chapter \ref{ch:steering}; if DP $\rightarrow$ DP add-on
\end{itemize}

\section{Prime-Mover Add-On Checklist (apply to any type)}

\textbf{If ELECTRIC-motor driven:}
\begin{itemize}[label=$\square$]
\item Drive-motor data sheet (\fa{})
\item Frequency converter (VFD) specification \& arrangement (\fa{})
\item Harmonic / power-quality analysis (\fa{})
\item Supply transformer \& switchboard interface (\fa{})
\item Power cable sizing, routing \& segregation (\fa{})
\item Motor protection, monitoring \& cooling (\fa{})
\end{itemize}

\textbf{If DIESEL-engine driven:}
\begin{itemize}[label=$\square$]
\item Diesel engine approval / type data (\fa{})
\item \textbf{Torsional vibration calculation of the complete train} (\fa{})
\item Clutch \& flexible coupling (\fa{})
\item Engine foundation \& chocking (\fa{})
\item Fuel oil / lube oil / cooling water systems (\fa{})
\item Starting system \& exhaust gas system (\fa{})
\item Engine control, safety \& monitoring (\fa{})
\item NOx Technical File / EIAPP (\fic{})
\end{itemize}

\section{Ice Class / Winterisation Add-On Checklist (apply if notation carried)}
\begin{itemize}[label=$\square$]
\item Ice-class notation basis \& design conditions (\fa{})
\item Propeller ice-load strength calculation (\fa{})
\item Shaft line \& ice-torque verification (\fa{})
\item Gear \& azimuth/nozzle ice-load substantiation (\fa{})
\item Thruster lower-unit ice protection (\fa{})
\item Tunnel/inlet grid ice arrangement (\fa{})
\item Winterisation of hydraulic/cooling/seal systems (\fa{})
\item Low-temperature material certification (\fa{})
\end{itemize}

\section{DP Thrusters - Additional Checklist Items}

\textbf{For ANY thruster type used in DP system, add:}
\begin{itemize}[label=$\square$]
\item DP System General Arrangement (\fa{})
\item Thruster Configuration Drawing (\fa{})
\item DP Capability Plot (\fa{})
\item FMEA Documentation (\fa{})
\item Redundancy Concept (\fa{})
\item Power Supply Segregation (\fa{})
\item DP Control System Architecture (\fa{})
\item Position Reference System Configuration (\fa{})
\item PMS Configuration (\fa{})
\item UPS Configuration (\fa{})
\item DP FMEA Proving Trials Plan (\fa{})
\item DP Trials Reports (\fic{})
\item DP Operations Manual (\fic{})
\end{itemize}

\newpage

% Chapter: Submission Cover Sheets by Ship Type
\chapter{Submission Cover Sheets by Ship Type}
\label{ch:coversheets}

\textbf{How to use:} one cover sheet per page. Pick the sheet matching the
vessel's service notation, complete the identification block, then tick each
line as the corresponding drawing/document is received and found acceptable.
Every sheet shares a \emph{common core} (base package, type-specific package,
prime mover, testing, manuals) and then lists the \textbf{ship-type-specific
extras} --- the items most often missing for that service. Use together with the
detailed checklists of Chapter \ref{ch:checklists} and the overlays of
Chapters \ref{ch:primemover}--\ref{ch:shiptype}.
Legend: \fa{} = For Approval, \fic{} = For Information.

\CoverSheet{Generic Template (any ship / any thruster type)}
\textbf{Core package (all thrusters $>$110\,kW):}
\begin{itemize}[label=$\square$,nosep]
\CoreCoverItems
\end{itemize}
\textbf{Functional overlays --- tick those applicable:}
\begin{itemize}[label=$\square$,nosep]
\item Steering function $\rightarrow$ Chapter \ref{ch:steering} (FMEA, redundancy, SOLAS II-1/29)
\item Part of main shaft line $\rightarrow$ Chapter \ref{ch:shafting} (+ MON-SHAFT, Section \ref{sec:monshaft}, if notated)
\item Dynamic positioning $\rightarrow$ Chapter \ref{ch:dp} (FMEA, capability plot, proving trials)
\item Hydraulic actuation $\rightarrow$ Chapter \ref{ch:hydraulics}
\end{itemize}
\textbf{Additional-notation \& ice overlays --- tick those in the notation string:}
\begin{itemize}[label=$\square$,nosep]
\item Ice class / Polar / COLD $\rightarrow$ Chapter \ref{ch:ice}
\item DP / COMF / AUT / CLEANSHIP / RP / MON-SHAFT $\rightarrow$ Chapter \ref{ch:notations}
\end{itemize}

\CoverSheet{Harbour / Escort Tug}
\begin{itemize}[label=$\square$,nosep]
\CoreCoverItems
\end{itemize}
\textbf{Ship-type-specific extras (tug):}
\begin{itemize}[label=$\square$,nosep]
\item Azimuth (Z-drive) or Voith-Schneider unit --- full type package (\fa{})
\item \textbf{Full steering-gear package} --- Chapter \ref{ch:steering}: FMEA, redundancy, control stations, response time, SOLAS II-1/29 (\fa{})
\item \textbf{Escort / indirect-towing loads} on steering \& structure; unit \& foundation substantiation for these loads, not free-running only (\fa{})
\item Uprated slewing duty \& brake/locking; stability interaction documentation (\fa{})
\item Drive-train torsional vibration calculation (diesel Z-drive) (\fa{})
\end{itemize}

\CoverSheet{Offshore Supply / PSV / AHTS (DP)}
\begin{itemize}[label=$\square$,nosep]
\CoreCoverItems
\end{itemize}
\textbf{Ship-type-specific extras (OSV / DP):}
\begin{itemize}[label=$\square$,nosep]
\item Bow tunnel(s) + retractable + stern azimuth --- type packages (\fa{})
\item \textbf{DYNAPOS / DP (usually DP-2)} --- Chapter \ref{ch:dp}: FMEA, redundancy concept, power segregation, capability plot, thruster allocation, proving-trials plan (\fa{})
\item Retractable/take-me-home thruster package incl.\ emergency lowering (\fa{})
\item FiFi and other service systems interfaces where fitted (\fic{})
\item Steering-function compliance for azimuth units (\fa{})
\end{itemize}

\CoverSheet{Passenger Ship / RoPax / Ferry}
\begin{itemize}[label=$\square$,nosep]
\CoreCoverItems
\end{itemize}
\textbf{Ship-type-specific extras (passenger):}
\begin{itemize}[label=$\square$,nosep]
\item Bow (+ stern) tunnel thrusters / azimuth --- type packages (\fa{})
\item \textbf{Safe Return to Port} redundancy \& segregation of propulsion/steering (ships $\geq$120\,m or 3 fire zones); systems-availability analysis (\fa{})
\item SOLAS steering-gear compliance (II-1/29) where thrusters steer (\fa{})
\item COMF-NOISE / COMF-VIB noise \& vibration prediction (\fa{})
\item Emergency/standby power to retained thruster after casualty (\fa{})
\end{itemize}

\CoverSheet{Cruise Ship}
\begin{itemize}[label=$\square$,nosep]
\CoreCoverItems
\end{itemize}
\textbf{Ship-type-specific extras (cruise):}
\begin{itemize}[label=$\square$,nosep]
\item Podded propulsors + multiple bow/stern tunnels --- type packages (\fa{})
\item \textbf{Safe Return to Port} redundancy \& systems availability (\fa{})
\item \textbf{COMF-NOISE / COMF-VIB} --- blade-rate excitation, cavitation noise, resilient mounting (\fa{})
\item Pod electrical-machine, cooling \& monitoring; slip-ring/power-through-strut (\fa{})
\item DP package if DP-notated (\fa{})
\end{itemize}

\CoverSheet{Double-Ended Ferry}
\begin{itemize}[label=$\square$,nosep]
\CoreCoverItems
\end{itemize}
\textbf{Ship-type-specific extras (double-ended):}
\begin{itemize}[label=$\square$,nosep]
\item Azimuth or Voith-Schneider unit \emph{at each end} --- type packages (\fa{})
\item \textbf{Steering redundancy at both ends} + control-station changeover logic (\fa{})
\item Response time / manoeuvring in both running directions (\fa{})
\item COMF noise \& vibration (passenger comfort) (\fa{})
\end{itemize}

\CoverSheet{Dredger (TSHD / CSD)}
\begin{itemize}[label=$\square$,nosep]
\CoreCoverItems
\end{itemize}
\textbf{Ship-type-specific extras (dredger):}
\begin{itemize}[label=$\square$,nosep]
\item Bow tunnels / azimuth --- type packages (\fa{})
\item DREDGING service duty; robust manoeuvring provisions (\fa{})
\item DP package where station-keeping/tracking notated (\fa{})
\item Duty-cycle / abrasion considerations in bearing \& seal life (\fic{})
\end{itemize}

\CoverSheet{Cable / Pipe Layer \& Offshore Construction (DP)}
\begin{itemize}[label=$\square$,nosep]
\CoreCoverItems
\end{itemize}
\textbf{Ship-type-specific extras (construction / DP):}
\begin{itemize}[label=$\square$,nosep]
\item Retractable + azimuth + tunnel thrusters --- type packages (\fa{})
\item \textbf{DYNAPOS DP-2 / DP-3} --- Chapter \ref{ch:dp}: FMEA, redundancy, power segregation, capability plot, proving trials (\fa{})
\item DP-3 \textbf{fire / flood zone separation} of thruster groups \& power (\fa{})
\item Worst-case-failure design intent documentation (\fa{})
\end{itemize}

\CoverSheet{Drillship / Semi-Submersible (MODU)}
\begin{itemize}[label=$\square$,nosep]
\CoreCoverItems
\end{itemize}
\textbf{Ship-type-specific extras (MODU / DP-3):}
\begin{itemize}[label=$\square$,nosep]
\item Multiple azimuth thrusters --- type packages (\fa{})
\item \textbf{DYNAPOS DP-3} station keeping --- full FMEA, redundancy, proving trials (\fa{})
\item \textbf{Fire \& flood separation} of thruster/power groups (DP-3) (\fa{})
\item Compliance with \textbf{offshore-unit / MODU rules} (hull, machinery, hazardous-area) (\fa{})
\end{itemize}

\CoverSheet{Oil / Chemical Tanker}
\begin{itemize}[label=$\square$,nosep]
\CoreCoverItems
\end{itemize}
\textbf{Ship-type-specific extras (tanker):}
\begin{itemize}[label=$\square$,nosep]
\item Bow thruster (usually electric) --- type package (\fa{})
\item \textbf{Hazardous-area classification} of thruster room \& boundaries (\fa{})
\item \textbf{Certified safe-type (Ex) equipment} for motor/sensors/junction boxes where in or adjacent to dangerous zones, with certificates (\fa{})
\item Cable entries, penetrations \& ventilation preserving zone integrity (\fa{})
\item Confirm motor location vs.\ hazardous zone \emph{before} accepting non-Ex equipment (\fa{})
\end{itemize}

\CoverSheet{Gas Carrier (LNG / LPG)}
\begin{itemize}[label=$\square$,nosep]
\CoreCoverItems
\end{itemize}
\textbf{Ship-type-specific extras (gas carrier):}
\begin{itemize}[label=$\square$,nosep]
\item Bow thruster / azimuth --- type package (\fa{})
\item \textbf{IGC Code} hazardous-zone requirements in addition to SOLAS II-2 (\fa{})
\item Gas-safe / certified (Ex) electrical equipment \& certificates (\fa{})
\item Hazardous-area drawings for thruster space \& boundaries (\fa{})
\end{itemize}

\CoverSheet{Ice-Class / Polar Trader}
\begin{itemize}[label=$\square$,nosep]
\CoreCoverItems
\end{itemize}
\textbf{Ship-type-specific extras (ice) --- Chapter \ref{ch:ice}:}
\begin{itemize}[label=$\square$,nosep]
\item Ice-class notation basis \& design conditions (\fa{})
\item \textbf{Propeller ice-load strength} calculation (\fa{})
\item \textbf{Shaft line \& gear ice-torque} verification (blade-milling / blade-loss) (\fa{})
\item Azimuth/pod lower-unit \& slewing gear vs.\ ice-block impact (\fa{})
\item Winterisation (low-temp oils, materials, sensor heating; anti-icing) (\fa{})
\item Low-temperature material certification (\fa{})
\end{itemize}

\CoverSheet{Research / Naval / Patrol}
\begin{itemize}[label=$\square$,nosep]
\CoreCoverItems
\end{itemize}
\textbf{Ship-type-specific extras (research / naval):}
\begin{itemize}[label=$\square$,nosep]
\item Rim-driven / podded / azimuth units --- type packages (\fa{})
\item \textbf{Low underwater-noise / signature} \& COMF (silent) documentation (\fa{})
\item Cavitation-noise control measures (\fa{})
\item Naval vessels: compliance with applicable \textbf{naval ship rules} (\fa{})
\end{itemize}

\CoverSheet{Large Yacht}
\begin{itemize}[label=$\square$,nosep]
\CoreCoverItems
\end{itemize}
\textbf{Ship-type-specific extras (yacht):}
\begin{itemize}[label=$\square$,nosep]
\item Bow/stern tunnel thrusters and/or pods --- type packages (\fa{})
\item \textbf{Comfort} (noise \& vibration) documentation (\fa{})
\item Retractable-thruster \& stabiliser interface where fitted (\fa{})
\item Applicable yacht-specific rules (\fa{})
\end{itemize}

\CoverSheet{Container / Bulk / General Cargo}
\begin{itemize}[label=$\square$,nosep]
\CoreCoverItems
\end{itemize}
\textbf{Ship-type-specific extras (cargo):}
\begin{itemize}[label=$\square$,nosep]
\item Bow tunnel thruster for manoeuvring --- type package (\fa{})
\item Typically electric drive --- confirm motor, cabling \& protection (\fa{})
\item No special service overlay unless ice class / other notation carried (verify string) (\fa{})
\end{itemize}

\newpage

% Chapter 11: Common Errors
\chapter{Common Submission Errors \& Surveyor Notes}
\label{ch:errors}

\section{Frequent Documentation Gaps}

\subsection{Missing or Incomplete Calculations}
\begin{itemize}[leftmargin=*]
\item Shaft strength calculations missing fatigue analysis
\item Bearing life calculations without proper load factors
\item Gear calculations not per recognized standards (ISO 6336)
\item Vibration analysis incomplete or missing
\item Hydraulic system calculations (accumulator sizing, pressure drop)
\end{itemize}

\subsection{Inadequate Drawings}
\begin{itemize}[leftmargin=*]
\item General arrangements without sufficient detail
\item Material specifications not clearly indicated
\item Welding details missing or incomplete
\item Foundation drawings without load specifications
\item Control system diagrams without protection details
\end{itemize}

\subsection{Steering Gear Overlap Not Addressed}
\begin{itemize}[leftmargin=*]
\item Azimuth or podded thrusters used for steering without steering gear compliance documentation
\item Missing FMEA for steering function
\item Inadequate redundancy analysis
\item Missing emergency steering provisions
\item SOLAS compliance not demonstrated
\end{itemize}

\subsection{DP System Gaps}
\begin{itemize}[leftmargin=*]
\item FMEA not covering all failure modes or worst-case scenarios
\item Inadequate power segregation for DP-2/DP-3
\item Missing fire/flood zone analysis for DP-3
\item Thruster interaction effects not analyzed
\item Incomplete proving trials procedures
\end{itemize}

\subsection{Hydraulic System Issues}
\begin{itemize}[leftmargin=*]
\item Undersized accumulators
\item Inadequate pressure relief provisions
\item Missing filtration specifications
\item Incomplete piping stress analysis
\item Emergency provisions not clearly defined
\end{itemize}

\subsection{Prime-Mover Gaps}
\begin{itemize}[leftmargin=*]
\item Diesel-driven Z-drive submitted \emph{without} a torsional vibration calculation for the complete engine--clutch--shaft--thruster train (most common tug/OSV gap)
\item Missing barred-speed range or clutch-in/misfiring conditions in the TVC
\item Electric-driven thruster without harmonic / power-quality study or cable-sizing \& segregation
\item Hybrid PTI/PTO arrangements with undocumented driving paths or mode-changeover logic
\end{itemize}

\subsection{Ice-Class \& Notation Gaps}
\begin{itemize}[leftmargin=*]
\item Ice-classed vessel with propeller/shaft/gear checked only for open-water rating, not ice peak-torque / blade-milling loads
\item Azimuth/pod lower unit and slewing gear not substantiated against ice-block impact
\item Winterisation (low-temperature oils, materials, sensor heating) not addressed
\item Comfort (COMF) notation carried but no thruster noise/vibration prediction
\item Requested notation string not reconciled with the submitted thruster package
\end{itemize}

\subsection{Cycloidal (VSP) \& Combined Propulsion/Steering Gaps}
\begin{itemize}[leftmargin=*]
\item VSP mechanical package submitted without the steering-function compliance set (FMEA, redundancy, SOLAS II-1/29, response time)
\item Blade guard/bottom plate and its hull attachment not substantiated for grounding/impact
\item Steerable water-jet / nozzle / grid units missing steering-gear documentation
\end{itemize}

\section{Surveyor Review Tips}

\subsection{Initial Review}
\begin{enumerate}
\item Verify thruster power >110 kW → Section 15 applies
\item Identify thruster type → Apply appropriate checklist (Chapters \ref{ch:transverse}--\ref{ch:othertypes})
\item Identify the prime mover (electric / diesel / hydraulic) → Add Chapter \ref{ch:primemover} submissions
\item Check for steering function → Add Chapter \ref{ch:steering} requirements
\item Check for main shafting classification → Add Chapter \ref{ch:shafting} requirements
\item Check for DP application → Add Chapter \ref{ch:dp} requirements
\item Check ice class / winterisation notation → Add Chapter \ref{ch:ice} requirements
\item Reconcile the full class-notation string → Add Chapter \ref{ch:notations} requirements
\end{enumerate}

\subsection{Calculation Review}
\begin{itemize}[leftmargin=*]
\item Verify all calculations use BV-approved formulas or recognized standards
\item Check safety factors applied
\item Verify material properties used in calculations match specifications
\item Ensure all loading conditions considered (normal, extreme, emergency)
\end{itemize}

\subsection{Drawing Review}
\begin{itemize}[leftmargin=*]
\item Check for consistency between drawings
\item Verify dimensions and tolerances
\item Ensure material specifications clearly indicated
\item Check for proper welding symbols and details
\item Verify accessibility for maintenance
\end{itemize}

\section{Red Flags for Surveyors}

\subsection{Immediate Rejection Criteria}
\begin{itemize}[leftmargin=*]
\item Missing critical calculations (shaft strength, bearing life, gear strength)
\item Material specifications not compliant with BV rules
\item Welding procedures not approved
\item Control system without adequate safety provisions
\item Steering gear requirements not addressed when applicable
\item DP FMEA incomplete or inadequate
\end{itemize}

\subsection{Require Clarification}
\begin{itemize}[leftmargin=*]
\item Ambiguous material specifications
\item Incomplete or unclear drawings
\item Calculations without clear assumptions
\item Testing procedures without acceptance criteria
\item Manuals not in working language of crew
\end{itemize}

\newpage

% Chapter 12: Quick Reference
\chapter{Quick Reference Tables}
\label{ch:quickref}

\section{Thruster Type vs. Applicable Rule Sections}

\begin{table}[H]
\centering
\caption{Applicability Matrix}
\begin{tabular}{|l|c|c|c|c|c|}
\hline
\rowcolor{bvblue}
\textcolor{white}{\textbf{Thruster Type}} & \textcolor{white}{\textbf{Sec.15}} & \textcolor{white}{\textbf{Sec.14}} & \textcolor{white}{\textbf{Sec.7}} & \textcolor{white}{\textbf{DP (Pt F)}} & \textcolor{white}{\textbf{Hydraulics}} \\
\hline
Tunnel (Fixed) & ✓ & Optional & Optional & Optional & If hydraulic \\
\hline
Retractable & ✓ & Optional & Optional & Optional & ✓ Required \\
\hline
Azimuth FP & ✓ & Often ✓ & Typically ✓ & Optional & If hydraulic \\
\hline
Azimuth CP & ✓ & Often ✓ & Typically ✓ & Optional & ✓ Required \\
\hline
Podded & ✓ & Often ✓ & Typically ✓ & Optional & If CP \\
\hline
Water-Jet & ✓ & Often ✓ & ✓ Required & Optional & ✓ Required \\
\hline
Voith-Schneider & ✓ & ✓ Required & Typically ✓ & Optional & ✓ Required \\
\hline
Rim-Driven & ✓ & Optional & Optional & Optional & Rarely \\
\hline
Pump-Jet / Ducted & ✓ & If steering & ✓ Required & Optional & If steering \\
\hline
Gill / Hydro-Jet & ✓ & Optional & Optional & Optional & If hydraulic \\
\hline
Steering Grid / Nozzle & ✓ & ✓ Required & ✓ Required & Optional & ✓ Required \\
\hline
\end{tabular}

\vspace{0.2cm}
{\footnotesize Sec.15 = Thrusters; Sec.14 = Steering Gear; Sec.7 = Shaft
lines/propellers (NR467 Pt C, Ch 1). ``Often/Typically ✓'' means applicable when
the unit provides steering or forms part of the main shaft line. Overlay
Chapter \ref{ch:primemover} (prime mover), Chapter \ref{ch:ice} (ice class) and
Chapter \ref{ch:notations} (notations) on every row as applicable.}
\end{table}

\section{Prime Mover \& Additional-Notation Quick Guide}

\begin{table}[H]
\centering
\caption{Prime Mover --- Key Additional Submissions}
\begin{tabular}{|p{3.2cm}|p{11.5cm}|}
\hline
\rowcolor{bvblue}
\textcolor{white}{\textbf{Prime mover}} & \textcolor{white}{\textbf{Headline additional submissions (see Chapter \ref{ch:primemover})}} \\
\hline
Electric motor & Motor data sheet; frequency converter (VFD); harmonic/power-quality study; transformer \& switchboard interface; cable sizing \& segregation; motor protection/monitoring/cooling. \\
\hline
Diesel engine & Engine approval (Sec.2); \textbf{full-train torsional vibration calculation}; clutch \& flexible coupling; foundation/chocking; FO/LO/cooling/starting/exhaust systems; engine safety \& controls; NOx file. \\
\hline
Hydraulic motor & Per Chapter \ref{ch:hydraulics} (HPU, circuit, accumulators, relief, filtration, monitoring). \\
\hline
\end{tabular}
\end{table}

\begin{table}[H]
\centering
\caption{Additional Notations --- Trigger Reference}
\begin{tabular}{|p{4cm}|p{10.7cm}|}
\hline
\rowcolor{bvblue}
\textcolor{white}{\textbf{Notation carried}} & \textcolor{white}{\textbf{Adds}} \\
\hline
DP / DYNAPOS & Chapter \ref{ch:dp}: FMEA, redundancy, power segregation, capability plot, proving trials (NR467 Pt F, DYNAPOS). \\
\hline
Ice class / Polar / COLD & Chapter \ref{ch:ice}: ice-load propeller/shaft/gear strengthening, winterisation, low-temp materials (NR467 Pt E, NR527, IACS UR I3). \\
\hline
COMF-NOISE / COMF-VIB & Chapter \ref{ch:notations}: noise \& vibration prediction, blade-rate excitation, resilient mounting. \\
\hline
AUT-UMS / AUT-CCS & Chapter \ref{ch:notations}: extended alarms/monitoring, safe remote control, fail-safe. \\
\hline
CLEANSHIP & Chapter \ref{ch:notations}: EAL seal/hydraulic fluids, leakage monitoring. \\
\hline
MON-SHAFT & Section \ref{sec:monshaft}: oil-lubricated bearing/seal, lube-oil analysis \& weardown monitoring, condition-based shaft survey. \\
\hline
TUG / ESCORT TUG & Chapter \ref{ch:notations}: uprated steering \& structural loads. \\
\hline
\end{tabular}
\end{table}

\section{Document Approval Categories Summary}

\begin{table}[H]
\centering
\caption{Approval Categories}
\begin{tabular}{|l|p{4cm}|p{3cm}|p{4cm}|}
\hline
\rowcolor{bvblue}
\textcolor{white}{\textbf{Category}} & \textcolor{white}{\textbf{Meaning}} & \textcolor{white}{\textbf{Timing}} & \textcolor{white}{\textbf{Surveyor Action}} \\
\hline
\fa{} & Must be reviewed and approved before construction & Submit before construction & Detailed technical review \\
\hline
\fic{} & Submitted for record, not requiring formal approval & Submit as specified & Review for completeness \\
\hline
During Construction & Certificates/reports generated during manufacturing & Submit as work progresses & Verify during surveys \\
\hline
\end{tabular}
\end{table}

\newpage

% Chapter 13: Regulatory References
\chapter{Regulatory References}
\label{ch:regulations}

\section{Bureau Veritas Rules}

\subsection{NR467 - Rules for Classification of Steel Ships}

\textbf{Part A: Classification and Surveys} --- Chapter 1 (``Principles of
Classification and Class Notations'') is where every \textbf{service notation}
and \textbf{additional class notation} is formally \emph{listed and defined}
(the master index of ship types and notations), together with the principles
of assignment/maintenance of class (Ch.\,2) and survey scope (Ch.\,3--5). The
\emph{detailed technical requirements} for each notation then sit in Parts
D/E (service notations) and F (additional class notations) as set out below
--- Part A tells the surveyor \emph{which} notations exist and apply; Parts
D/E/F tell the surveyor \emph{what to demand} for each one.

\textbf{Part B: Hull and Stability}
\begin{itemize}[leftmargin=*]
\item Chapter 2: Subdivision and Stability
\item Chapter 3: Materials
\item Chapter 4: Welding
\item Chapter 5: Hull Structure
\end{itemize}

\textbf{Part C: Machinery, Electricity, Automation and Fire Protection}
\begin{itemize}[leftmargin=*]
\item Chapter 1 --- Machinery:
\begin{itemize}
\item Section 2: Diesel Engines (prime-mover-driven thrusters)
\item Sections 7--8: Shaft Lines and Propellers
\item Section 14: Steering Gear
\item Section 15: Thrusters (and water-jets)
\end{itemize}
\item Chapter 2 --- Electrical Installations (drive motors, converters, cables)
\item Chapter 3 --- Automation (control, alarms, AUT notations)
\item Chapter 4 --- Fire Protection
\end{itemize}
\textit{Note: exact section numbering varies with rule edition --- always confirm
against the current NR467 in Rules Explorer.}

\textbf{Part D: Service Notations} --- general ship-type service notations
(e.g.\ passenger ship, oil tanker, dredger, container ship). Determines the
ship-type requirements overlaid on the thruster package.

\textbf{Part E: Service Notations for Offshore Service Vessels and Tugs} ---
e.g.\ \emph{tug}, \emph{escort tug}, \emph{supply}, anchor-handling. Source of
the uprated escort/towing steering \& structural loads for tug/OSV thrusters.

\textbf{Part F: Additional Class Notations} --- the additional (voluntary)
notations, in multiple chapters:
\begin{itemize}[leftmargin=*]
\item Ice class (IA Super / IA / IB / IC) and Polar Class (PC1--PC7); COLD winterisation
\item Comfort on board --- COMF-NOISE, COMF-VIB
\item Automated machinery spaces / integrated systems --- AUT-UMS, AUT-CCS, AUT-PORT
\item Dynamic positioning --- \textbf{DYNAPOS} (AM / AT / ATR, SAM / SAT)
\item Environmental protection --- CLEANSHIP and related (EAL lubricants)
\item Redundant propulsion --- RP / AVM notations
\item Shaft-line monitoring --- MON-SHAFT
\end{itemize}

\vspace{0.2cm}
\noindent\fcolorbox{warningred}{lightgray}{\begin{minipage}{0.95\textwidth}
\small\textbf{\textcolor{warningred}{Reference correction:}} Dynamic positioning is
an \emph{additional class notation} (\textbf{DYNAPOS}) defined in \textbf{NR467
Part F}. There is no ``NR216 DP rule'': \textbf{NR216} is the \emph{Rules on
Materials and Welding} (see below). Earlier drafts of this catalog mis-cited
NR216 for DP --- read DP references as \textbf{NR467 Part F (DYNAPOS)}.
\end{minipage}}

\subsection{Related BV Rule Notes (as applicable)}
\begin{itemize}[leftmargin=*]
\item \textbf{NR216} --- Rules on Materials and Welding for the classification of marine units
\item \textbf{NR266} --- Requirements for Survey of Materials and Equipment
\item \textbf{NR320} --- Certification scheme of materials and equipment for marine units
\item \textbf{NR527} --- Ships operating in polar waters / icebreakers (Polar Class)
\item \textbf{NR500} --- Classification and certification of \emph{yachts}
\item \textbf{NR566} --- Hull arrangement, stability and systems for \emph{ships $<$ 500 GT}
\item \textbf{NR483} --- Classification of \emph{naval ships}
\item \textbf{NR445} --- Classification of \emph{offshore units} (MODU / drillships)
\end{itemize}

\section{International Standards}

\subsection{ISO Standards}
\begin{itemize}[leftmargin=*]
\item \textbf{ISO 6336}: Calculation of load capacity of spur and helical gears
\item \textbf{ISO 22547}: Ships and marine technology --- azimuth thrusters
\item \textbf{ISO 484}: Shipbuilding --- ship screw propellers (manufacturing tolerances)
\item \textbf{ISO 4413}: Hydraulic fluid power --- general rules for systems
\item \textbf{ISO 20283 / ISO 6954}: Mechanical vibration --- measurement \& evaluation (supports COMF-VIB)
\end{itemize}

\subsection{IEC Standards}
\begin{itemize}[leftmargin=*]
\item \textbf{IEC 60092 Series}: Electrical installations in ships
\begin{itemize}
\item IEC 60092-101: Definitions and general requirements
\item IEC 60092-201: System design --- general
\item IEC 60092-301: Equipment - Generators and motors
\item IEC 60092-352: Choice and installation of cables
\item IEC 60092-504: Control and instrumentation
\end{itemize}
\item \textbf{IEC 60034 Series}: Rotating electrical machines (ratings, cooling, temperature rise)
\item \textbf{IEC 61508}: Functional safety (for DP / thruster control systems)
\end{itemize}

\subsection{IACS Unified Requirements \& Codes}
\begin{itemize}[leftmargin=*]
\item \textbf{IACS UR I1--I3}: Polar Class --- ships and machinery requirements (ice loads on propeller/shaft/gear)
\item \textbf{IACS UR M}: Machinery requirements (shafting, gears, TVC principles)
\item \textbf{IACS UR E22 / E26 / E27}: Programmable / cyber-resilient control systems
\item \textbf{IMO Polar Code}: Ships operating in polar waters
\item \textbf{Finnish-Swedish Ice Class Rules (TRAFICOM)}: Baltic ice classes basis
\end{itemize}

\subsection{SOLAS Convention}
\begin{itemize}[leftmargin=*]
\item \textbf{Chapter II-1, Regulation 29}: Steering gear
\item \textbf{Chapter II-1, Reg.\ 8-1 \& II-2, Reg.\ 21/22}: Safe Return to Port / systems availability after a casualty (passenger ships)
\item \textbf{Chapter II-2}: Fire protection; hazardous-area boundaries
\item \textbf{Chapter V, Regulation 28}: Maneuverability
\end{itemize}

\subsection{Ship-Type / Service Rule Sets \& Codes}
\begin{itemize}[leftmargin=*]
\item \textbf{IEC 60079 series}: Explosion protection --- equipment in hazardous areas (tankers, gas carriers, bow-thruster rooms)
\item \textbf{IGC Code}: Gas carriers --- hazardous-zone requirements
\item \textbf{IMO MODU Code / offshore-unit rules}: Mobile offshore drilling units and offshore service craft
\item \textbf{HSC Code}: High-speed craft (may modify/replace NR467 machinery requirements)
\item \textbf{Naval ship rules} (e.g.\ BV NR483): Naval and patrol vessels --- signature/silent operation
\item \textbf{NR467 Part D}: Service notations (tug, escort tug, supply, dredging, passenger, tanker, etc.)
\end{itemize}

\subsection{IMO Resolutions}
\begin{itemize}[leftmargin=*]
\item \textbf{IMO Res.A.468(XII)}: Code on Noise Levels
\item \textbf{IMO Res.MSC.337(91)}: Code on Noise Levels on Board Ships (mandatory)
\item \textbf{IMO MSC/Circ.645 \& MSC.1/Circ.1580}: Guidelines for vessels/units with DP systems
\item \textbf{IMO MSC.1/Circ.1369}: Interim explanatory notes for SRtP assessment
\end{itemize}

\section{Online Resources}

\subsection{Bureau Veritas Marine \& Offshore}
\begin{itemize}[leftmargin=*]
\item \textbf{Rules Explorer}: \url{https://rulesexplorer.bureauveritas.com/}
\item \textbf{NR467 Rules}: \url{https://marine-offshore.bureauveritas.com/nr467-rules-classification-steel-ships}
\item \textbf{All Rules \& Guidance Notes}: \url{https://marine-offshore.bureauveritas.com/rules-classification-rule-notes-and-guidance-notes}
\end{itemize}

\newpage

% Chapter 14: Contact Information
\chapter{Contact Information \& Document Control}
\label{ch:contact}

\section{For Technical Queries}

\textbf{Bureau Veritas Marine \& Offshore Division}
\begin{itemize}[leftmargin=*]
\item Regional Technical Support Offices
\item Email: Contact through official BV channels
\item Rules Explorer Support: Available through BV website
\end{itemize}

\section{For Urgent Classification Matters}

\begin{itemize}[leftmargin=*]
\item Contact your assigned Plan Approval Manager
\item Regional Marine \& Offshore Director
\end{itemize}

\section{Document Revision History}

\begin{table}[H]
\centering
\caption{Revision History}
\begin{tabular}{|c|c|p{6cm}|p{3cm}|}
\hline
\rowcolor{bvblue}
\textcolor{white}{\textbf{Version}} & \textcolor{white}{\textbf{Date}} & \textcolor{white}{\textbf{Changes}} & \textcolor{white}{\textbf{Author}} \\
\hline
1.0 & 2025 & Initial comprehensive catalog based on NR467 Sections 7, 14, 15 and NR216 & Jules (BV AI Assistant) \\
\hline
1.1 & 2026 & Added Voith-Schneider (cycloidal) and other propulsor types (rim-driven, pump-jet, gill/jet, steering grid); prime-mover chapter (diesel-engine vs.\ electric-motor drive incl.\ full-train TVC); ice class \& winterisation chapter; additional BV notations chapter (COMF, AUT, CLEANSHIP, RP, service notations); expanded checklists, quick-reference tables and standards. Preamble/compile fixes. & Plan Approval Reference (expanded ed.) \\
\hline
1.2 & 2026 & Added Ship-Type \& Service-Notation chapter (ship-type $\rightarrow$ thruster expectation matrix; Safe Return to Port for passenger ships; hazardous-area/Ex for tankers \& gas carriers; offshore/MODU \& DP-intensive vessels; tugs \& escort tugs; thruster-compartment requirements). Added ship-type/service rule sets \& codes to references (IEC 60079, IGC, MODU, HSC, naval, SRtP). & Plan Approval Reference (expanded ed.) \\
\hline
1.3 & 2026 & Added \textbf{MON-SHAFT} shaft-line condition-monitoring notation (Section \ref{sec:monshaft}) with submissions table and cross-references in the notation matrix/quick-guide; added hydraulic-system rule-basis note (Sec.14/Sec.15/piping, ISO 4413, fail-safe pitch). Added \textbf{per-ship-type submission cover sheets} (one-page tick-off) for 15 ship types plus a generic template. & Plan Approval Reference (expanded ed.) \\
\hline
1.4 & 2026 & Readability/accessibility: switched to scalable Latin Modern fonts (crisp, selectable, searchable text) with microtype; enabled navigable PDF bookmarks (open, numbered) and FitH view; added a ``Navigating this PDF'' guide; slightly opened line spacing. & Plan Approval Reference (expanded ed.) \\
\hline
1.5 & 2026 & \textbf{Rule-reference correction} (verified against NR467 Jan 2026): DP is additional class notation \textbf{DYNAPOS} in \textbf{Part F} --- the earlier ``NR216 DP'' citation was wrong (NR216 = Materials \& Welding). Corrected Part D (service notations) / Part E (OSV \& tugs) / Part F (additional class notations) throughout; added NR216/NR266/NR320/NR500/NR566/NR483/NR445 to references. Added ``two axes'' section --- \textbf{manoeuvring vs.\ propulsion} duty and \textbf{open vs.\ ducted (nozzle)} propeller, incl.\ fixed vs.\ steering nozzle. & Plan Approval Reference (expanded ed.) \\
\hline
1.6 & 2026 & Added \textbf{Part A} (Classification and Surveys) to references --- clarifies that notations/ship types are \emph{defined} in Part A Ch.\,1 while Parts D/E/F carry the technical requirements. Companion interactive tool rebuilt as a fully offline local app with a robust clipboard-based Excel/Sheets export (replacing the mobile-unreliable file-download-only export), a spec-capture bar, and a \textbf{contra-rotating propeller (CRP)} option for azimuth/podded units. & Plan Approval Reference (expanded ed.) \\
\hline
\end{tabular}
\end{table}

\section{Important Disclaimers}

\begin{center}
\fcolorbox{warningred}{lightgray}{%
\begin{minipage}{0.9\textwidth}
\textbf{\textcolor{warningred}{⚠ CRITICAL NOTICE FOR SURVEYORS:}}

\begin{enumerate}
\item \textbf{Verification Required}: This catalog is a reference guide compiled from Bureau Veritas rules. Always verify against the latest official BV rules and guidance notes.

\item \textbf{Rule Updates}: BV rules are updated periodically. Ensure you are working with the current version applicable to the vessel's build date and classification notation.

\item \textbf{Complex Cases}: For unusual thruster configurations, hybrid systems, or novel technologies, consult with BV technical specialists before approval.

\item \textbf{Project-Specific Requirements}: Some projects may have additional requirements beyond standard rules (e.g., owner specifications, flag state requirements, special notations).

\item \textbf{Overlapping Requirements}: When thrusters serve multiple functions (propulsion + steering, or DP + steering), ALL applicable requirements must be met.

\item \textbf{FMEA Critical}: For DP systems, the FMEA is the cornerstone document. Inadequate FMEA is grounds for rejection.

\item \textbf{Steering Gear Compliance}: If thrusters provide steering, full SOLAS steering gear requirements apply.

\item \textbf{Material Traceability}: All critical components must have traceable material certificates.

\item \textbf{Software Validation}: Control system software, especially for DP and steering functions, requires validation documentation.

\item \textbf{Cybersecurity}: Modern thruster control systems must address cybersecurity per latest BV guidance and IMO requirements.
\end{enumerate}
\end{minipage}}
\end{center}

\vspace{1cm}

\begin{center}
\textbf{⚠ PLEASE CAREFULLY REVIEW THIS DOCUMENT}

\textit{It may contain errors or omissions. Always verify against official BV rules.}
\end{center}

\newpage

% Back cover
\thispagestyle{empty}
\vspace*{\fill}
\begin{center}
{\Huge\bfseries\color{bvblue} Bureau Veritas}

\vspace{1cm}

{\Large Thruster Requirements}

{\Large Comprehensive Surveyor Catalog}

\vspace{2cm}

{\large Document Version 1.6}

{\large 2026}

\vspace{2cm}

\textit{For Bureau Veritas Plan Approval Surveyors}

\vspace{1cm}

\textbf{Source Links:}

\small
\url{https://marine-offshore.bureauveritas.com/nr467-rules-classification-steel-ships}

\url{https://marine-offshore.bureauveritas.com/rules-classification-rule-notes-and-guidance-notes}

\url{https://rulesexplorer.bureauveritas.com/}

\end{center}
\vspace*{\fill}

\end{document}